% Pentagame-Orakel: Das Buch der Trümpfe
% Mit XeLaTeX kompilieren:  xelatex pentagame-orakel.tex
\documentclass[12pt,paper=a5,DIV=13,headings=small,twoside=false]{scrbook}

\usepackage{fontspec}
\usepackage[headsepline]{scrlayer-scrpage}
\pagestyle{scrheadings}
\manualmark   % Marken von Hand: \markboth{Kapitel}{} nach \addchap, \markright{...} für Abschnitte
\clearpairofpagestyles
% Auf allen Seiten: Kapitel links, Abschnitt (letzte rechte Marke der Seite) rechts
\lehead{\leftmark}\rehead{\rightbotmark}
\lohead{\leftmark}\rohead{\rightbotmark}
\cfoot*{\pagemark}
\setkomafont{pageheadfoot}{\small\sffamily}
\setkomafont{pagenumber}{\small\sffamily}
\usepackage[ngerman]{babel}
% \usepackage[a4paper,margin=2.5cm]{geometry}
\usepackage{booktabs}
\usepackage{tabularx}
\usepackage{array}
\usepackage{enumitem}
\usepackage{needspace}
\usepackage{ragged2e}
\usepackage[hidelinks]{hyperref}

% Schrift für die chinesischen Zeichen (beliebige installierte CJK-Schrift)
\newfontfamily\cjk{Noto Serif CJK TC}

\usepackage{xcolor}
\usepackage{tikz}
\usetikzlibrary{arrows.meta}
\usepackage[simple]{pentpack}  % github.com/penta-jan/pentpack.sty-2.0

\newcommand{\ra}{\,$\rightarrow$\,\allowbreak}
\setlength{\emergencystretch}{2em}
\newcolumntype{L}{>{\raggedright\arraybackslash}X}

% ---------------------------------------------------------------
% Farben der Elemente: nur hier ändern, gilt im ganzen Dokument.
% Die Namen der Farben sind die Namen der Elemente.
% ---------------------------------------------------------------
\colorlet{Holz}{green!70!black}
\colorlet{Feuer}{red!85!black}
\colorlet{Erde}{yellow!85!orange}
\colorlet{Metall}{white}
\colorlet{Wasser}{blue!75!black}

% Wie stark die Farbe hinter dem Wort ist und wie dunkel der Rand
\newcommand{\elementfill}{35}  % Prozent der Elementfarbe im Hintergrund
\newcommand{\elementframe}{65} % Prozent der Elementfarbe im Rand (Rest: schwarz)

% \element{Holz}         -> farbig hinterlegtes "Holz"
% \element[Grün]{Holz}   -> farbig hinterlegtes "Grün" in der Farbe von Holz
% \element*{Holz}        -> so schmal wie das Wort, für Fließtext
% Ein Wort pro Etikett: Zwischen den Wörtern bricht die Zeile normal um.
% Der Rand sorgt dafür, dass auch Weiß (Metall) sichtbar ist.
% Gemeinsamer Stil für alle Element-Etiketten (Text und Abbildung)
\tikzset{element/.style={
    fill=#1!\elementfill, draw=#1!\elementframe!black,
    rounded corners=2pt, inner xsep=2pt, inner ysep=1pt, line width=0.4pt,
    minimum width=12ex}}

\NewDocumentCommand{\element}{s O{#3} m}{%
  \IfBooleanTF{#1}%
    {\tikz[baseline=(e.base)] \node[element=#3, minimum width=0pt] (e) {\vphantom{Äg}#2};}%
    {\tikz[baseline=(e.base)] \node[element=#3] (e) {\vphantom{Äg}#2};}%
}

% Chinesisches Zeichen und Schriftfarbe im Kreis, je Element
\newcommand{\zeichenHolz}{木}   \colorlet{HolzZeichen}{white}
\newcommand{\zeichenFeuer}{火}  \colorlet{FeuerZeichen}{white}
\newcommand{\zeichenErde}{土}   \colorlet{ErdeZeichen}{black}
\newcommand{\zeichenMetall}{金} \colorlet{MetallZeichen}{black}
\newcommand{\zeichenWasser}{水} \colorlet{WasserZeichen}{white}

% Ein Kreis mit Zeichen: \elementkreis{Radius}{x-Position}{Element}
\newcommand{\elementkreis}[3]{%
  \filldraw[fill=#3, draw=black!60, line width=0.5pt] (#2,0) circle (#1);
  \pgfmathsetmacro{\zeichenskala}{#1/1em}%
  \node[text=#3Zeichen, scale=\zeichenskala, inner sep=0pt] at (#2,0)
    {\cjk\csname zeichen#3\endcsname};%
}

% \mycombination[Radius]{Holz}{Feuer}{Erde} -> drei Kreise mit Zeichen
\newcommand{\mycombination}[4][1em]{%
  \tikz[baseline=-0.6ex]{%
    \elementkreis{#1}{0}{#2}
    \elementkreis{#1}{2.3*#1}{#3}
    \elementkreis{#1}{4.6*#1}{#4}
  }%
}

% ---------------------------------------------------------------
% Pentagame-Brett (pentpack.sty, 50 Felder, vereinfacht)
% Im Uhrzeigersinn liegen die Farben in der Reihenfolge des Nährens:
% Weiß (Metall), Blau (Wasser), Grün (Holz), Rot (Feuer), Gelb (Erde).
% ---------------------------------------------------------------
\renewcommand{\mywhite}{Metall}
\renewcommand{\myblue}{Wasser}
\renewcommand{\mygreen}{Holz}
\renewcommand{\myred}{Feuer}
\renewcommand{\myyellow}{Erde}

% Pfeilfarben: von der Ecke (außen) schwarz, von der Kreuzung (innen) rot
\colorlet{pfeilaussen}{black}
\colorlet{pfeilinnen}{red!80!black}
\newcommand{\pfeilfarbe}{pfeilaussen}
\tikzset{musterpfeil/.style={-{Stealth[scale=0.9]}, color=\pfeilfarbe}}
\tikzset{musterpfeilrueck/.style={{Stealth[scale=0.9]}-, color=\pfeilfarbe}} % Spitze am Anfang des Pfads
\tikzset{schrittnr/.style={pos=0.5, auto, swap, inner sep=1.5pt, text=black,
                           font=\footnotesize\sffamily\bfseries}}

% Ein Schritt im Wandlungskreis als Linie auf dem Brett.
% Aktuelle Stelle: \musterpos (Farbposition 0-4), \mustertyp (c = Ecke, j = Kreuzung).
% +1 / +4: zur Nachbarfarbe (Ring bzw. inneres Fünfeck), gleicher Typ.
% +2 / +3: über eine lange Linie zur Zielfarbe, Ecke <-> Kreuzung.
\newcommand{\musterschritt}[2]{% #1 = Schritte im Wandlungskreis, #2 = Nummer
  \pgfmathtruncatemacro{\musterziel}{mod(\musterpos-#1+10,5)}%
  \if c\mustertyp
    \ifnum#1=1 \pbogen{\musterziel}{musterpfeilrueck, shorten <=\cvar} node[schrittnr]{#2};\fi
    \ifnum#1=4 \pbogen{\musterpos}{musterpfeil, shorten >=\cvar} node[schrittnr]{#2};\fi
    \ifnum#1=2 \plinie{\musterpos}{\musterziel a}{musterpfeil, shorten >=\jvar} node[schrittnr]{#2};\gdef\mustertyp{j}\fi
    \ifnum#1=3 \plinie{\musterpos}{\musterziel a}{musterpfeil, shorten >=\jvar} node[schrittnr]{#2};\gdef\mustertyp{j}\fi
  \else
    \ifnum#1=1 \plinie{\musterpos a}{\musterziel a}{musterpfeil, shorten >=\jvar} node[schrittnr]{#2};\fi
    \ifnum#1=4 \plinie{\musterpos a}{\musterziel a}{musterpfeil, shorten >=\jvar} node[schrittnr]{#2};\fi
    \ifnum#1=2 \plinie{\musterpos a}{\musterziel}{musterpfeil, shorten >=\cvar} node[schrittnr]{#2};\gdef\mustertyp{c}\fi
    \ifnum#1=3 \plinie{\musterpos a}{\musterziel}{musterpfeil, shorten >=\cvar} node[schrittnr]{#2};\gdef\mustertyp{c}\fi
  \fi
  \xdef\musterpos{\musterziel}%
}

% \musterbrett{Schritte 1}{Schritte 2}: der Weg ab Ecke 0 (schwarz) und ab Kreuzung 0a (rot),
% auf einem kleinen, grauen 100er-Brett (kleinere Ecken, Pfeile besser sichtbar)
\newcommand{\mustergrau}{black!30}
\colorlet{dreieckfarbe}{yellow!22}
\newcommand{\musterbrett}[2]{{%
  \boardwhite
  \renewcommand{\mywhite}{\mustergrau}\renewcommand{\myblue}{\mustergrau}%
  \renewcommand{\mygreen}{\mustergrau}\renewcommand{\myred}{\mustergrau}%
  \renewcommand{\myyellow}{\mustergrau}%
  \pgfmathtruncatemacro{\dreieckb}{mod(10-#1,5)}%
  \pgfmathtruncatemacro{\dreieckc}{mod(10-#1-#2,5)}%
  \begin{pentdiag}[stops=100, radius=0.16\textwidth, background=white]
    % das Dreieck der drei Elemente, aus den Ecken gebildet, hellgelb hinterlegt
    \begin{scope}[blend mode=multiply]
      \fill[dreieckfarbe] (0.center) -- (\dreieckb.center) -- (\dreieckc.center) -- cycle;
    \end{scope}
    \renewcommand{\pfeilfarbe}{pfeilaussen}%
    \gdef\musterpos{0}\gdef\mustertyp{c}%   von der Ecke oben
    \musterschritt{#1}{}
    \musterschritt{#2}{}
    \renewcommand{\pfeilfarbe}{pfeilinnen}%
    \gdef\musterpos{0}\gdef\mustertyp{j}%   von der Kreuzung unten
    \musterschritt{#1}{}
    \musterschritt{#2}{}
  \end{pentdiag}}}

% \beziehungsbrett{Schritte}: dieselbe Beziehung zweimal, ab Ecke 0 (schwarz) und ab Kreuzung 0a (rot)
\newcommand{\beziehungsbrett}[1]{{%
  \boardwhite
  \renewcommand{\mywhite}{\mustergrau}\renewcommand{\myblue}{\mustergrau}%
  \renewcommand{\mygreen}{\mustergrau}\renewcommand{\myred}{\mustergrau}%
  \renewcommand{\myyellow}{\mustergrau}%
  \begin{pentdiag}[stops=100, radius=0.15\textwidth, background=white]
    \renewcommand{\pfeilfarbe}{pfeilaussen}%
    \gdef\musterpos{0}\gdef\mustertyp{c}%   von der Ecke oben
    \musterschritt{#1}{}
    \renewcommand{\pfeilfarbe}{pfeilinnen}%
    \gdef\musterpos{0}\gdef\mustertyp{j}%   von der Kreuzung unten
    \musterschritt{#1}{}
  \end{pentdiag}}}

% \paarbrett{A}{B}{Schritte}: graues Brett (Geometrie 100, Farbfolge 50), nur Ecke und Kreuzung von A und B in ihren Farben;
% die Beziehung schwarz von der Ecke, rot von der Kreuzung von A aus
% Lage der Elemente auf dem Brett (gegen den Uhrzeigersinn ab Holz oben)
\newcommand{\elposHolz}{0}\newcommand{\elposWasser}{1}\newcommand{\elposMetall}{2}
\newcommand{\elposErde}{3}\newcommand{\elposFeuer}{4}
% welche pentpack-Farbe zu welchem Element gehört
\newcommand{\pfarbHolz}{mygreen}\newcommand{\pfarbFeuer}{myred}\newcommand{\pfarbErde}{myyellow}
\newcommand{\pfarbMetall}{mywhite}\newcommand{\pfarbWasser}{myblue}
\newcommand{\paarbrett}[3]{{%
  \boardwhite
  \renewcommand{\mywhite}{\mustergrau}\renewcommand{\myblue}{\mustergrau}%
  \renewcommand{\mygreen}{\mustergrau}\renewcommand{\myred}{\mustergrau}%
  \renewcommand{\myyellow}{\mustergrau}%
  \expandafter\renewcommand\csname\csname pfarb#1\endcsname\endcsname{#1}%
  \expandafter\renewcommand\csname\csname pfarb#2\endcsname\endcsname{#2}%
  \pentcolours[100]{green,red,yellow,white,blue}% Farbfolge im Uhrzeigersinn ab Holz oben
  \begin{pentdiag}[stops=100, radius=0.17\textwidth, background=white]
    \renewcommand{\pfeilfarbe}{pfeilaussen}%
    \xdef\musterpos{\csname elpos#1\endcsname}\gdef\mustertyp{c}%
    \musterschritt{#3}{}
    \renewcommand{\pfeilfarbe}{pfeilinnen}%
    \xdef\musterpos{\csname elpos#1\endcsname}\gdef\mustertyp{j}%
    \musterschritt{#3}{}
  \end{pentdiag}}}

% \tripelbrett{A}{B}{C}{Muster}: graues Brett, Ecken und Kreuzungen der drei Elemente farbig,
% beide Schritte des Musters schwarz von der Ecke von A, rot von der Kreuzung von A
\expandafter\def\csname musterschritte I\endcsname{{1}{1}}
\expandafter\def\csname musterschritte II\endcsname{{1}{2}}
\expandafter\def\csname musterschritte III\endcsname{{1}{3}}
\expandafter\def\csname musterschritte VII\endcsname{{2}{1}}
\expandafter\def\csname musterschritte IX\endcsname{{2}{2}}
\expandafter\def\csname musterschritte VIII\endcsname{{2}{4}}
\expandafter\def\csname musterschritte X\endcsname{{3}{1}}
\expandafter\def\csname musterschritte XII\endcsname{{3}{3}}
\expandafter\def\csname musterschritte XI\endcsname{{3}{4}}
\expandafter\def\csname musterschritte V\endcsname{{4}{2}}
\expandafter\def\csname musterschritte VI\endcsname{{4}{3}}
\expandafter\def\csname musterschritte IV\endcsname{{4}{4}}
\newcommand{\tripelbrettweg}[2]{%
    \renewcommand{\pfeilfarbe}{pfeilaussen}%
    \xdef\musterpos{\tripelstart}\gdef\mustertyp{c}%
    \musterschritt{#1}{}\musterschritt{#2}{}%
    \renewcommand{\pfeilfarbe}{pfeilinnen}%
    \xdef\musterpos{\tripelstart}\gdef\mustertyp{j}%
    \musterschritt{#1}{}\musterschritt{#2}{}}
\newcommand{\tripelbrett}[4]{{%
  \boardwhite
  \renewcommand{\mywhite}{\mustergrau}\renewcommand{\myblue}{\mustergrau}%
  \renewcommand{\mygreen}{\mustergrau}\renewcommand{\myred}{\mustergrau}%
  \renewcommand{\myyellow}{\mustergrau}%
  \expandafter\renewcommand\csname\csname pfarb#1\endcsname\endcsname{#1}%
  \expandafter\renewcommand\csname\csname pfarb#2\endcsname\endcsname{#2}%
  \expandafter\renewcommand\csname\csname pfarb#3\endcsname\endcsname{#3}%
  \pentcolours[100]{green,red,yellow,white,blue}%
  \xdef\tripelstart{\csname elpos#1\endcsname}%
  \begin{pentdiag}[stops=100, radius=0.17\textwidth, background=white]
    \begin{scope}[blend mode=multiply]% Dreieck der drei Elemente, hellgelb
      \fill[dreieckfarbe] (\csname elpos#1\endcsname.center) -- (\csname elpos#2\endcsname.center) -- (\csname elpos#3\endcsname.center) -- cycle;
    \end{scope}
    \expandafter\expandafter\expandafter\tripelbrettweg\csname musterschritte #4\endcsname
  \end{pentdiag}}}
\def\tripelmuster#1 #2\endtripelmuster{#1}

% \kreis{Holz}: ein einzelner Kreis mit Zeichen (Schlüssel)
\newcommand{\kreis}[2][0.6em]{\tikz[baseline=-0.6ex]{\elementkreis{#1}{0}{#2}}}

% \mypair[Radius]{Holz}{Feuer}: zwei Kreise mit Pfeil dazwischen
\newcommand{\mypair}[3][0.5em]{%
  \tikz[baseline=-0.6ex]{%
    \elementkreis{#1}{0}{#2}
    \draw[-{Stealth[scale=0.6]}, black!70] ({1.25*#1},0) -- ({2.35*#1},0);
    \elementkreis{#1}{3.6*#1}{#3}
  }%
}

% \mypairlink[Radius]{Holz}{Feuer}: wie \mypair, jeder Kreis ein Link auf seine Elementseite
% \XeTeXLinkBox (hyperref): Klickfläche über den ganzen Kreis, nicht nur über das Zeichen
\setlength{\XeTeXLinkMargin}{0pt}
\newcommand{\kreislink}[2]{\hyperlink{element:#2}{\XeTeXLinkBox{\kreis[#1]{#2}}}}
\newcommand{\mypairlink}[3][0.5em]{%
  \kreislink{#1}{#2}%
  \tikz[baseline=-0.6ex]{\path (0,0) -- ({1.6*#1},0);
    \draw[-{Stealth[scale=0.6]}, black!70] ({0.25*#1},0) -- ({1.35*#1},0);}%
  \kreislink{#1}{#3}}

% Farbe der Überschriften von Beziehungen und Mustern (Nummern bleiben schwarz)
\colorlet{ueberschriftfarbe}{red!80!black}

% \blink{Nähren}: Name einer Beziehung als Link auf ihre Seite
\expandafter\def\csname bznr Nähren\endcsname{1}
\expandafter\def\csname bznr Bändigen\endcsname{2}
\expandafter\def\csname bznr Sich stellen\endcsname{3}
\expandafter\def\csname bznr Annehmen\endcsname{4}
% \bnlink{Nähren}: nur der Name der Richtung, als Link
\newcommand{\bnlink}[1]{\hyperlink{beziehung:\csname bznr #1\endcsname}{#1}}
\newcommand{\blink}[1]{\hyperlink{beziehung:\csname bznr #1\endcsname}{\csname bzlabel\csname bznr #1\endcsname\endcsname~#1}}

% Dreiecksform jeder Wendung: stumpf = drei Elemente hintereinander im Kreis, spitz = eines ausgelassen
\expandafter\def\csname musterdreieck I\endcsname{Yin}
\expandafter\def\csname musterdreieck II\endcsname{Yang}
\expandafter\def\csname musterdreieck III\endcsname{Yin}
\expandafter\def\csname musterdreieck IV\endcsname{Yin}
\expandafter\def\csname musterdreieck V\endcsname{Yin}
\expandafter\def\csname musterdreieck VI\endcsname{Yang}
\expandafter\def\csname musterdreieck VII\endcsname{Yang}
\expandafter\def\csname musterdreieck VIII\endcsname{Yin}
\expandafter\def\csname musterdreieck IX\endcsname{Yang}
\expandafter\def\csname musterdreieck X\endcsname{Yin}
\expandafter\def\csname musterdreieck XI\endcsname{Yang}
\expandafter\def\csname musterdreieck XII\endcsname{Yang}
% \dreieckbild{(p1)--(p2)--(p3)--cycle}: kleines Fünfeck mit hervorgehobenem Dreieck
\newcommand{\dreieckbild}[1]{%
  \begin{tikzpicture}[scale=0.9]
    \foreach [count=\i] \e in {Holz,Feuer,Erde,Metall,Wasser}
      \coordinate (p\i) at ({90-72*(\i-1)}:1.3cm);
    \draw[gray!50] (p1)--(p2)--(p3)--(p4)--(p5)--cycle;
    \draw[gray!50] (p1)--(p3)--(p5)--(p2)--(p4)--cycle;
    \draw[very thick] #1;
    \foreach [count=\i] \e in {Holz,Feuer,Erde,Metall,Wasser}
      \filldraw[fill=\e!60, draw=black!60] (p\i) circle (0.17cm);
  \end{tikzpicture}}

% Kopf eines Musters, aufgebaut wie ein Spruch:
% Nummer, Bild, Überschrift, dann die Schritte
% \musterkopf{Nummer}{Name}{Schritte 1}{Schritte 2}{Text 1.}{Text 2.}{Anfang->Ende}
\newcommand{\musterkopf}[7]{%
  \hypertarget{muster:#1}{{\Large\bfseries #1}}\par
  \smallskip
  \musterbrett{#3}{#4}\par
  \smallskip
  {\large\bfseries\color{ueberschriftfarbe} #2}\par
  \smallskip
  {\small\bnlink{#5}, dann \bnlink{#6}: \bnlink{#7}.\\
  Ein \csname musterdreieck #1\endcsname-Dreieck.\par}
  \smallskip}

% Blocksatz-Einleitung innerhalb der zentrierten Kapitel
\newlength{\normalparindent}
\AtBeginDocument{\setlength{\normalparindent}{\parindent}}
\newenvironment{einleitung}
  {\par\justifying\setlength{\parindent}{\normalparindent}\setlength{\parskip}{0pt}}
  {\par}

% \lichtschatten{Licht}{Schatten}: zweiteilige Deutung im Blocksatz (Yin und Yang)
\newcommand{\lichtschatten}[2]{%
  {\justifying\setlength{\parindent}{0pt}%
   \noindent\textit{Im Licht:} #1\par
   \smallskip
   \noindent\textit{Im Schatten:} #2\par}}

% Ein Block (Beziehung, Muster, Spruch) bleibt immer auf einer Seite
% Oben ausgerichtet mit festem Abstand zur Kopfzeile; \linkteil schiebt die Links an den Seitenfuß
\newcommand{\blockoben}{\medskipamount}
\newenvironment{orakelblock}
  {\par\vspace*{\blockoben}\begingroup\centering}
  {\par\endgroup}
\newcommand{\linkteil}{\par\vfill}

% Position im Wandlungskreis, für die Pfeil-Links
\newcommand{\elnumHolz}{0}\newcommand{\elnumFeuer}{1}\newcommand{\elnumErde}{2}
\newcommand{\elnumMetall}{3}\newcommand{\elnumWasser}{4}
% \verbpfeil{Holz}{Feuer}: Pfeil mit dem Verb der Beziehung klein darüber ("nährt" usw.)
\expandafter\def\csname bzverb1\endcsname{nährt}
\expandafter\def\csname bzverb2\endcsname{bändigt}
\expandafter\def\csname bzverb3\endcsname{stellt sich}
\expandafter\def\csname bzverb4\endcsname{nimmt an}
\newcommand{\verbpfeil}[2]{%
  \pgfmathtruncatemacro{\vpschritt}{mod(\csname elnum#2\endcsname-\csname elnum#1\endcsname+5,5)}%
  \,\tikz[baseline=-0.6ex]{%
    \node[font=\tiny, inner sep=0.6pt, anchor=south] (vp) at (0,0.4pt) {\csname bzverb\vpschritt\endcsname};
    \draw[-{Stealth[length=3pt]}, line width=0.45pt] ([xshift=-2pt]vp.south west) -- ([xshift=2pt]vp.south east);}\,}
% \rlink{Holz}{Feuer}: "Holz ->" als ein Link auf das Paar Holz -> Feuer;
% Etikett und Pfeil sind zusammen die Klickfläche
\newcommand{\rlink}[2]{%
  \hyperlink{paar:#1-#2}{\element{#1}\verbpfeil{#1}{#2}}\allowbreak}

% \paar{A}{B}{Titel}{Urteil}{Licht}{Schatten}{Tripel}: Eintrag eines der 20 Paare,
% Nummer (rot, arabisch) = 4 * Position von A im Wandlungskreis + Nummer der Beziehung (+1, −1, +2, −2 = 1..4)
\newcommand{\paar}[7]{%
  \pgfmathtruncatemacro{\paarnr}{mod(\csname elnum#2\endcsname-\csname elnum#1\endcsname+5,5)}%
  \pgfmathtruncatemacro{\paarnum}{4*\csname elnum#1\endcsname+\csname bzidx\paarnr\endcsname}%
  \begin{orakelblock}
   \hypertarget{paar:#1-#2}{{\Large\bfseries\color{ueberschriftfarbe}\paarnum}}\par
   \smallskip
   \paarbrett{#1}{#2}{\paarnr}\par
   \smallskip
   \mypairlink[1em]{#1}{#2}\par
   \smallskip
   {\small\element{#1}\verbpfeil{#1}{#2}\element{#2}}\par
   \medskip
   {\large\bfseries #3}\par
   {\small Richtung:}~\hyperlink{beziehung:\paarnr}{\textit{\csname bzlabel\paarnr\endcsname~\csname bzname\paarnr\endcsname}}\par
   \smallskip
   \textbf{#4}\par
   \smallskip
   \lichtschatten{#5}{#6}
   \linkteil
   #7
  \end{orakelblock}
}
% Beziehungen intern als Schritte im Wandlungskreis (1..4), gezeigt als +1, −1, +2, −2;
% Reihenfolge im Buch: +1, −1, +2, −2 (bzidx)
\expandafter\def\csname bzlabel1\endcsname{+1}\expandafter\def\csname bzidx1\endcsname{1}
\expandafter\def\csname bzlabel4\endcsname{−1}\expandafter\def\csname bzidx4\endcsname{2}
\expandafter\def\csname bzlabel2\endcsname{+2}\expandafter\def\csname bzidx2\endcsname{3}
\expandafter\def\csname bzlabel3\endcsname{−2}\expandafter\def\csname bzidx3\endcsname{4}
\expandafter\def\csname bzname1\endcsname{Nähren}
\expandafter\def\csname bzname2\endcsname{Bändigen}
\expandafter\def\csname bzname3\endcsname{Sich stellen}
\expandafter\def\csname bzname4\endcsname{Annehmen}

% \musterref XI Rückhalt\endmusterref -> "Muster: XI. Der Rückhalt" als Link zum Muster XI
\expandafter\def\csname mustername I\endcsname{Der Strom}
\expandafter\def\csname mustername II\endcsname{Die Prüfung}
\expandafter\def\csname mustername III\endcsname{Die Unterweisung}
\expandafter\def\csname mustername VII\endcsname{Der Damm}
\expandafter\def\csname mustername IX\endcsname{Der strenge Weg}
\expandafter\def\csname mustername VIII\endcsname{Die Abkürzung}
\expandafter\def\csname mustername X\endcsname{Die Läuterung}
\expandafter\def\csname mustername XII\endcsname{Der Aufstand}
\expandafter\def\csname mustername XI\endcsname{Die List}
\expandafter\def\csname mustername V\endcsname{Die gesammelte Kraft}
\expandafter\def\csname mustername VI\endcsname{Der Rückhalt}
\expandafter\def\csname mustername IV\endcsname{Die Wurzelsuche}
% Ausgabe: "Muster: II. Die Prüfung", Nummer und Name als Link
\def\musterref#1 #2\endmusterref{{\small Wendung:}~\hyperlink{muster:#1}{\textit{#1.~\csname mustername #1\endcsname}}}

% Kleine Kreispfeile für die Drehrichtung
\newcommand{\uzs}{\tikz[baseline=-0.6ex]\draw[-{Stealth[scale=0.6]}, line width=0.7pt] (90:1.1ex) arc (90:-210:1.1ex);}  % im Uhrzeigersinn
\newcommand{\guzs}{\tikz[baseline=-0.6ex]\draw[-{Stealth[scale=0.6]}, line width=0.7pt] (90:1.1ex) arc (90:390:1.1ex);}  % gegen den Uhrzeigersinn

% \elementmitfarbe{Holz}{Grün}: Etikett, darunter klein der Farbname (Tabelle der Elemente)
\newcommand{\elementmitfarbe}[2]{%
  \begin{tabular}[t]{@{}c@{}}\element{#1}\\[-0.3ex]{\scriptsize #2}\end{tabular}}

% \spruch{Nr}{Element 1}{Element 2}{Element 3}{Muster}{Name}{Urteil}{Licht}{Schatten}
% Der Spruch (fett, zentriert) ist das knappe Urteil, die Deutung (Blocksatz) die ausführliche Auslegung
% Nummer, Kreise, Elemente, Überschrift, Muster, Spruch: untereinander und zentriert
\newcommand{\spruch}[9]{%
  \begin{orakelblock}
   \hypertarget{spruch:#1}{{\Large\bfseries #1}}\par
   \smallskip
   \tripelbrett{#2}{#3}{#4}{\tripelmuster#5\endtripelmuster}\par
   \smallskip
   \mycombination{#2}{#3}{#4}\par
   \smallskip
   {\small\rlink{#2}{#3}\rlink{#3}{#4}\element{#4}}\par
   \medskip
   {\large\bfseries #6}\par
   \musterref#5\endmusterref\par
   \smallskip
   \textbf{#7}\par
   \smallskip
   \lichtschatten{#8}{#9}
  \end{orakelblock}
}

% \wandlungskreis[art]: die fünf Elemente im Kreis, Holz oben, im Uhrzeigersinn.
%   ohne Option: nur der Kreis des Hervorbringens (sheng)
%   [ke]:     dazu innen das Pentagramm des Bändigens (ke)
%   [gegen]:  die Gegenrichtungen: außen xie, innen wu
\newcommand{\wkscale}{0.82}% Größe des Wandlungskreises (Radius 3 cm mal \wkscale), überall gleich
\newcommand{\wandlungskreis}[1][]{%
  \begin{tikzpicture}[>=latex]
    \pgfmathsetmacro{\R}{3*\wkscale}\pgfmathsetmacro{\RN}{\R+1.25}\pgfmathsetmacro{\RL}{\R+0.55}%
    \pgfmathsetmacro{\d}{deg(0.8/\R)}%
    \def\myradius{\R cm}
    \def\art{#1}\def\artke{ke}\def\artgegen{gegen}%
    \foreach [count=\i] \zeichen/\py/\e in {木/mù/Holz,火/huǒ/Feuer,土/tǔ/Erde,金/jīn/Metall,水/shuǐ/Wasser}{
      \pgfmathsetmacro{\w}{90-72*(\i-1)}
      \node[circle, draw, fill=\e!40, minimum size=1.4cm, align=center, inner sep=1pt] (w\i) at (\w:\myradius)
        {{\Large\cjk\zeichen}\\[-1pt]{\scriptsize\py}};
      \node at (\w:\RN cm) {\scriptsize\e};
      \ifx\art\artgegen
        \draw[<-, thick] ({\w-\d}:\myradius) arc[start angle={\w-\d}, end angle={\w-72+\d}, radius=\myradius];
        \ifnum\i=1 \node[font=\scriptsize, fill=white, inner sep=1pt] at ({\w-36}:\RL cm) {{\cjk 泄} xiè};\fi
      \else
        \draw[->, thick] ({\w-\d}:\myradius) arc[start angle={\w-\d}, end angle={\w-72+\d}, radius=\myradius];
        \ifnum\i=1 \node[font=\scriptsize, fill=white, inner sep=1pt] at ({\w-36}:\RL cm) {{\cjk 生} shēng};\fi
      \fi
    }
    \ifx\art\artke
      \foreach \i/\ii in {1/3,2/4,3/5,4/1,5/2}
        \draw[->, red!75!black] (w\i) -- (w\ii);
      \node[inner sep=1pt, font=\scriptsize] at (0,0) {{\cjk 克} kè};
    \fi
    \ifx\art\artgegen
      \foreach \i/\ii in {3/1,4/2,5/3,1/4,2/5}
        \draw[->, red!75!black] (w\i) -- (w\ii);
      \node[inner sep=1pt, font=\scriptsize] at (0,0) {{\cjk 侮} wǔ};
    \fi
  \end{tikzpicture}}

\title{Pentagame-Orakel\\[0.3em]\large Das Buch der Trümpfe}
\author{Jan Suchanek}
\date{2026}

\makeatletter\renewcommand\@pnumwidth{2em}\makeatother % Platz für dreistellige Seitenzahlen im Inhaltsverzeichnis

% \beziehungsuebersicht: volles, farbiges 50er-Brett (Holz oben), von der grünen Ecke (schwarz) und
% der grünen Kreuzung (rot) die vier Beziehungen; Punkt am Anfang, Pfeil bis in die Mitte des Ziels
\tikzset{bezlabel/.style={font=\scriptsize, align=center, inner sep=1pt, text=black, fill=white, fill opacity=0.85, text opacity=1},
  bezinnen/.style={font=\scriptsize, inner sep=1pt, text=pfeilinnen, fill=white, fill opacity=0.85, text opacity=1},
  bezpfeil/.style={{Circle[length=3.2pt]}-{Stealth[scale=0.9]}, color=\pfeilfarbe},
  bezpfeilrueck/.style={{Stealth[scale=0.9]}-{Circle[length=3.2pt]}, color=\pfeilfarbe}}
\newcommand{\bezabstand}{6.5pt}% Abstand der Pfeilanfänge vom Mittelpunkt
\makeatletter
\newcommand{\beziehungsuebersicht}{{%
  \boardcoloured
  \pentcolours[50]{green,red,yellow,white,blue}% Holz oben
  \def\pp@lw{1.4pt}%
  \renewcommand{\pfeilfarbe}{pfeilaussen}%
  \begin{pentdiag}[stops=50, radius=0.31\textwidth, background=white]
    \pbogen{4}{bezpfeilrueck, shorten >=\bezabstand}
      node[bezlabel, pos=0.4, above right=1pt] {\hyperlink{beziehung:1}{\XeTeXLinkBox{\begin{tabular}{@{}c@{}}+1 Nähren\\{\cjk 生} shēng\end{tabular}}}};
    \pbogen{0}{bezpfeil, shorten <=\bezabstand}
      node[bezlabel, pos=0.6, above left=1pt] {\hyperlink{beziehung:4}{\XeTeXLinkBox{\begin{tabular}{@{}c@{}}−1 Annehmen\\{\cjk 泄} xiè\end{tabular}}}};
    \plinie{0}{3a}{bezpfeil, shorten <=\bezabstand}
      node[bezlabel, pos=0.7, above left=0pt] {\hyperlink{beziehung:2}{\XeTeXLinkBox{\begin{tabular}{@{}c@{}}+2 Bändigen\\{\cjk 克} kè\end{tabular}}}};
    \plinie{0}{2a}{bezpfeil, shorten <=\bezabstand}
      node[bezlabel, pos=0.7, above right=0pt] {\hyperlink{beziehung:3}{\XeTeXLinkBox{\begin{tabular}{@{}c@{}}−2 Sich stellen\\{\cjk 侮} wǔ\end{tabular}}}};
    \renewcommand{\pfeilfarbe}{pfeilinnen}%
    \plinie{0a}{4a}{bezpfeil, shorten <=\bezabstand}
      node[bezinnen, pos=0.42, below left=4pt] {\hyperlink{beziehung:1}{\XeTeXLinkBox{+1 {\cjk 生}}}};
    \plinie{0a}{1a}{bezpfeil, shorten <=\bezabstand}
      node[bezinnen, pos=0.42, below right=4pt] {\hyperlink{beziehung:4}{\XeTeXLinkBox{−1 {\cjk 泄}}}};
    \plinie{0a}{3}{bezpfeil, shorten <=\bezabstand}
      node[bezinnen, pos=0.66, above right=2.5pt] {\hyperlink{beziehung:2}{\XeTeXLinkBox{+2 {\cjk 克}}}};
    \plinie{0a}{2}{bezpfeil, shorten <=\bezabstand}
      node[bezinnen, pos=0.66, above left=2.5pt] {\hyperlink{beziehung:3}{\XeTeXLinkBox{−2 {\cjk 侮}}}};
      % Klickflächen über den Farben: Ecke und Kreuzung führen zur Seite des Elements
    \foreach \pos/\el in {0/Holz,1/Wasser,2/Metall,3/Erde,4/Feuer}{
      \node[inner sep=0pt] at (\pos) {\hyperlink{element:\el}{\XeTeXLinkBox{\hbox to 2\cvar{\vrule height\cvar depth\cvar width0pt\hfil}}}};
      \node[inner sep=0pt] at (\pos a) {\hyperlink{element:\el}{\XeTeXLinkBox{\hbox to 2\jvar{\vrule height\jvar depth\jvar width0pt\hfil}}}};}
  \end{pentdiag}}}
\makeatother

\begin{document}
\frontmatter
\maketitle
% Rückseite des Titels: KOMA lässt sie einseitig weg, daher von Hand
\thispagestyle{empty}
\null\vfill
{\centering
  \copyright\ 2026 Jan Suchanek\\
  Lizenz: \href{https://creativecommons.org/licenses/by-nc/4.0/deed.de}{CC BY-NC 4.0}\\
  \href{https://pentagame.org/}{pentagame.org}\par}
\bigskip
{\small\centering
  Konzept, Aufbau und Auswahl: Jan Suchanek.\\
  Bei der Ausarbeitung der Texte, der Übersetzungen\\
  und des Satzes half das KI-Modell Claude (Anthropic).\\
  Alle Texte wurden vom Autor geprüft und verantwortet.\par}
\clearpage
\tableofcontents
\mainmatter

\addchap{Einführung}\markboth{Einführung}{\mbox{}}

\addsec{Über dieses Buch}\markright{Über dieses Buch}

Wer Pentagame spielt, kennt diesen Moment: Die dritte Figur erreicht ihr Ziel und zieht heraus, die Partie ist gewonnen, und im Zentrum stehen drei Figuren verschiedener Farbe in einer bestimmten Reihenfolge. Jede Farbe steht für eine der fünf Wandlungsphasen des Wu Xing: \element*{Holz}, \element*{Feuer}, \element*{Erde}, \element*{Metall} und \element*{Wasser}. Dieses Buch erzählt, was solche Reihenfolgen bedeuten können, und es baut diese Bedeutung Schritt für Schritt aus den fünf Elementen auf.

\subsection*{Vom Element zum Tripel}

\begin{description}[leftmargin=0pt, itemsep=0.5ex]
  \item[Die fünf Elemente.] Im Wu Xing folgen sie in fester Reihenfolge aufeinander, und jedes geht aus dem vorigen hervor. Zu jedem gehört ein Thema mit einem hellen und einem dunklen Aspekt.
  \item[Die vier Richtungen.] Von einem Element zu einem anderen führt immer genau eine von vier Richtungen: \emph{Nähren} (+1), \emph{Annehmen} (−1), \emph{Bändigen} (+2) oder \emph{Sich stellen} (−2). Welche es ist, hängt nur davon ab, ob man in der Reihenfolge der Elemente einen oder zwei Schritte vor- oder zurückgeht.
  \item[Die 20 Paare.] Ein Element und eine Richtung ergeben ein Paar: zwei Elemente in fester Reihenfolge, zum Beispiel \element*{Holz}\verbpfeil{Holz}{Feuer}\element*{Feuer}. Zu jedem der fünf Elemente gibt es vier Paare: $5 \times 4 = 20$.
  \item[Die zwölf Wendungen.] Zwei Richtungen hintereinander ergeben eine Wendung. Für die erste Richtung gibt es vier Möglichkeiten, für die zweite drei, denn sie darf die erste nicht aufheben: Auf +1 folgt nie −1, auf +2 nie −2. Das ergibt $4 \times 3 = 12$.
  \item[Die 60 Tripel.] Ein Element und eine Wendung ergeben ein Tripel, eine Folge aus drei verschiedenen Elementen, etwa \element*{Holz}\verbpfeil{Holz}{Feuer}\element*{Feuer}\verbpfeil{Feuer}{Erde}\element*{Erde}. Das sind $5 \times 12 = 60$, genau die Folgen, mit denen man beim Pentagame traditionell gewinnt. Jedes Tripel besteht aus zwei Paaren, hier \textit{Holz nährt Feuer} und \textit{Feuer nährt Erde}.
\end{description}

\begin{center}
\small\renewcommand{\arraystretch}{1.15}
\begin{tabular}{@{}lrl@{}}
\toprule
\textbf{Stufe} & \textbf{Anzahl} & \textbf{entsteht aus} \\
\midrule
\hyperref[kap:elemente]{Elemente}    & 5  & dem Wu Xing \\
\hyperref[kap:beziehungen]{Richtungen}  & 4  & 1 oder 2 vor oder zurück \\
\hyperref[kap:paare]{Paare}          & 20 & Element $\times$ Richtung \\
\hyperref[kap:muster]{Wendungen}     & 12 & Richtung $\times$ Richtung \\
\hyperref[kap:tripel]{Tripel}        & 60 & Element $\times$ Wendung \\
\bottomrule
\end{tabular}
\end{center}

Zu jeder Erscheinung gibt es einen Spruch: zu jedem Element, jeder Richtung, jedem Paar, jeder Wendung und jedem Tripel. Ein Spruch besteht aus einem knappen Urteil, fett gesetzt, und einer Deutung darunter. Die Deutung zeigt immer beide Seiten: das Licht und den Schatten.

Die Bilder der Sprüche bauen aufeinander auf. Jedes Paar hat einen eigenen Vorgang aus Natur und Handwerk, wie er tatsächlich geschieht: Holz brennt, der Deich hält das Wasser, die Axt fällt den Baum. Jede Wendung beschreibt eine Lage, wie Menschen sie erleben, etwa die Prüfung oder den Aufstand. Ein Tripel fügt die Vorgänge seiner beiden Paare zu einer Szene zusammen, und seine Wendung gibt ihr den Sinn. Wer die Paare kennt, erkennt sie in den Tripeln wieder.

Das Vorbild ist das \href{https://en.wikipedia.org/wiki/Lingqijing}{Lingqijing} ({\cjk 靈棋經}, „Klassiker der göttlichen Spielsteine“), ein chinesisches Orakel, das wohl in der frühen Wei-Jin-Zeit entstanden ist. Dort ergeben drei Positionen (Himmel, Mensch, Erde) mit je fünf Zuständen 125 Sprüche. Beim Pentagame-Orakel kommt keine Farbe doppelt vor, deshalb sind es 60 Tripel.

Die Sprüche sind keine Übersetzung eines alten Textes, sondern eine moderne Deutung entlang der Wu-Xing-Zyklen. Überliefert sind die fünf Elemente und ihre Reihenfolge, die beiden Zyklen des Hervorbringens ({\cjk 生} shēng) und des Bändigens ({\cjk 克} kè) mit ihren Gegenrichtungen xiè und wǔ sowie die Zuordnungen von Jahreszeit, Himmelsrichtung, Tugend und Gemütsbewegung, Organ, Geschmack und Planet. Nähren und Bändigen sind auch im Deutschen gebräuchliche Namen; Annehmen und Sich stellen sind für dieses Buch gewählt. Die Zählung +1, −1, +2, −2 stammt aus der Notation des \textit{Pentagame Compendium}. Neu sind die zwölf Wendungen mit ihren Namen, die Positionen Wurzel, Weg und Frucht, die Übertragung auf das Pentagame-Brett und alle Sprüche.

Das Wort Trumpf kommt von Triumph: So hießen auch die Bildkarten des ersten Tarots im Italien des 15.~Jahrhunderts, die \textit{trionfi}. Jede gewonnene Partie Pentagame ist ein solcher Triumph, und jeder Spruch dieses Buches ein Trumpf, den man in der Hand hält.

\addsec{Befragung}\markright{Befragung}

Am einfachsten befragst du das Orakel mit einem Beutel. Lege je eine Figur jeder Farbe hinein, etwa eine ganze Mannschaft, formuliere deine Frage und ziehe, ohne hinzusehen:

\begin{enumerate}[leftmargin=*]
  \item \textbf{Zwei Figuren:} Schlag das Paar nach (\hyperref[kap:paare]{„Die 20 Paare“}). Jedes Paar ist gleich wahrscheinlich (1 zu 20).
  \item \textbf{Drei Figuren:} Schlag das Tripel nach (\hyperref[kap:tripel]{„Die 60 Tripel“}). Jedes Tripel ist gleich wahrscheinlich (1 zu 60).
\end{enumerate}

Die Reihenfolge, in der du die Figuren ziehst, ist wichtig, denn jede Position hat ihre eigene Bedeutung:

\begin{center}
\begin{tabular}{lll}
\toprule
\textbf{Position} & \textbf{Bedeutung} & \textbf{Frage dazu} \\
\midrule
1.\ Farbe & \textbf{Wurzel:} die Ausgangslage & Wo stehe ich? \\
2.\ Farbe & \textbf{Weg:} die Wandlung & Was bewegt sich? \\
3.\ Farbe & \textbf{Frucht:} das Ergebnis & Wohin führt es? \\
\bottomrule
\end{tabular}
\end{center}

Beim Paar fallen Weg und Frucht zusammen: Die zweite Farbe zeigt, wohin die Wurzel führt.

Der Beutel ist aber nur der kürzeste Weg. Brett und Buch beruhen auf denselben fünf Farben und denselben Linien. Deshalb lässt sich fast alles, was in einer Partie Pentagame geschieht, hier nachschlagen und deuten:

\begin{itemize}[leftmargin=*]
  \item \textbf{Richtungen:} Jeder Zug, der auf einem farbigen Feld beginnt, einer Ecke oder einer Kreuzung, und auf einem andersfarbigen endet, entspricht einer der vier \hyperref[kap:beziehungen]{Richtungen}. Welche es ist, hängt nur von den beiden Farben ab.
  \item \textbf{Paare:} Ein solcher Zug verbindet zwei Farben in fester Reihenfolge, von wo nach wohin, und ist damit schon ein \hyperref[kap:paare]{Paar}. Ebenso bilden zwei Figuren verschiedener Farbe ein Paar, etwa wenn sie beim Schlagen ihre Plätze tauschen.
  \item \textbf{Wendungen:} Ein Doppelzug mit derselben Figur besteht aus zwei Richtungen und ergibt so eine der zwölf \hyperref[kap:muster]{Wendungen}.
  \item \textbf{Tripel:} Die drei Figuren, mit denen eine Partie gewonnen wird, ergeben in der Reihenfolge, in der sie herausgezogen wurden, ein \hyperref[kap:tripel]{Tripel}. So wird jede Partie zu einem Spruch. Ebenso lassen sich die ersten drei geschlagenen Blöcke lesen, nach der Farbe der Felder, auf denen sie standen: Sie zeigen, unter welchem Zeichen das Spiel beginnt.
\end{itemize}

Das Spiel bringt noch viel reichere Verbindungen hervor: Eine grüne Figur auf einer weißen Kreuzung tauscht mit einer blauen auf einer gelben Ecke, und schon begegnen sich zwei Paare. Wollte man all das deuten, bräuchte es Sprüche für Paare von Paaren und mehr. Dieses Buch beschränkt sich bewusst auf die Grundformen, aus denen sich alles Weitere lesen lässt.

Zum Deuten lies zuerst den Spruch. Von dort aus kannst du auf jeder Stufe weitergehen, denn jeder Eintrag verweist auf die verwandten Einträge der anderen Stufen. Bei jedem Element stehen seine vier Paare. Bei jeder Richtung findest du die fünf Paare, in denen sie vorkommt. Bei jedem Paar stehen seine beiden Elemente, seine Richtung und die sechs Tripel, in denen es steckt. Bei jeder Wendung findest du ihre Richtungen und ihre fünf Tripel, und bei jedem Tripel seine Wendung und die beiden Paare, aus denen es besteht. Am schnellsten findest du eine Kombination über den \hyperref[kap:register]{Schlüssel} am Ende des Buches.

Die Sprüche geben Orientierung, keine Vorhersage.

\newpage
\begingroup
\centering
\setlength{\parindent}{0pt}
\setlength{\parskip}{0.3ex}
\renewcommand{\raggedsection}{\centering}

\addchap{Die fünf Elemente}\markboth{Die fünf Elemente}{\mbox{}}\label{kap:elemente}

\begin{einleitung}
Die Lehre von den fünf Elementen heißt auf Chinesisch Wu Xing ({\cjk 五行} wǔxíng). Wörtlich bedeutet das „fünf Gänge“: \textit{Xing} heißt gehen, sich wandeln. Gemeint sind deshalb weniger fünf Stoffe als fünf Phasen oder Kräfte, die einander ablösen. Holz, Feuer, Erde, Metall und Wasser sind Bilder für diese Wandlungen.

Die Lehre ist mehr als zwei Jahrtausende alt, ihre Wurzeln reichen noch weiter zurück. Die erste ausführliche Aufzählung steht im Kapitel „Der große Plan“ (\textit{Hongfan}, {\cjk 洪範}) des \textit{Buchs der Urkunden} (\textit{Shujing}). Dort hat jedes Element seine Natur: Wasser befeuchtet und sinkt, Feuer lodert und steigt, Holz biegt sich und richtet sich auf, Metall folgt und wandelt sich, die Erde nimmt die Saat auf und gibt die Ernte. Im 3.~Jahrhundert v.\,Chr. baute der Gelehrte Zou Yan daraus eine Lehre von den Kreisläufen der Natur und der Geschichte; spätere Dynastien beriefen sich auf sie, um ihre Herrschaft zu begründen.

In der Han-Zeit (202 v.\,Chr.\ bis 220 n.\,Chr.) wurde Wu Xing zu einem umfassenden Weltsystem. Das \textit{Huainanzi} ordnet den Kosmos nach den fünf Phasen, das \textit{Chunqiu fanlu}, das Dong Zhongshu zugeschrieben wird, verbindet sie mit der konfuzianischen Ethik, und der \textit{Innere Klassiker des Gelben Kaisers} (\textit{Huangdi neijing}) legt sie der Medizin zugrunde: Organe, Gefühle und Geschmäcker gehören je einem Element. Bis heute prägt Wu Xing die chinesische Medizin, den Kalender, die Kampfkünste und das Feng Shui.

Die fünf Elemente ordnet man im Kreis an, im Lauf der Sonne, also im Uhrzeigersinn: Holz, Feuer, Erde, Metall, Wasser. Jedes geht aus dem vorigen hervor und nährt das nächste. Holz nährt das Feuer, das Feuer hinterlässt Asche und damit Erde, in der Erde wächst das Erz, am kühlen Metall schlägt sich Wasser nieder, und das Wasser lässt das Holz wachsen. Dieses Hervorbringen heißt {\cjk 生} shēng, und der Kreis, den es bildet, ist der Wandlungskreis.
\end{einleitung}

\begin{center}
\begin{minipage}{\linewidth}\centering
\wandlungskreis\par
\smallskip
{\small\itshape Der Wandlungskreis: Die fünf Elemente bringen einander hervor ({\cjk 生} shēng).\par}
\end{minipage}
\end{center}

\begin{einleitung}
Im Pentagame steht jedes Element für eine der fünf Farben. Die Farben des heutigen Bretts sind bewusst nach dem Wu Xing gewählt und stehen in seiner Reihenfolge; nur Wasser ist blau statt schwarz. Die Tabelle nennt zu jedem Element seine Zeit und Lage, sein Thema und seine beiden Seiten: die Tugend als Licht, die Gemütsbewegung als Schatten.
\end{einleitung}

\begin{center}
\footnotesize\setlength{\tabcolsep}{2.5pt}
\begin{tabularx}{\textwidth}{@{}lLLLL@{}}
\toprule
\textbf{Element} & \textbf{Zeit, Lage} & \textbf{Thema} & \textbf{Licht:} \newline \textbf{Tugend} & \textbf{Schatten:} \newline \textbf{Emotion} \\
\midrule
\elementmitfarbe{Holz}{Grün} & Frühling, Osten    & Aufbruch, Plan, Wachstum           & Güte \newline \mbox{({\cjk 仁} rén)} & Zorn \newline \mbox{({\cjk 怒} nù)} \\
\elementmitfarbe{Feuer}{Rot} & Sommer, Süden      & Ausdruck, Begeisterung, Begegnung  & Anstand \newline \mbox{({\cjk 禮} lǐ)} & Übermut \newline \mbox{({\cjk 喜} xǐ)} \\
\elementmitfarbe{Erde}{Gelb} & Spätsommer, Übergänge, Mitte & Sammlung, Fürsorge, Halt           & Vertrauen \newline \mbox{({\cjk 信} xìn)} & Grübeln \newline \mbox{({\cjk 思} sī)} \\
\elementmitfarbe{Metall}{Weiß} & Herbst, Westen     & Ernte, Klarheit, Abschied          & Gerechtigkeit \newline \mbox{({\cjk 義} yì)} & Trauer \newline \mbox{({\cjk 憂} yōu)} \\
\elementmitfarbe{Wasser}{Blau} & Winter, Norden     & Ruhe, Tiefe, Ursprung              & Weisheit \newline \mbox{({\cjk 智} zhì)} & Furcht \newline \mbox{({\cjk 恐} kǒng)} \\
\bottomrule
\end{tabularx}
\end{center}

\begin{einleitung}
{\small Die Spalte „Licht“ nennt die fünf konfuzianischen Tugenden ({\cjk 五常} wǔcháng). Zur Tugend des Holzes gehört auch die Liebe: Konfuzius erklärt {\cjk 仁} rén als „die Menschen lieben“ ({\cjk 愛人} ài rén; \textit{Lunyu} 12.22). Die Spalte „Schatten“ folgt den fünf Gemütsbewegungen der chinesischen Medizin ({\cjk 五志} wǔzhì; \textit{Suwen} 5): Zorn, Freude, Nachdenken, Trauer und Furcht. Keine von ihnen ist an sich schlecht. Zum Schatten wird eine Gemütsbewegung erst, wenn sie übermäßig wird oder sich staut; darum steht beim Feuer der Übermut als Schatten der Freude ({\cjk 喜} xǐ) und bei der Erde das Grübeln als Schatten des Nachdenkens ({\cjk 思} sī). Die Erde hat keine eigene Jahreszeit: Ihr gehören der Spätsommer und die Übergänge am Ende jeder Jahreszeit. Die folgenden Seiten stellen die Elemente einzeln vor.\par}
\end{einleitung}
\newpage
\markright{\element{Holz}}
\begin{orakelblock}
\hypertarget{element:Holz}{\tikz{\elementkreis{1.6em}{0}{Holz}}}\par
\smallskip
{\large\bfseries\color{ueberschriftfarbe} Holz}\par
{\cjk 木}~mù\par
\smallskip
\textit{„Holz biegt sich und richtet sich wieder auf.“}\par
\smallskip
{\small\renewcommand{\arraystretch}{1.1}%
\begin{tabular}{@{}r@{\quad}l@{}}
Jahreszeit & Frühling \\
Himmelsrichtung & Osten \\
Farbe & Grün \\
Tugend & Güte ({\cjk 仁} rén) \\
Gemütsbewegung & Zorn ({\cjk 怒} nù) \\
Organ, Geschmack & Leber, sauer \\
Planet & Jupiter
\end{tabular}\par}
\medskip
\textbf{Das Holz ist der Aufbruch. Was keimt, drängt ans Licht.}\par
\smallskip
\lichtschatten{Holz ist die Kraft des Anfangs: der Plan, der Wille, das Wachsen. Im Frühling bricht der Keim durch die Erde, biegsam und doch beharrlich. Wer im Zeichen des Holzes steht, hat Visionen und den Mut, sie anzugehen, und die Güte, die Liebe zu den Menschen, die auch andere wachsen lässt.}
  {Wo das Wachsen keinen Raum findet, wird es zu Ungeduld und Zorn. Darin steckt nur Kraft, die einen Weg sucht. Gib ihr ein Ziel und etwas Zeit, dann wird aus dem Zorn wieder Wachstum.}
\linkteil
{\small Die vier Paare:}\\*[0.5ex] \mbox{\hyperlink{paar:Holz-Feuer}{\mypair{Holz}{Feuer}}} \hfill \mbox{\hyperlink{paar:Holz-Wasser}{\mypair{Holz}{Wasser}}} \hfill \mbox{\hyperlink{paar:Holz-Erde}{\mypair{Holz}{Erde}}} \hfill \mbox{\hyperlink{paar:Holz-Metall}{\mypair{Holz}{Metall}}}\par
\end{orakelblock}
\newpage
\markright{\element{Feuer}}
\begin{orakelblock}
\hypertarget{element:Feuer}{\tikz{\elementkreis{1.6em}{0}{Feuer}}}\par
\smallskip
{\large\bfseries\color{ueberschriftfarbe} Feuer}\par
{\cjk 火}~huǒ\par
\smallskip
\textit{„Feuer lodert und steigt empor.“}\par
\smallskip
{\small\renewcommand{\arraystretch}{1.1}%
\begin{tabular}{@{}r@{\quad}l@{}}
Jahreszeit & Sommer \\
Himmelsrichtung & Süden \\
Farbe & Rot \\
Tugend & Anstand ({\cjk 禮} lǐ) \\
Gemütsbewegung & Freude ({\cjk 喜} xǐ), im Übermaß Übermut \\
Organ, Geschmack & Herz, bitter \\
Planet & Mars
\end{tabular}\par}
\medskip
\textbf{Das Feuer ist die Fülle. Was gewachsen ist, zeigt sich.}\par
\smallskip
\lichtschatten{Feuer ist Ausdruck, Begeisterung und Begegnung. Im Sommer steht alles in voller Blüte. Wer im Zeichen des Feuers steht, wärmt andere, steckt sie an und bringt Menschen zusammen. Der Anstand gibt der Flamme ihre Form, damit sie leuchtet, ohne zu verbrennen.}
  {Feuer kann sich verausgaben oder in Übermut umschlagen. Das ist die Kehrseite großer Freude. Gönn dir die Glut, und sorge für Holz zum Nachlegen.}
\linkteil
{\small Die vier Paare:}\\*[0.5ex] \mbox{\hyperlink{paar:Feuer-Erde}{\mypair{Feuer}{Erde}}} \hfill \mbox{\hyperlink{paar:Feuer-Holz}{\mypair{Feuer}{Holz}}} \hfill \mbox{\hyperlink{paar:Feuer-Metall}{\mypair{Feuer}{Metall}}} \hfill \mbox{\hyperlink{paar:Feuer-Wasser}{\mypair{Feuer}{Wasser}}}\par
\end{orakelblock}
\newpage
\markright{\element{Erde}}
\begin{orakelblock}
\hypertarget{element:Erde}{\tikz{\elementkreis{1.6em}{0}{Erde}}}\par
\smallskip
{\large\bfseries\color{ueberschriftfarbe} Erde}\par
{\cjk 土}~tǔ\par
\smallskip
\textit{„Erde nimmt die Saat auf und gibt die Ernte.“}\par
\smallskip
{\small\renewcommand{\arraystretch}{1.1}%
\begin{tabular}{@{}r@{\quad}l@{}}
Jahreszeit & Spätsommer, Übergänge \\
Himmelsrichtung & Mitte \\
Farbe & Gelb \\
Tugend & Vertrauen ({\cjk 信} xìn) \\
Gemütsbewegung & Nachdenken ({\cjk 思} sī), im Übermaß Grübeln \\
Organ, Geschmack & Milz, süß \\
Planet & Saturn
\end{tabular}\par}
\medskip
\textbf{Die Erde ist die Mitte. Sie trägt alles und verbindet die Zeiten.}\par
\smallskip
\lichtschatten{Erde ist Sammlung, Fürsorge und Halt. Sie hat keine eigene Jahreszeit, sondern gehört zu jedem Übergang, und sie steht in der Mitte der vier Himmelsrichtungen. Wer im Zeichen der Erde steht, gibt anderen Boden unter den Füßen und ist verlässlich.}
  {Wer immer trägt, kann ins Grübeln geraten oder schwer werden. Auch die Erde braucht Lockerung. Lass dich bewegen; fruchtbarer Boden ist nie ganz fest.}
\linkteil
{\small Die vier Paare:}\\*[0.5ex] \mbox{\hyperlink{paar:Erde-Metall}{\mypair{Erde}{Metall}}} \hfill \mbox{\hyperlink{paar:Erde-Feuer}{\mypair{Erde}{Feuer}}} \hfill \mbox{\hyperlink{paar:Erde-Wasser}{\mypair{Erde}{Wasser}}} \hfill \mbox{\hyperlink{paar:Erde-Holz}{\mypair{Erde}{Holz}}}\par
\end{orakelblock}
\newpage
\markright{\element{Metall}}
\begin{orakelblock}
\hypertarget{element:Metall}{\tikz{\elementkreis{1.6em}{0}{Metall}}}\par
\smallskip
{\large\bfseries\color{ueberschriftfarbe} Metall}\par
{\cjk 金}~jīn\par
\smallskip
\textit{„Metall folgt und wandelt sich.“}\par
\smallskip
{\small\renewcommand{\arraystretch}{1.1}%
\begin{tabular}{@{}r@{\quad}l@{}}
Jahreszeit & Herbst \\
Himmelsrichtung & Westen \\
Farbe & Weiß \\
Tugend & Gerechtigkeit ({\cjk 義} yì) \\
Gemütsbewegung & Trauer, Kummer ({\cjk 憂} yōu) \\
Organ, Geschmack & Lunge, scharf \\
Planet & Venus
\end{tabular}\par}
\medskip
\textbf{Das Metall ist die Ernte. Was reif ist, wird geschieden und bewahrt.}\par
\smallskip
\lichtschatten{Metall ist Klarheit, Maß und Abschied. Im Herbst wird geerntet und getrennt: das Korn von der Spreu, das Wichtige vom Überflüssigen. Wer im Zeichen des Metalls steht, urteilt gerecht und schafft Ordnung.}
  {Klarheit kann streng werden, und jeder Abschied bringt Trauer. Beides gehört zur Ernte. Lass die Trauer zu; sie macht Platz für das, was wertvoll bleibt.}
\linkteil
{\small Die vier Paare:}\\*[0.5ex] \mbox{\hyperlink{paar:Metall-Wasser}{\mypair{Metall}{Wasser}}} \hfill \mbox{\hyperlink{paar:Metall-Erde}{\mypair{Metall}{Erde}}} \hfill \mbox{\hyperlink{paar:Metall-Holz}{\mypair{Metall}{Holz}}} \hfill \mbox{\hyperlink{paar:Metall-Feuer}{\mypair{Metall}{Feuer}}}\par
\end{orakelblock}
\newpage
\markright{\element{Wasser}}
\begin{orakelblock}
\hypertarget{element:Wasser}{\tikz{\elementkreis{1.6em}{0}{Wasser}}}\par
\smallskip
{\large\bfseries\color{ueberschriftfarbe} Wasser}\par
{\cjk 水}~shuǐ\par
\smallskip
\textit{„Wasser befeuchtet und sinkt nach unten.“}\par
\smallskip
{\small\renewcommand{\arraystretch}{1.1}%
\begin{tabular}{@{}r@{\quad}l@{}}
Jahreszeit & Winter \\
Himmelsrichtung & Norden \\
Farbe & Schwarz, hier Blau \\
Tugend & Weisheit ({\cjk 智} zhì) \\
Gemütsbewegung & Furcht ({\cjk 恐} kǒng) \\
Organ, Geschmack & Niere, salzig \\
Planet & Merkur
\end{tabular}\par}
\medskip
\textbf{Das Wasser ist der Ursprung. In der Stille sammelt sich die Kraft für den nächsten Kreis.}\par
\smallskip
\lichtschatten{Wasser ist Ruhe, Tiefe und Ausdauer. Im Winter zieht sich alles zurück, und unter dem Eis liegt die Quelle des Frühlings. Wer im Zeichen des Wassers steht, ist weise und geduldig und findet immer einen Weg.}
  {In der Tiefe kann Furcht wohnen, und Rückzug kann einsam machen. Das ist die Stille vor dem Neubeginn. Vertrau ihr; aus jedem Winter wird ein Frühling.}
\linkteil
{\small Die vier Paare:}\\*[0.5ex] \mbox{\hyperlink{paar:Wasser-Holz}{\mypair{Wasser}{Holz}}} \hfill \mbox{\hyperlink{paar:Wasser-Metall}{\mypair{Wasser}{Metall}}} \hfill \mbox{\hyperlink{paar:Wasser-Feuer}{\mypair{Wasser}{Feuer}}} \hfill \mbox{\hyperlink{paar:Wasser-Erde}{\mypair{Wasser}{Erde}}}\par
\end{orakelblock}

\newpage
\addchap{Die vier Richtungen}\markboth{Die vier Richtungen}{\mbox{}}\label{kap:beziehungen}

\begin{einleitung}
Verbindet man die fünf Elemente im Wandlungskreis miteinander, dann führt von jedem Element eine Richtung zu jedem der vier anderen. Ein Schritt im Sonnenlauf vorwärts (+1) ist, wie schon beschrieben, das Hervorbringen oder Nähren ({\cjk 生} shēng). Zwei Schritte vorwärts (+2), also über ein Element hinweg, führt eine Diagonale durch den Kreis: Das ist das Bändigen ({\cjk 克} kè). Zusammen bilden die Diagonalen ein Pentagramm. Diese beiden Richtungen ergeben bereits das vollständige Bild des Wu Xing.
\end{einleitung}

\begin{center}
\begin{minipage}{\linewidth}\centering
\wandlungskreis[ke]\par
\smallskip
{\small\itshape Außen der Kreis des Hervorbringens ({\cjk 生} shēng), innen das Pentagramm des Bändigens ({\cjk 克} kè).\par}
\end{minipage}
\end{center}

\begin{einleitung}
Zu shēng und kè gibt es jeweils ein Gegenstück in der Gegenrichtung. Einen Schritt zurück (−1) liegt das Annehmen ({\cjk 泄} xiè), bei dem das Kind an der Mutter zehrt. Zwei Schritte zurück (−2) liegt das Sich stellen ({\cjk 侮} wǔ), die Auflehnung des Gebändigten gegen seinen Bändiger.
\end{einleitung}

\begin{center}
\begin{minipage}{\linewidth}\centering
\wandlungskreis[gegen]\par
\smallskip
{\small\itshape Die Gegenrichtungen: außen das Annehmen ({\cjk 泄} xiè), innen das Sich stellen ({\cjk 侮} wǔ).\par}
\end{minipage}
\end{center}

\begin{einleitung}
So ergeben sich vier Richtungen, in diesem Buch kurz +1, −1, +2 und −2: Die Zahl sagt, wie weit, das Vorzeichen, ob vor oder zurück. Dieselbe Zählung verwendet die Notation der Pentagame-Züge im \textit{Pentagame Compendium}: ±1 für Züge auf Ring oder Pentagon, ±2 für Züge entlang des Sterns.

\textbf{Vom Wandlungskreis zum Brett.} Im Pentagame-Brett erscheint der Wandlungskreis zweimal: außen auf den Ecken des Rings und innen auf den Kreuzungen des Pentagons, dort am Mittelpunkt gespiegelt. Damit kommt etwas Neues ins Spiel: Die langen Linien des Pentagramms, die Wege von kè und wǔ, verbinden nicht mehr eine Ecke mit einer Ecke, sondern führen von außen nach innen oder von innen nach außen. Wer auf einer solchen Linie zieht, geht im Wandlungskreis zwei Schritte vor oder zurück.

\begin{center}
\begin{minipage}{\linewidth}\centering
\beziehungsuebersicht\par
{\small\itshape Die vier Richtungen von Holz aus: schwarz von der grünen Ecke, rot von der grünen Kreuzung.\par}
\end{minipage}
\end{center}

\needspace{4\baselineskip}\noindent So ordnen sich die vier Richtungen:

\begin{itemize}[leftmargin=*]
  \item \textbf{Nähren und Annehmen} (+1 und −1, {\cjk 生} shēng und {\cjk 泄} xiè) sind Bewegungen \emph{auf} einem der beiden Kreise, entweder auf dem Ring oder auf dem Pentagon: im Uhrzeigersinn vor, gegen den Uhrzeigersinn zurück.
  \item \textbf{Bändigen und Sich stellen} (+2 und −2, {\cjk 克} kè und {\cjk 侮} wǔ) sind Bewegungen \emph{zwischen} den Kreisen, auf einer langen Linie vom Ring zum Pentagon oder umgekehrt: Bändigen gegen den Uhrzeigersinn, Sich stellen im Uhrzeigersinn. Das Vorzeichen bezieht sich also auf den Wandlungskreis, nicht auf die Drehrichtung auf dem Brett.
\end{itemize}
\end{einleitung}

\begin{center}
\footnotesize\setlength{\tabcolsep}{2pt}\renewcommand{\arraystretch}{1.3}
\begin{tabularx}{\textwidth}{@{}lLcll@{}}
\toprule
\textbf{Richtung} & \textbf{Im Kreis} & \textbf{Drehsinn} & \textbf{Weg} & \textbf{Strecke} \\
\midrule
\hyperlink{beziehung:1}{\textbf{+1} Nähren}       & 1 vor \newline\mbox{({\cjk 生} shēng)} & \uzs  & innen oder außen & kurz \\
\hyperlink{beziehung:4}{\textbf{−1} Annehmen}     & 1 zurück \newline\mbox{({\cjk 泄} xiè)} & \guzs & innen oder außen & kurz \\
\hyperlink{beziehung:2}{\textbf{+2} Bändigen}     & 2 vor \newline\mbox{({\cjk 克} kè)} & \guzs & innen und außen  & lang \\
\hyperlink{beziehung:3}{\textbf{−2} Sich stellen} & 2 zurück \newline\mbox{({\cjk 侮} wǔ)} & \uzs  & innen und außen  & lang \\
\bottomrule
\end{tabularx}
\end{center}

\begin{einleitung}
Deshalb zeigen die Bilder auf den folgenden Seiten Bändigen und Sich stellen nicht als Diagonalen in einem Kreis, sondern als Übergänge zwischen außen und innen. Das heißt nicht, dass nur diese Züge dazugehören. Wer auf dem Ring zwei Schritte vorwärts geht, kommt ebenfalls beim übernächsten Element an, und auch das ist Bändigen; diese beiden Schritte bilden die Wendung \hyperlink{muster:I}{I, Der Strom}. Die lange Linie ist nur der direkte Weg.

Jedes Bild zeigt seine Richtung zweimal: schwarz von außen, von der Ecke oben, und rot von innen, von der Kreuzung unten in der Mitte. Die Bilder sind bewusst grau: Die Ecke oben kann jedes der fünf Elemente sein, und die Richtung sieht für alle gleich aus.
\end{einleitung}


\newpage
\markright{+1 \quad Nähren}
\begin{orakelblock}
\hypertarget{beziehung:1}{{\Large\bfseries +1}}\par
\smallskip
\beziehungsbrett{1}\par
\smallskip
{\large\bfseries\color{ueberschriftfarbe} Nähren}\par
{\cjk 生}~shēng\par
\smallskip

\textit{Im Wandlungskreis: ein Schritt vorwärts, der Sheng-Zyklus.}\\
\textit{Auf dem Brett: im Uhrzeigersinn außen auf dem Ring oder innen auf dem Pentagon voran.}\par
\smallskip

\textbf{Du gibst Kraft an dein „Kind“ weiter. Was du nährst, wächst durch dich.}\par
\smallskip
\lichtschatten{Nähren ist die einfachste und schönste Bewegung: Du gibst weiter, was du hast, und etwas wächst daran. Wie die Mutter das Kind, wie das Holz das Feuer, wie der Regen den Baum. Wer nährt, verliert nichts. Er wird Teil von etwas, das größer ist als er selbst.}
  {Nähren wird zum Erschöpfen: Du gibst und gibst, bis nichts mehr da ist. Oder es wird zum Verwöhnen, das dem anderen das eigene Wachsen abnimmt. Manchmal steckt hinter dem Geben auch der Wunsch, gebraucht zu werden. Nähre, aber vergiss nicht, selbst zu trinken.}
\linkteil
{\small Alle fünf Paare:}\\*[0.5ex] \mbox{\hyperlink{paar:Holz-Feuer}{\mypair{Holz}{Feuer}}} \hfill \mbox{\hyperlink{paar:Feuer-Erde}{\mypair{Feuer}{Erde}}} \hfill \mbox{\hyperlink{paar:Erde-Metall}{\mypair{Erde}{Metall}}} \hfill \mbox{\hyperlink{paar:Metall-Wasser}{\mypair{Metall}{Wasser}}} \hfill \mbox{\hyperlink{paar:Wasser-Holz}{\mypair{Wasser}{Holz}}}\par
\end{orakelblock}
\newpage

\markright{−1 \quad Annehmen}
\begin{orakelblock}
\hypertarget{beziehung:4}{{\Large\bfseries −1}}\par
\smallskip
\beziehungsbrett{4}\par
\smallskip
{\large\bfseries\color{ueberschriftfarbe} Annehmen}\par
{\cjk 泄}~xiè\par
\smallskip

\textit{Im Wandlungskreis: ein Schritt zurück, gegen den Sheng-Zyklus.}\\
\textit{Auf dem Brett: gegen den Uhrzeigersinn außen auf dem Ring oder innen auf dem Pentagon zurück.}\par
\smallskip

\textbf{Du nimmst an, was deine „Mutter“ dir gibt, die Quelle, die dich nährt.}\par
\smallskip
\lichtschatten{Annehmen heißt, sich nähren zu lassen: von dem, was dich hervorgebracht hat, von Menschen, Orten, Erinnerungen, von der Stille. Zugleich nimmst du der Mutter etwas von ihrer Fülle ab: In der Tradition ist das der sanfteste Weg, ein Übermaß auszugleichen.}
  {Das Kind zehrt an der Mutter: Wer immer nur nimmt, erschöpft die Quelle. Manchmal fällt es auch schwer, überhaupt etwas anzunehmen. Nimm an, was dir gegeben wird, und danke dafür; dann bleibt die Quelle lebendig.}
\linkteil
{\small Alle fünf Paare:}\\*[0.5ex] \mbox{\hyperlink{paar:Holz-Wasser}{\mypair{Holz}{Wasser}}} \hfill \mbox{\hyperlink{paar:Feuer-Holz}{\mypair{Feuer}{Holz}}} \hfill \mbox{\hyperlink{paar:Erde-Feuer}{\mypair{Erde}{Feuer}}} \hfill \mbox{\hyperlink{paar:Metall-Erde}{\mypair{Metall}{Erde}}} \hfill \mbox{\hyperlink{paar:Wasser-Metall}{\mypair{Wasser}{Metall}}}\par
\end{orakelblock}
\newpage

\markright{+2 \quad Bändigen}
\begin{orakelblock}
\hypertarget{beziehung:2}{{\Large\bfseries +2}}\par
\smallskip
\beziehungsbrett{2}\par
\smallskip
{\large\bfseries\color{ueberschriftfarbe} Bändigen}\par
{\cjk 克}~kè\par
\smallskip

\textit{Im Wandlungskreis: zwei Schritte vorwärts, der Ke-Zyklus.}\\
\textit{Auf dem Brett: gegen den Uhrzeigersinn vom Ring nach innen oder von innen zum Ring.}\par
\smallskip

\textbf{Du setzt Grenzen und gibst Form. Was du bändigst, braucht deine Ordnung.}\par
\smallskip
\lichtschatten{Bändigen heißt Form geben. Das Wasser braucht ein Ufer, das Feuer einen Herd, das Wachstum einen Schnitt. Wer bändigt, schützt das Lebendige davor, sich selbst zu zerstören. Gute Grenzen sind ein Ausdruck von Fürsorge.}
  {Aus Bändigen wird Beherrschen, in der Tradition {\cjk 乘} chéng, das Überwältigen. Wer nur noch kontrolliert, erstickt, was er schützen wollte. Auch das Gegenteil gilt: Wer gar keine Grenzen setzt, lässt das Gegenüber allein. Frag dich ehrlich, ob deine Ordnung dient oder ob sie nur dich beruhigt.}
\linkteil
{\small Alle fünf Paare:}\\*[0.5ex] \mbox{\hyperlink{paar:Holz-Erde}{\mypair{Holz}{Erde}}} \hfill \mbox{\hyperlink{paar:Feuer-Metall}{\mypair{Feuer}{Metall}}} \hfill \mbox{\hyperlink{paar:Erde-Wasser}{\mypair{Erde}{Wasser}}} \hfill \mbox{\hyperlink{paar:Metall-Holz}{\mypair{Metall}{Holz}}} \hfill \mbox{\hyperlink{paar:Wasser-Feuer}{\mypair{Wasser}{Feuer}}}\par
\end{orakelblock}
\newpage

\markright{−2 \quad Sich stellen}
\begin{orakelblock}
\hypertarget{beziehung:3}{{\Large\bfseries −2}}\par
\smallskip
\beziehungsbrett{3}\par
\smallskip
{\large\bfseries\color{ueberschriftfarbe} Sich stellen}\par
{\cjk 侮}~wǔ\par
\smallskip

\textit{Im Wandlungskreis: zwei Schritte zurück, gegen den Ke-Zyklus.}\\
\textit{Auf dem Brett: im Uhrzeigersinn vom Ring nach innen oder von innen zum Ring.}\par
\smallskip

\textbf{Du gehst auf das zu, was dich bändigt. Dort wirst du geprüft.}\par
\smallskip
\lichtschatten{Sich stellen heißt, der eigenen Prüfung nicht auszuweichen. Du gehst auf das zu, was dich begrenzt, und schaust ihm in die Augen. Oft verliert das, was uns bändigt, seine Macht, sobald wir es ansehen. Die Tradition sieht in dieser Richtung vor allem eine Störung; hier ist sie auch Mut zur eigenen Prüfung.}
  {Aus dem Sich-Stellen wird Auflehnung um jeden Preis, wie die Tradition sie beschreibt: Das Gebändigte verhöhnt seinen Bändiger und fordert jede Grenze heraus, auch die schützende. Oder du stellst dich nie und lässt dich ein Leben lang bändigen. Wähle deine Prüfungen mit Bedacht.}
\linkteil
{\small Alle fünf Paare:}\\*[0.5ex] \mbox{\hyperlink{paar:Holz-Metall}{\mypair{Holz}{Metall}}} \hfill \mbox{\hyperlink{paar:Feuer-Wasser}{\mypair{Feuer}{Wasser}}} \hfill \mbox{\hyperlink{paar:Erde-Holz}{\mypair{Erde}{Holz}}} \hfill \mbox{\hyperlink{paar:Metall-Feuer}{\mypair{Metall}{Feuer}}} \hfill \mbox{\hyperlink{paar:Wasser-Erde}{\mypair{Wasser}{Erde}}}\par
\end{orakelblock}
\newpage

\addchap{Die 20 Paare}\markboth{Die 20 Paare}{\mbox{}}\label{kap:paare}

\begin{einleitung}
Ein Element und eine Richtung ergeben ein Paar: zwei Elemente in fester Reihenfolge. Zu jedem der fünf Elemente gibt es vier Paare, eines für jede Richtung, zusammen also 20. Zwei Paare, die aneinander anschließen, ergeben später ein Tripel: Holz, Feuer, Erde etwa besteht aus \textit{Holz nährt Feuer} und \textit{Feuer nährt Erde}.

\textbf{Zehn Verbindungen, zwanzig Paare.} Im Wandlungskreis gibt es zehn Verbindungen: fünf Bögen des Hervorbringens und fünf Linien des Bändigens. Jede lässt sich von beiden Enden aus lesen, und so werden daraus 20 Paare. \textit{Holz nährt Feuer} (+1) und \textit{Feuer nimmt Holz an} (−1) sind dieselbe Verbindung, einmal von der Mutter und einmal vom Kind aus gesehen. So hat jedes Paar sein Gegenstück. Ein Gegensatz wie Hell und Dunkel ist ein Paar trotzdem nicht, sondern eine gerichtete Beziehung zwischen zwei verschiedenen Kräften.

\textbf{Mutter und Kind.} Die Tradition beschreibt jedes Element über seine Beziehungen zu den anderen. Das \textit{Nanjing}, ein Klassiker der Medizin, vergleicht das Hervorbringen mit Mutter und Kind (\textit{Nanjing} 69 und 75): Was mich hervorbringt, ist meine Mutter, was ich hervorbringe, mein Kind. Beim Bändigen nennt der \textit{Neijing} (\textit{Suwen} 67) das, was ich bändige, „das von mir Besiegte“ ({\cjk 所勝} suǒ shèng), und das, was mich bändigt, „das von mir nicht Besiegte“ ({\cjk 所不勝} suǒ bù shèng). Ein Paar ist darum nie gleichrangig: Einer gibt, der andere empfängt; einer formt, der andere wird geformt. Die chinesische Schicksalskunde liest diese Beziehungen sogar als Rollen im menschlichen Leben, etwa als Eltern und Kinder oder als Obrigkeit.

\textbf{Gleichgewicht.} Hervorbringen und Bändigen brauchen einander. Dong Zhongshu fasst das Schema in einen Satz: „Benachbarte bringen einander hervor, übernächste besiegen einander“ ({\cjk 比相生而間相勝}; \textit{Chunqiu fanlu} 58). Der \textit{Neijing} ergänzt: „Übermaß schadet; was folgt, setzt ihm eine Grenze; durch die Grenze entstehen Wachsen und Wandlung“ ({\cjk 亢則害，承乃制，制則生化}; \textit{Suwen} 68). Auch das Nähren kann zu viel sein: Zu viel Wasser lässt das Holz faulen. Und das Bändigen ist nicht nur Zwang: Erst im Feuer wird das Eisen formbar, erst unter der Axt der Baum zum Balken. Darum zeigt sich in jedem Paar beides, das Gelingen und die Gefahr.

\textbf{Bilder.} Jedes Paar hat ein eigenes Bild: einen Vorgang aus Natur oder Handwerk, wie er tatsächlich geschieht. Holz brennt, der Keim bricht durch die Erde, das Holz klemmt die Säge. Das Urteil nennt zuerst diesen Vorgang und dann, was er für dich bedeutet. Dieselben Bilder kehren in den Tripeln wieder.
\end{einleitung}

\newpage
\markright{\element{Holz}}
\paar{Holz}{Feuer}{Das Scheit}{Holz brennt und wird zu Licht und Wärme. Was du dir vorgenommen hast, entfacht Begeisterung.}
  {Was du geplant hast, fängt Feuer, bei dir und bei anderen. Jetzt ist die Zeit, es zu zeigen und zu teilen. Der Plan wird dabei verbraucht, und das ist gut so: Dafür war er da.}
  {Ein Feuer, das zu schnell brennt, frisst das Holz, bevor es Wärme gibt. Ein Scheit, das nie ins Feuer kommt, bleibt nur Holz.}
  {{\small Alle sechs Tripel:}\par\smallskip
   \begin{tabular*}{\linewidth}{@{}l@{\extracolsep{\fill}}ll@{}}
   \mbox{\hyperlink{spruch:1}{\mycombination[0.5em]{Holz}{Feuer}{Erde}~1}} & \mbox{\hyperlink{spruch:2}{\mycombination[0.5em]{Holz}{Feuer}{Metall}~2}} & \mbox{\hyperlink{spruch:3}{\mycombination[0.5em]{Holz}{Feuer}{Wasser}~3}} \\[0.3ex]
   \mbox{\hyperlink{spruch:34}{\mycombination[0.5em]{Erde}{Holz}{Feuer}~34}} & \mbox{\hyperlink{spruch:43}{\mycombination[0.5em]{Metall}{Holz}{Feuer}~43}} & \mbox{\hyperlink{spruch:49}{\mycombination[0.5em]{Wasser}{Holz}{Feuer}~49}}
   \end{tabular*}\par}
\newpage
\paar{Holz}{Wasser}{Die Wurzel}{Die Wurzel zieht Wasser aus der Tiefe. Bevor du weiterwächst, trinkst du.}
  {Du wendest dich dem zu, was dich nährt: der Stille, dem Wissen, deiner Herkunft. Das ist kein Stillstand. Die Wurzel arbeitet, auch wenn oben nichts zu sehen ist.}
  {Ein Baum, der nur nach oben wächst, vertrocknet im ersten heißen Sommer. Einer, der nur in die Tiefe treibt, kommt nie ans Licht.}
  {{\small Alle sechs Tripel:}\par\smallskip
   \begin{tabular*}{\linewidth}{@{}l@{\extracolsep{\fill}}ll@{}}
   \mbox{\hyperlink{spruch:4}{\mycombination[0.5em]{Holz}{Wasser}{Metall}~4}} & \mbox{\hyperlink{spruch:5}{\mycombination[0.5em]{Holz}{Wasser}{Feuer}~5}} & \mbox{\hyperlink{spruch:6}{\mycombination[0.5em]{Holz}{Wasser}{Erde}~6}} \\[0.3ex]
   \mbox{\hyperlink{spruch:16}{\mycombination[0.5em]{Feuer}{Holz}{Wasser}~16}} & \mbox{\hyperlink{spruch:35}{\mycombination[0.5em]{Erde}{Holz}{Wasser}~35}} & \mbox{\hyperlink{spruch:44}{\mycombination[0.5em]{Metall}{Holz}{Wasser}~44}}
   \end{tabular*}\par}
\newpage
\paar{Holz}{Erde}{Der Keim}{Der Keim bricht durch die verkrustete Erde. Dein Wille lockert, was festgefahren war.}
  {Was lange fest lag, kommt in Bewegung, weil etwas in dir wachsen will: Gewohnheiten, Abläufe, ein Boden, auf dem nichts mehr geschah. Ein Keim ist zart und bricht doch die Kruste. Wo er durch ist, dringen Luft und Regen in den Boden.}
  {Ein Keim, der zu tief liegt, verbraucht seine Kraft, bevor er das Licht erreicht. Einer, der die Kruste nicht durchbricht, erstickt darunter.}
  {{\small Alle sechs Tripel:}\par\smallskip
   \begin{tabular*}{\linewidth}{@{}l@{\extracolsep{\fill}}ll@{}}
   \mbox{\hyperlink{spruch:7}{\mycombination[0.5em]{Holz}{Erde}{Metall}~7}} & \mbox{\hyperlink{spruch:8}{\mycombination[0.5em]{Holz}{Erde}{Feuer}~8}} & \mbox{\hyperlink{spruch:9}{\mycombination[0.5em]{Holz}{Erde}{Wasser}~9}} \\[0.3ex]
   \mbox{\hyperlink{spruch:17}{\mycombination[0.5em]{Feuer}{Holz}{Erde}~17}} & \mbox{\hyperlink{spruch:45}{\mycombination[0.5em]{Metall}{Holz}{Erde}~45}} & \mbox{\hyperlink{spruch:50}{\mycombination[0.5em]{Wasser}{Holz}{Erde}~50}}
   \end{tabular*}\par}
\newpage
\paar{Holz}{Metall}{Die Säge}{Das Holz klemmt die Säge. Du gibst nicht nach, wo man dich zurechtschneiden will.}
  {Man will dich kürzen, stutzen, passend machen: mit Regeln, Kritik, Erwartungen. Du hältst dagegen, nicht mit Gewalt, sondern mit deiner eigenen Spannung. Was gewachsen ist, lässt sich nicht beliebig zerteilen.}
  {Ein Stamm, der jede Säge klemmt, wird am Ende mit der Axt gefällt. Wer jedem Schnitt nachgibt, wird zu Brettern nach fremdem Maß.}
  {{\small Alle sechs Tripel:}\par\smallskip
   \begin{tabular*}{\linewidth}{@{}l@{\extracolsep{\fill}}ll@{}}
   \mbox{\hyperlink{spruch:10}{\mycombination[0.5em]{Holz}{Metall}{Wasser}~10}} & \mbox{\hyperlink{spruch:11}{\mycombination[0.5em]{Holz}{Metall}{Erde}~11}} & \mbox{\hyperlink{spruch:12}{\mycombination[0.5em]{Holz}{Metall}{Feuer}~12}} \\[0.3ex]
   \mbox{\hyperlink{spruch:18}{\mycombination[0.5em]{Feuer}{Holz}{Metall}~18}} & \mbox{\hyperlink{spruch:36}{\mycombination[0.5em]{Erde}{Holz}{Metall}~36}} & \mbox{\hyperlink{spruch:51}{\mycombination[0.5em]{Wasser}{Holz}{Metall}~51}}
   \end{tabular*}\par}
\newpage
\markright{\element{Feuer}}
\paar{Feuer}{Erde}{Die Asche}{Das Feuer brennt herunter, und seine Asche düngt den Boden. Aus deiner Begeisterung wird etwas, das trägt.}
  {Was du mit Leidenschaft begonnen hast, kommt zur Ruhe und wird fruchtbar. Die Flammen sind kleiner geworden, aber nichts ist verloren. Auf diesem Boden können andere säen.}
  {Wer nur noch von der Asche lebt, sät nichts mehr. Wer immer nur lodert, düngt keinen Boden.}
  {{\small Alle sechs Tripel:}\par\smallskip
   \begin{tabular*}{\linewidth}{@{}l@{\extracolsep{\fill}}ll@{}}
   \mbox{\hyperlink{spruch:13}{\mycombination[0.5em]{Feuer}{Erde}{Metall}~13}} & \mbox{\hyperlink{spruch:14}{\mycombination[0.5em]{Feuer}{Erde}{Wasser}~14}} & \mbox{\hyperlink{spruch:15}{\mycombination[0.5em]{Feuer}{Erde}{Holz}~15}} \\[0.3ex]
   \mbox{\hyperlink{spruch:1}{\mycombination[0.5em]{Holz}{Feuer}{Erde}~1}} & \mbox{\hyperlink{spruch:46}{\mycombination[0.5em]{Metall}{Feuer}{Erde}~46}} & \mbox{\hyperlink{spruch:55}{\mycombination[0.5em]{Wasser}{Feuer}{Erde}~55}}
   \end{tabular*}\par}
\newpage
\paar{Feuer}{Holz}{Die Flamme}{Die Flamme greift nach dem Scheit. Du holst dir neue Nahrung für dein Feuer.}
  {Deine Begeisterung braucht Stoff, und du weißt, wo du ihn findest: in Plänen und Ideen, in den Menschen und Vorhaben, die dich einmal entzündet haben. Greif zu. Was dich nährt, darf dich nähren.}
  {Ein Feuer, das alles Holz in Reichweite verschlingt, steht bald im Leeren. Eines, das nicht zugreift, geht aus, obwohl das Holz daneben liegt.}
  {{\small Alle sechs Tripel:}\par\smallskip
   \begin{tabular*}{\linewidth}{@{}l@{\extracolsep{\fill}}ll@{}}
   \mbox{\hyperlink{spruch:16}{\mycombination[0.5em]{Feuer}{Holz}{Wasser}~16}} & \mbox{\hyperlink{spruch:17}{\mycombination[0.5em]{Feuer}{Holz}{Erde}~17}} & \mbox{\hyperlink{spruch:18}{\mycombination[0.5em]{Feuer}{Holz}{Metall}~18}} \\[0.3ex]
   \mbox{\hyperlink{spruch:28}{\mycombination[0.5em]{Erde}{Feuer}{Holz}~28}} & \mbox{\hyperlink{spruch:47}{\mycombination[0.5em]{Metall}{Feuer}{Holz}~47}} & \mbox{\hyperlink{spruch:56}{\mycombination[0.5em]{Wasser}{Feuer}{Holz}~56}}
   \end{tabular*}\par}
\newpage
\paar{Feuer}{Metall}{Die Esse}{Das Feuer macht das Eisen weich. In deiner Glut wird formbar, was starr war.}
  {Vor dir liegt etwas Hartes: eine Regel, eine verhärtete Haltung, eine festgefahrene Lage. Mit Gewalt richtest du nichts aus, mit Wärme schon. Halte es in deine Glut, bis es nachgibt, und dann gib ihm Form.}
  {Zu viel Hitze verbrennt das Eisen. Bei zu wenig schmiedet man kaltes Eisen, und kaltes Eisen bricht.}
  {{\small Alle sechs Tripel:}\par\smallskip
   \begin{tabular*}{\linewidth}{@{}l@{\extracolsep{\fill}}ll@{}}
   \mbox{\hyperlink{spruch:19}{\mycombination[0.5em]{Feuer}{Metall}{Wasser}~19}} & \mbox{\hyperlink{spruch:20}{\mycombination[0.5em]{Feuer}{Metall}{Erde}~20}} & \mbox{\hyperlink{spruch:21}{\mycombination[0.5em]{Feuer}{Metall}{Holz}~21}} \\[0.3ex]
   \mbox{\hyperlink{spruch:2}{\mycombination[0.5em]{Holz}{Feuer}{Metall}~2}} & \mbox{\hyperlink{spruch:29}{\mycombination[0.5em]{Erde}{Feuer}{Metall}~29}} & \mbox{\hyperlink{spruch:57}{\mycombination[0.5em]{Wasser}{Feuer}{Metall}~57}}
   \end{tabular*}\par}
\newpage
\paar{Feuer}{Wasser}{Der Schauer}{Das Feuer zischt im Regen und brennt weiter. Du hältst deine Begeisterung gegen die Ernüchterung.}
  {Von außen kommt Kühle: Zweifel, nüchterne Einwände, ein kaltes Wort. Du weichst nicht aus und lässt dich nicht löschen. Was nach dem Regen noch brennt, ist echt.}
  {Ein Feuer, das gegen jeden Schauer anlodert, verbraucht sein Holz im Dampf. Eines, das beim ersten Tropfen ausgeht, war nur Zunder.}
  {{\small Alle sechs Tripel:}\par\smallskip
   \begin{tabular*}{\linewidth}{@{}l@{\extracolsep{\fill}}ll@{}}
   \mbox{\hyperlink{spruch:22}{\mycombination[0.5em]{Feuer}{Wasser}{Holz}~22}} & \mbox{\hyperlink{spruch:23}{\mycombination[0.5em]{Feuer}{Wasser}{Metall}~23}} & \mbox{\hyperlink{spruch:24}{\mycombination[0.5em]{Feuer}{Wasser}{Erde}~24}} \\[0.3ex]
   \mbox{\hyperlink{spruch:3}{\mycombination[0.5em]{Holz}{Feuer}{Wasser}~3}} & \mbox{\hyperlink{spruch:30}{\mycombination[0.5em]{Erde}{Feuer}{Wasser}~30}} & \mbox{\hyperlink{spruch:48}{\mycombination[0.5em]{Metall}{Feuer}{Wasser}~48}}
   \end{tabular*}\par}
\newpage
\markright{\element{Erde}}
\paar{Erde}{Metall}{Das Erz}{Im Gestein wächst über lange Zeit das Erz. Aus deiner Geduld und Sorgfalt entsteht etwas Kostbares.}
  {Du hast gepflegt, gehütet, durchgetragen, oft ohne zu sehen, wozu. Jetzt zeigt sich, was darin herangereift ist: etwas Klares, Wertvolles, das bleibt. Lass es ans Licht.}
  {Ein Berg, der zu schnell ausgebeutet wird, ist bald leer. Einer, in dem nie gegraben wird, behält seinen Schatz für sich.}
  {{\small Alle sechs Tripel:}\par\smallskip
   \begin{tabular*}{\linewidth}{@{}l@{\extracolsep{\fill}}ll@{}}
   \mbox{\hyperlink{spruch:25}{\mycombination[0.5em]{Erde}{Metall}{Wasser}~25}} & \mbox{\hyperlink{spruch:26}{\mycombination[0.5em]{Erde}{Metall}{Holz}~26}} & \mbox{\hyperlink{spruch:27}{\mycombination[0.5em]{Erde}{Metall}{Feuer}~27}} \\[0.3ex]
   \mbox{\hyperlink{spruch:7}{\mycombination[0.5em]{Holz}{Erde}{Metall}~7}} & \mbox{\hyperlink{spruch:13}{\mycombination[0.5em]{Feuer}{Erde}{Metall}~13}} & \mbox{\hyperlink{spruch:58}{\mycombination[0.5em]{Wasser}{Erde}{Metall}~58}}
   \end{tabular*}\par}
\newpage
\paar{Erde}{Feuer}{Das Gefäß}{Im Feuer wird der Ton zum Gefäß. Du nimmst die Wärme an, und dein Halt bekommt Gestalt.}
  {Du bist verlässlich, tragfähig, geduldig. Jetzt lässt du dich erwärmen: von Freude, von Nähe, von einer Begeisterung, die nicht deine eigene sein muss. Erst diese Hitze macht aus weichem Ton ein Gefäß, das hält.}
  {Ton, der zu schnell ins Feuer kommt, springt. Ton, der nie gebrannt wird, zerfällt im ersten Regen.}
  {{\small Alle sechs Tripel:}\par\smallskip
   \begin{tabular*}{\linewidth}{@{}l@{\extracolsep{\fill}}ll@{}}
   \mbox{\hyperlink{spruch:28}{\mycombination[0.5em]{Erde}{Feuer}{Holz}~28}} & \mbox{\hyperlink{spruch:29}{\mycombination[0.5em]{Erde}{Feuer}{Metall}~29}} & \mbox{\hyperlink{spruch:30}{\mycombination[0.5em]{Erde}{Feuer}{Wasser}~30}} \\[0.3ex]
   \mbox{\hyperlink{spruch:8}{\mycombination[0.5em]{Holz}{Erde}{Feuer}~8}} & \mbox{\hyperlink{spruch:40}{\mycombination[0.5em]{Metall}{Erde}{Feuer}~40}} & \mbox{\hyperlink{spruch:59}{\mycombination[0.5em]{Wasser}{Erde}{Feuer}~59}}
   \end{tabular*}\par}
\newpage
\paar{Erde}{Wasser}{Der Deich}{Der Deich hält das Wasser. Du gibst dem Fließenden Halt und Richtung.}
  {Um dich herum ist viel in Bewegung: Gefühle, Ideen, Menschen, die nicht wissen, wohin. Bei dir finden sie eine Grenze, an der sie sich sammeln können. Das ist keine Enge, sondern Schutz.}
  {Ein Deich ohne Auslass bricht irgendwann unter dem Druck des Wassers. Ohne Deich läuft das Land voll.}
  {{\small Alle sechs Tripel:}\par\smallskip
   \begin{tabular*}{\linewidth}{@{}l@{\extracolsep{\fill}}ll@{}}
   \mbox{\hyperlink{spruch:31}{\mycombination[0.5em]{Erde}{Wasser}{Holz}~31}} & \mbox{\hyperlink{spruch:32}{\mycombination[0.5em]{Erde}{Wasser}{Metall}~32}} & \mbox{\hyperlink{spruch:33}{\mycombination[0.5em]{Erde}{Wasser}{Feuer}~33}} \\[0.3ex]
   \mbox{\hyperlink{spruch:9}{\mycombination[0.5em]{Holz}{Erde}{Wasser}~9}} & \mbox{\hyperlink{spruch:14}{\mycombination[0.5em]{Feuer}{Erde}{Wasser}~14}} & \mbox{\hyperlink{spruch:41}{\mycombination[0.5em]{Metall}{Erde}{Wasser}~41}}
   \end{tabular*}\par}
\newpage
\paar{Erde}{Holz}{Der feste Grund}{Der Boden hält die Wurzel fest. Du lässt dich nicht von allem umgraben, was in dich hineinwachsen will.}
  {Etwas Neues drängt zu dir herein: ein Plan, ein Mensch, eine Veränderung. Du gibst nicht einfach nach. Aber gerade indem du festhältst, gibst du dem Neuen den Halt, den es zum Stehen braucht.}
  {Boden, der zu fest ist, lässt keine Wurzel ein, und nichts wächst darauf. Boden, der überall nachgibt, hält keinen Baum im Sturm.}
  {{\small Alle sechs Tripel:}\par\smallskip
   \begin{tabular*}{\linewidth}{@{}l@{\extracolsep{\fill}}ll@{}}
   \mbox{\hyperlink{spruch:34}{\mycombination[0.5em]{Erde}{Holz}{Feuer}~34}} & \mbox{\hyperlink{spruch:35}{\mycombination[0.5em]{Erde}{Holz}{Wasser}~35}} & \mbox{\hyperlink{spruch:36}{\mycombination[0.5em]{Erde}{Holz}{Metall}~36}} \\[0.3ex]
   \mbox{\hyperlink{spruch:15}{\mycombination[0.5em]{Feuer}{Erde}{Holz}~15}} & \mbox{\hyperlink{spruch:42}{\mycombination[0.5em]{Metall}{Erde}{Holz}~42}} & \mbox{\hyperlink{spruch:60}{\mycombination[0.5em]{Wasser}{Erde}{Holz}~60}}
   \end{tabular*}\par}
\newpage
\markright{\element{Metall}}
\paar{Metall}{Wasser}{Der Tau}{Auf kaltem Metall sammelt sich Tau. Aus deiner Klarheit werden Ruhe und Tiefe.}
  {Du hast entschieden, getrennt, aufgeräumt. Nun setzt sich etwas ab, still und klar wie Tropfen am Morgen. Was du geklärt hast, beginnt zu fließen.}
  {Metall, das nie abkühlt, sammelt keinen Tau. Ist es zu kalt, setzt sich statt Tau nur Reif an.}
  {{\small Alle sechs Tripel:}\par\smallskip
   \begin{tabular*}{\linewidth}{@{}l@{\extracolsep{\fill}}ll@{}}
   \mbox{\hyperlink{spruch:37}{\mycombination[0.5em]{Metall}{Wasser}{Holz}~37}} & \mbox{\hyperlink{spruch:38}{\mycombination[0.5em]{Metall}{Wasser}{Feuer}~38}} & \mbox{\hyperlink{spruch:39}{\mycombination[0.5em]{Metall}{Wasser}{Erde}~39}} \\[0.3ex]
   \mbox{\hyperlink{spruch:10}{\mycombination[0.5em]{Holz}{Metall}{Wasser}~10}} & \mbox{\hyperlink{spruch:19}{\mycombination[0.5em]{Feuer}{Metall}{Wasser}~19}} & \mbox{\hyperlink{spruch:25}{\mycombination[0.5em]{Erde}{Metall}{Wasser}~25}}
   \end{tabular*}\par}
\newpage
\paar{Metall}{Erde}{Der Berg}{Das Erz kommt aus dem Berg. Du wendest dich dem zu, was dich hervorgebracht hat.}
  {Hinter deiner Klarheit und Strenge liegt ein Grund, der dich getragen hat: Menschen, die für dich gesorgt haben, ein Ort, eine lange Geduld. Kehr dorthin zurück und nimm an, was er dir gibt. Deine Härte bekommt dadurch Wärme.}
  {Wer im Berg bleibt, wird nie zu Metall. Wer den Berg vergisst, weiß nicht mehr, woraus er gemacht ist.}
  {{\small Alle sechs Tripel:}\par\smallskip
   \begin{tabular*}{\linewidth}{@{}l@{\extracolsep{\fill}}ll@{}}
   \mbox{\hyperlink{spruch:40}{\mycombination[0.5em]{Metall}{Erde}{Feuer}~40}} & \mbox{\hyperlink{spruch:41}{\mycombination[0.5em]{Metall}{Erde}{Wasser}~41}} & \mbox{\hyperlink{spruch:42}{\mycombination[0.5em]{Metall}{Erde}{Holz}~42}} \\[0.3ex]
   \mbox{\hyperlink{spruch:11}{\mycombination[0.5em]{Holz}{Metall}{Erde}~11}} & \mbox{\hyperlink{spruch:20}{\mycombination[0.5em]{Feuer}{Metall}{Erde}~20}} & \mbox{\hyperlink{spruch:52}{\mycombination[0.5em]{Wasser}{Metall}{Erde}~52}}
   \end{tabular*}\par}
\newpage
\paar{Metall}{Holz}{Die Axt}{Die Axt fällt den Baum. Du triffst eine klare Entscheidung und schneidest ab, was zu viel ist.}
  {Etwas ist zu groß geworden, zu wild, zu verzweigt. Du siehst es und handelst: ein Schnitt, ein Nein, ein Ende. Das ist nicht kalt. Wo ein Baum fällt, kommt Licht auf den Boden, und Neues kann wachsen.}
  {Wer jeden Baum fällt, steht im Kahlschlag. Wer nie schlägt, dem wächst der Wald zu.}
  {{\small Alle sechs Tripel:}\par\smallskip
   \begin{tabular*}{\linewidth}{@{}l@{\extracolsep{\fill}}ll@{}}
   \mbox{\hyperlink{spruch:43}{\mycombination[0.5em]{Metall}{Holz}{Feuer}~43}} & \mbox{\hyperlink{spruch:44}{\mycombination[0.5em]{Metall}{Holz}{Wasser}~44}} & \mbox{\hyperlink{spruch:45}{\mycombination[0.5em]{Metall}{Holz}{Erde}~45}} \\[0.3ex]
   \mbox{\hyperlink{spruch:21}{\mycombination[0.5em]{Feuer}{Metall}{Holz}~21}} & \mbox{\hyperlink{spruch:26}{\mycombination[0.5em]{Erde}{Metall}{Holz}~26}} & \mbox{\hyperlink{spruch:53}{\mycombination[0.5em]{Wasser}{Metall}{Holz}~53}}
   \end{tabular*}\par}
\newpage
\paar{Metall}{Feuer}{Das Eisen}{Das Eisen hält der Glut stand. Du setzt dich der Hitze aus und behältst deine Form.}
  {Um dich herum wird es heiß: Leidenschaft, Druck, Gefühle, die dich verformen könnten. Du weichst nicht aus. Du bleibst, was du bist, und wirst dabei nicht spröde, sondern lebendiger.}
  {Wer jeder Hitze standhalten will, wird spröde wie zu hart gehärteter Stahl. Wer jeder Hitze nachgibt, verliert seine Form.}
  {{\small Alle sechs Tripel:}\par\smallskip
   \begin{tabular*}{\linewidth}{@{}l@{\extracolsep{\fill}}ll@{}}
   \mbox{\hyperlink{spruch:46}{\mycombination[0.5em]{Metall}{Feuer}{Erde}~46}} & \mbox{\hyperlink{spruch:47}{\mycombination[0.5em]{Metall}{Feuer}{Holz}~47}} & \mbox{\hyperlink{spruch:48}{\mycombination[0.5em]{Metall}{Feuer}{Wasser}~48}} \\[0.3ex]
   \mbox{\hyperlink{spruch:12}{\mycombination[0.5em]{Holz}{Metall}{Feuer}~12}} & \mbox{\hyperlink{spruch:27}{\mycombination[0.5em]{Erde}{Metall}{Feuer}~27}} & \mbox{\hyperlink{spruch:54}{\mycombination[0.5em]{Wasser}{Metall}{Feuer}~54}}
   \end{tabular*}\par}
\newpage
\markright{\element{Wasser}}
\paar{Wasser}{Holz}{Der Regen}{Der Regen tränkt die Wurzeln. Aus deiner Ruhe und deinem Wissen wächst etwas Neues.}
  {Du hast lange still gelegen, gesammelt, nachgedacht. Jetzt gibst du davon ab, und etwas fängt an zu wachsen: ein Plan, ein Mensch, ein Vorhaben. Du musst nicht daran ziehen. Es wächst von selbst, wenn es genug zu trinken hat.}
  {Zu viel Regen lässt die Wurzeln faulen. Bei zu wenig verdorrt der Trieb, bevor er stark ist.}
  {{\small Alle sechs Tripel:}\par\smallskip
   \begin{tabular*}{\linewidth}{@{}l@{\extracolsep{\fill}}ll@{}}
   \mbox{\hyperlink{spruch:49}{\mycombination[0.5em]{Wasser}{Holz}{Feuer}~49}} & \mbox{\hyperlink{spruch:50}{\mycombination[0.5em]{Wasser}{Holz}{Erde}~50}} & \mbox{\hyperlink{spruch:51}{\mycombination[0.5em]{Wasser}{Holz}{Metall}~51}} \\[0.3ex]
   \mbox{\hyperlink{spruch:22}{\mycombination[0.5em]{Feuer}{Wasser}{Holz}~22}} & \mbox{\hyperlink{spruch:31}{\mycombination[0.5em]{Erde}{Wasser}{Holz}~31}} & \mbox{\hyperlink{spruch:37}{\mycombination[0.5em]{Metall}{Wasser}{Holz}~37}}
   \end{tabular*}\par}
\newpage
\paar{Wasser}{Metall}{Die Quelle}{Die Quelle entspringt dem Fels. Du kehrst zu dem zurück, was du klar weißt.}
  {In der Unruhe suchst du etwas Festes, und du findest es in dir: Grundsätze, Gewissheiten, ein klares Ja und ein klares Nein. Daraus fließt dein Handeln, frisch wie Wasser aus dem Stein.}
  {Wasser, das zu lange im Fels steht, wird hart. Wasser, das keinen Fels findet, versickert.}
  {{\small Alle sechs Tripel:}\par\smallskip
   \begin{tabular*}{\linewidth}{@{}l@{\extracolsep{\fill}}ll@{}}
   \mbox{\hyperlink{spruch:52}{\mycombination[0.5em]{Wasser}{Metall}{Erde}~52}} & \mbox{\hyperlink{spruch:53}{\mycombination[0.5em]{Wasser}{Metall}{Holz}~53}} & \mbox{\hyperlink{spruch:54}{\mycombination[0.5em]{Wasser}{Metall}{Feuer}~54}} \\[0.3ex]
   \mbox{\hyperlink{spruch:4}{\mycombination[0.5em]{Holz}{Wasser}{Metall}~4}} & \mbox{\hyperlink{spruch:23}{\mycombination[0.5em]{Feuer}{Wasser}{Metall}~23}} & \mbox{\hyperlink{spruch:32}{\mycombination[0.5em]{Erde}{Wasser}{Metall}~32}}
   \end{tabular*}\par}
\newpage
\paar{Wasser}{Feuer}{Das Löschen}{Wasser löscht die Glut. Du bringst Ruhe in eine aufgeheizte Lage.}
  {Um dich herum brennt es: Streit, Aufregung, Leidenschaften, die außer Kontrolle geraten. Du bleibst ruhig, und deine Ruhe breitet sich aus. Was eben noch loderte, kommt zur Besinnung.}
  {Wer jedes Feuer löscht, sitzt bald im Kalten. Zu wenig Wasser zischt nur und verdampft.}
  {{\small Alle sechs Tripel:}\par\smallskip
   \begin{tabular*}{\linewidth}{@{}l@{\extracolsep{\fill}}ll@{}}
   \mbox{\hyperlink{spruch:55}{\mycombination[0.5em]{Wasser}{Feuer}{Erde}~55}} & \mbox{\hyperlink{spruch:56}{\mycombination[0.5em]{Wasser}{Feuer}{Holz}~56}} & \mbox{\hyperlink{spruch:57}{\mycombination[0.5em]{Wasser}{Feuer}{Metall}~57}} \\[0.3ex]
   \mbox{\hyperlink{spruch:5}{\mycombination[0.5em]{Holz}{Wasser}{Feuer}~5}} & \mbox{\hyperlink{spruch:33}{\mycombination[0.5em]{Erde}{Wasser}{Feuer}~33}} & \mbox{\hyperlink{spruch:38}{\mycombination[0.5em]{Metall}{Wasser}{Feuer}~38}}
   \end{tabular*}\par}
\newpage
\paar{Wasser}{Erde}{Das Flussbett}{Der Fluss gräbt sich sein Bett. Du stellst dich deinen Grenzen und findest in ihnen deinen Weg.}
  {Es gibt Grenzen, die du nicht einfach überspülen kannst. Du stellst dich ihnen, geduldig und beharrlich, Jahr um Jahr. Mit der Zeit formst du selbst das Ufer, das dich führt.}
  {Ein Fluss, der sich zu tief eingräbt, findet aus seiner Schlucht nicht mehr heraus. Einer, der jedem Hindernis ausweicht, versandet in der Ebene.}
  {{\small Alle sechs Tripel:}\par\smallskip
   \begin{tabular*}{\linewidth}{@{}l@{\extracolsep{\fill}}ll@{}}
   \mbox{\hyperlink{spruch:58}{\mycombination[0.5em]{Wasser}{Erde}{Metall}~58}} & \mbox{\hyperlink{spruch:59}{\mycombination[0.5em]{Wasser}{Erde}{Feuer}~59}} & \mbox{\hyperlink{spruch:60}{\mycombination[0.5em]{Wasser}{Erde}{Holz}~60}} \\[0.3ex]
   \mbox{\hyperlink{spruch:6}{\mycombination[0.5em]{Holz}{Wasser}{Erde}~6}} & \mbox{\hyperlink{spruch:24}{\mycombination[0.5em]{Feuer}{Wasser}{Erde}~24}} & \mbox{\hyperlink{spruch:39}{\mycombination[0.5em]{Metall}{Wasser}{Erde}~39}}
   \end{tabular*}\par}

\newpage

\addchap{Die zwölf Wendungen}\markboth{Die zwölf Wendungen}{\mbox{}}\label{kap:muster}

\begin{einleitung}
Eine Richtung führt von einem Element zu einem anderen. Folgen zwei Richtungen aufeinander, entsteht ein Doppelschritt: vom ersten Element über ein zweites zu einem dritten. Ein solcher Doppelschritt heißt hier Wendung. Mit ihm ist zugleich eine dritte Richtung gegeben, die vom Anfang unmittelbar zum Ende führt. Nach zweimal Nähren (+1, +1) etwa steht man zwei Schritte weiter im Kreis, und dorthin führt vom Anfang aus das Bändigen (+2).

Für die erste Richtung gibt es vier Möglichkeiten, für die zweite nur drei, denn sie darf die erste nicht aufheben: Auf +1 folgt nie −1, auf +2 nie −2, sonst stünde man wieder am Anfang. So ergeben sich zwölf Wendungen.

Eine Wendung ist an kein Element gebunden. Erst wenn sie bei einem bestimmten Element beginnt, wird aus ihr eines der 60 Tripel des folgenden Kapitels; unter jeder Wendung stehen deshalb die fünf Tripel, in denen sie vorkommt. Weil sie keinem Element gehört, sprechen ihre Sprüche nicht in Bildern aus der Natur, sondern von Lagen, wie Menschen sie erleben: von der Prüfung, der Unterweisung, dem Aufstand.

Das Bild über jeder Wendung zeigt ihre zwei Richtungen als Pfeile auf dem Brett, wie bei den Richtungen schwarz von außen und rot von innen; die Ecke oben kann jedes Element sein. Spielzüge sind damit nicht gemeint. Das Brett macht nur sichtbar, wohin zwei Schritte im Wandlungskreis führen.

\textbf{Dreiecke.} Die drei Elemente einer Wendung bilden ein Dreieck. Liegen sie im Wandlungskreis nebeneinander, ist das Dreieck stumpf: Zwei Seiten verbinden benachbarte Elemente (±1), eine Seite übernächste (±2). Es ist vom Hervorbringen geprägt, ein Stück vom Kreis. Ist zwischen ihnen ein Element ausgelassen, ist das Dreieck spitz, mit zwei Seiten ±2 und einer Seite ±1. Es ist vom Bändigen geprägt, ein Stück vom Stern. Im Bild über jeder Wendung ist ihr Dreieck hellgelb hinterlegt.

\end{einleitung}

\begin{center}
\begin{tabular}{c@{\hspace{3em}}c}
\dreieckbild{(p1)--(p2)--(p3)--cycle} & \dreieckbild{(p1)--(p2)--(p4)--cycle} \\[0.5ex]
{\small\begin{tabular}[t]{@{}c@{}}stumpf\\Holz, Feuer, Erde\end{tabular}} & {\small\begin{tabular}[t]{@{}c@{}}spitz\\Holz, Feuer, Metall\end{tabular}}
\end{tabular}
\end{center}

\begin{einleitung}
Die beiden Arten von Dreiecken legen einen Vergleich mit Yin und Yang nahe. Das stumpfe Dreieck, geprägt von Nähren und Empfangen, von Mutter und Kind, ist ein Yin-Dreieck; das spitze, geprägt von Grenze, Form und Spannung, ein Yang-Dreieck. Die Tradition ordnet Yin und Yang den Elementen zu, nicht den beiden Zyklen; diese Zuordnung ist eine Deutung dieses Buches. Mit Licht und Schatten ist sie nicht zu verwechseln: Jede Wendung hat beide Seiten, und das Dreieck sagt nur, ob in ihr mehr Yin oder mehr Yang liegt. Sechs Wendungen sind Yin, sechs Yang; welches Dreieck eine Wendung bildet, steht auf ihrer Seite.
\end{einleitung}



\newpage
\markright{I \quad Der Strom}
\begin{orakelblock}
\musterkopf{I}{Der Strom}{1}{1}{Nähren}{Nähren}{Bändigen}
\textbf{Was du gibst, wird weitergegeben und wächst über zwei Stufen. Am Ende bist du es, der ihm das Maß vorgibt, wie eine Großmutter dem Enkel.}\par
\smallskip
\lichtschatten{Du setzt etwas in Gang, und es läuft weiter, von dir zu anderen und von ihnen weiter. Du musst nicht alles selbst tun. Aber du bleibst nicht außen vor: Gerade weil du am Anfang standest, hört man auf dich, wenn es Form braucht.}
  {Wer alles weiterlaufen lässt, erkennt am Ende nicht wieder, was er angestoßen hat. Wer sich in jede Stufe einmischt, lässt den Strom nicht fließen.}
\linkteil
{\small Alle fünf Tripel:}\\*[0.5ex] \mbox{\hyperlink{spruch:1}{\mycombination[0.5em]{Holz}{Feuer}{Erde}~1}} \hfill \mbox{\hyperlink{spruch:13}{\mycombination[0.5em]{Feuer}{Erde}{Metall}~13}} \hfill \mbox{\hyperlink{spruch:25}{\mycombination[0.5em]{Erde}{Metall}{Wasser}~25}} \hfill \mbox{\hyperlink{spruch:37}{\mycombination[0.5em]{Metall}{Wasser}{Holz}~37}} \hfill \mbox{\hyperlink{spruch:49}{\mycombination[0.5em]{Wasser}{Holz}{Feuer}~49}}\par
\end{orakelblock}
\newpage

\markright{II \quad Die Prüfung}
\begin{orakelblock}
\musterkopf{II}{Die Prüfung}{1}{2}{Nähren}{Bändigen}{Sich stellen}
\textbf{Du stärkst, was das Strenge in Schach hält. Doch am Ende stehst du der Prüfung selbst gegenüber, und niemand kann sie für dich ablegen.}\par
\smallskip
\lichtschatten{Du hast andere unterstützt, und mit ihrer Hilfe ist manches Harte weicher geworden. Das hat dich vorbereitet. Jetzt kommt die Prüfung, und du trittst selbst an. Du bist besser gerüstet, als du denkst.}
  {Wer sich nur auf seine Helfer verlässt, steht in der Prüfung mit leeren Händen da. Wer jede Prüfung meidet, erfährt nie, was er kann.}
\linkteil
{\small Alle fünf Tripel:}\\*[0.5ex] \mbox{\hyperlink{spruch:2}{\mycombination[0.5em]{Holz}{Feuer}{Metall}~2}} \hfill \mbox{\hyperlink{spruch:14}{\mycombination[0.5em]{Feuer}{Erde}{Wasser}~14}} \hfill \mbox{\hyperlink{spruch:26}{\mycombination[0.5em]{Erde}{Metall}{Holz}~26}} \hfill \mbox{\hyperlink{spruch:38}{\mycombination[0.5em]{Metall}{Wasser}{Feuer}~38}} \hfill \mbox{\hyperlink{spruch:50}{\mycombination[0.5em]{Wasser}{Holz}{Erde}~50}}\par
\end{orakelblock}
\newpage

\markright{III \quad Die Unterweisung}
\begin{orakelblock}
\musterkopf{III}{Die Unterweisung}{1}{3}{Nähren}{Sich stellen}{Annehmen}
\textbf{Du förderst, was sich bewähren muss. Während es sich stellt, erinnerst du dich, woher du selbst kommst.}\par
\smallskip
\lichtschatten{Du gibst weiter, was du kannst, und siehst zu, wie es sich bewährt. Das ist mehr als Lehren. Wer erklärt, versteht selbst neu, und wer begleitet, erinnert sich an die eigenen Lehrer. Wer lehrt, lernt am meisten.}
  {Ein Lehrer, der nur gibt, wird belehrend. Einer, der seine Schüler nicht in die Prüfung lässt, lernt nichts von ihnen.}
\linkteil
{\small Alle fünf Tripel:}\\*[0.5ex] \mbox{\hyperlink{spruch:3}{\mycombination[0.5em]{Holz}{Feuer}{Wasser}~3}} \hfill \mbox{\hyperlink{spruch:15}{\mycombination[0.5em]{Feuer}{Erde}{Holz}~15}} \hfill \mbox{\hyperlink{spruch:27}{\mycombination[0.5em]{Erde}{Metall}{Feuer}~27}} \hfill \mbox{\hyperlink{spruch:39}{\mycombination[0.5em]{Metall}{Wasser}{Erde}~39}} \hfill \mbox{\hyperlink{spruch:51}{\mycombination[0.5em]{Wasser}{Holz}{Metall}~51}}\par
\end{orakelblock}
\newpage

\markright{IV \quad Die Wurzelsuche}
\begin{orakelblock}
\musterkopf{IV}{Die Wurzelsuche}{4}{4}{Annehmen}{Annehmen}{Sich stellen}
\textbf{Du gehst zurück zu deinen Ursprüngen und noch hinter sie. Dort findest du, was dich bis heute begrenzt, und stellst dich ihm.}\par
\smallskip
\lichtschatten{Es ist Zeit für einen Blick zurück, weiter als gewohnt. Hinter dem, was dich genährt hat, liegt etwas Älteres: ein Satz, eine Angst, eine Regel, die nie jemand ausgesprochen hat. Schau sie an. Was erkannt ist, verliert seine Macht.}
  {Wer immer tiefer gräbt, findet für alles eine Erklärung und für nichts einen Ausweg. Wer nie zurückschaut, wird von dem bestimmt, was er nicht kennt.}
\linkteil
{\small Alle fünf Tripel:}\\*[0.5ex] \mbox{\hyperlink{spruch:4}{\mycombination[0.5em]{Holz}{Wasser}{Metall}~4}} \hfill \mbox{\hyperlink{spruch:16}{\mycombination[0.5em]{Feuer}{Holz}{Wasser}~16}} \hfill \mbox{\hyperlink{spruch:28}{\mycombination[0.5em]{Erde}{Feuer}{Holz}~28}} \hfill \mbox{\hyperlink{spruch:40}{\mycombination[0.5em]{Metall}{Erde}{Feuer}~40}} \hfill \mbox{\hyperlink{spruch:52}{\mycombination[0.5em]{Wasser}{Metall}{Erde}~52}}\par
\end{orakelblock}
\newpage

\markright{V \quad Die gesammelte Kraft}
\begin{orakelblock}
\musterkopf{V}{Die gesammelte Kraft}{4}{2}{Annehmen}{Bändigen}{Nähren}
\textbf{Du schöpfst aus deiner Quelle, und sie bringt Ruhe in das, was zu heftig drängt. Dann nährst du es, diesmal aus der Fülle statt aus der Not.}\par
\smallskip
\lichtschatten{Du willst loslegen, aber etwas ist zu hitzig, zu eilig. Geh zuerst zurück zu dem, was dich trägt. Dort wird das Drängende still. Danach kannst du ihm geben, was es braucht, und es wächst stetig statt sprunghaft.}
  {Wer immer nur sammelt, findet nie den Moment zu geben. Wer gibt, ohne gesammelt zu haben, nährt die Unruhe statt das Werk.}
\linkteil
{\small Alle fünf Tripel:}\\*[0.5ex] \mbox{\hyperlink{spruch:5}{\mycombination[0.5em]{Holz}{Wasser}{Feuer}~5}} \hfill \mbox{\hyperlink{spruch:17}{\mycombination[0.5em]{Feuer}{Holz}{Erde}~17}} \hfill \mbox{\hyperlink{spruch:29}{\mycombination[0.5em]{Erde}{Feuer}{Metall}~29}} \hfill \mbox{\hyperlink{spruch:41}{\mycombination[0.5em]{Metall}{Erde}{Wasser}~41}} \hfill \mbox{\hyperlink{spruch:53}{\mycombination[0.5em]{Wasser}{Metall}{Holz}~53}}\par
\end{orakelblock}
\newpage

\markright{VI \quad Der Rückhalt}
\begin{orakelblock}
\musterkopf{VI}{Der Rückhalt}{4}{3}{Annehmen}{Sich stellen}{Bändigen}
\textbf{Du lehnst dich an deine Herkunft, und sie hält dem stand, was sie bedrängt. Am Ende gibst du selbst der Lage ihre Form.}\par
\smallskip
\lichtschatten{Du stehst nicht allein. Hinter dir stehen Menschen, Erfahrungen, eine Geschichte, und sie halten stand, wo du allein nachgeben würdest. Mit diesem Rückhalt wird aus Abwehr Gestaltung: Am Ende bestimmst du, wie es weitergeht.}
  {Wer nur mit Rückendeckung handelt, wagt nie etwas Eigenes. Wer jeden Rückhalt ablehnt, kämpft allein und reibt sich auf.}
\linkteil
{\small Alle fünf Tripel:}\\*[0.5ex] \mbox{\hyperlink{spruch:6}{\mycombination[0.5em]{Holz}{Wasser}{Erde}~6}} \hfill \mbox{\hyperlink{spruch:18}{\mycombination[0.5em]{Feuer}{Holz}{Metall}~18}} \hfill \mbox{\hyperlink{spruch:30}{\mycombination[0.5em]{Erde}{Feuer}{Wasser}~30}} \hfill \mbox{\hyperlink{spruch:42}{\mycombination[0.5em]{Metall}{Erde}{Holz}~42}} \hfill \mbox{\hyperlink{spruch:54}{\mycombination[0.5em]{Wasser}{Metall}{Feuer}~54}}\par
\end{orakelblock}
\newpage

\markright{VII \quad Der Damm}
\begin{orakelblock}
\musterkopf{VII}{Der Damm}{2}{1}{Bändigen}{Nähren}{Sich stellen}
\textbf{Du setzt eine Grenze, und hinter ihr sammelt sich etwas Neues. Es wächst, bis es dich selbst herausfordert.}\par
\smallskip
\lichtschatten{Wo du Grenzen setzt, entsteht Raum für Neues: Regeln schaffen Spielräume, Formen bringen Inhalte hervor. Was so entsteht, bleibt nicht klein. Irgendwann stellt es die Grenze infrage, die es hervorgebracht hat. Das ist kein Undank, sondern ein Zeichen, dass es lebt.}
  {Wer seine Grenze gegen alles Neue verteidigt, hält am Ende nur Stillstand fest. Wer jede Grenze aufgibt, sobald sie infrage steht, bringt nichts hervor.}
\linkteil
{\small Alle fünf Tripel:}\\*[0.5ex] \mbox{\hyperlink{spruch:7}{\mycombination[0.5em]{Holz}{Erde}{Metall}~7}} \hfill \mbox{\hyperlink{spruch:19}{\mycombination[0.5em]{Feuer}{Metall}{Wasser}~19}} \hfill \mbox{\hyperlink{spruch:31}{\mycombination[0.5em]{Erde}{Wasser}{Holz}~31}} \hfill \mbox{\hyperlink{spruch:43}{\mycombination[0.5em]{Metall}{Holz}{Feuer}~43}} \hfill \mbox{\hyperlink{spruch:55}{\mycombination[0.5em]{Wasser}{Feuer}{Erde}~55}}\par
\end{orakelblock}
\newpage

\markright{VIII \quad Die Abkürzung}
\begin{orakelblock}
\musterkopf{VIII}{Die Abkürzung}{2}{4}{Bändigen}{Annehmen}{Nähren}
\textbf{Du greifst ein und räumst etwas aus dem Weg. Dahinter liegt, was du ohnehin nähren wolltest, und nun erreichst du es direkt.}\par
\smallskip
\lichtschatten{Manchmal steht etwas zwischen dir und dem, was dir am Herzen liegt. Du räumst es beherzt beiseite. Dahinter wartet keine neue Hürde, sondern genau das, wofür du dich einsetzen wolltest. Der kurze Weg ist frei.}
  {Wer jeden Umweg abkürzt, räumt auch weg, was hätte bleiben dürfen. Wer nie eingreift, nimmt zu jedem Ziel den längsten Weg.}
\linkteil
{\small Alle fünf Tripel:}\\*[0.5ex] \mbox{\hyperlink{spruch:8}{\mycombination[0.5em]{Holz}{Erde}{Feuer}~8}} \hfill \mbox{\hyperlink{spruch:20}{\mycombination[0.5em]{Feuer}{Metall}{Erde}~20}} \hfill \mbox{\hyperlink{spruch:32}{\mycombination[0.5em]{Erde}{Wasser}{Metall}~32}} \hfill \mbox{\hyperlink{spruch:44}{\mycombination[0.5em]{Metall}{Holz}{Wasser}~44}} \hfill \mbox{\hyperlink{spruch:56}{\mycombination[0.5em]{Wasser}{Feuer}{Holz}~56}}\par
\end{orakelblock}
\newpage

\markright{IX \quad Der strenge Weg}
\begin{orakelblock}
\musterkopf{IX}{Der strenge Weg}{2}{2}{Bändigen}{Bändigen}{Annehmen}
\textbf{Ordnung folgt auf Ordnung, Regel auf Regel. Am Ende dieses strengen Weges wartet keine Härte. Dort liegt die Quelle, die dich nährt.}\par
\smallskip
\lichtschatten{Es ist eine Zeit für Disziplin: Tag für Tag dasselbe, sorgfältig und ohne Abkürzung. Das fühlt sich vielleicht eng an. Aber wie bei jeder guten Übung zeigt sich am Ende etwas Weiches: Du findest zu dem zurück, was dich trägt.}
  {Wer Regel auf Regel häuft, vergisst, wozu sie da waren. Wer jede Ordnung meidet, findet nie zu seinem Ursprung.}
\linkteil
{\small Alle fünf Tripel:}\\*[0.5ex] \mbox{\hyperlink{spruch:9}{\mycombination[0.5em]{Holz}{Erde}{Wasser}~9}} \hfill \mbox{\hyperlink{spruch:21}{\mycombination[0.5em]{Feuer}{Metall}{Holz}~21}} \hfill \mbox{\hyperlink{spruch:33}{\mycombination[0.5em]{Erde}{Wasser}{Feuer}~33}} \hfill \mbox{\hyperlink{spruch:45}{\mycombination[0.5em]{Metall}{Holz}{Erde}~45}} \hfill \mbox{\hyperlink{spruch:57}{\mycombination[0.5em]{Wasser}{Feuer}{Metall}~57}}\par
\end{orakelblock}
\newpage

\markright{X \quad Die Läuterung}
\begin{orakelblock}
\musterkopf{X}{Die Läuterung}{3}{1}{Sich stellen}{Nähren}{Annehmen}
\textbf{Du stellst dich dem, was dich begrenzt, und gerade das bringt hervor, was dich nährt. Aus der Krise wird ein Neubeginn.}\par
\smallskip
\lichtschatten{Etwas schränkt dich ein, verletzt dich vielleicht. Du weichst nicht aus. Und dann geschieht das Überraschende: Was dich begrenzt hat, führt dich zu neuer Kraft. Läuterung heißt nicht, dass es nicht wehtut. Es heißt, dass der Schmerz zu etwas wird.}
  {Wer Leid sucht, weil es läutert, verwechselt Schmerz mit Wachstum. Wer jeder Krise ausweicht, bleibt, wie er ist.}
\linkteil
{\small Alle fünf Tripel:}\\*[0.5ex] \mbox{\hyperlink{spruch:10}{\mycombination[0.5em]{Holz}{Metall}{Wasser}~10}} \hfill \mbox{\hyperlink{spruch:22}{\mycombination[0.5em]{Feuer}{Wasser}{Holz}~22}} \hfill \mbox{\hyperlink{spruch:34}{\mycombination[0.5em]{Erde}{Holz}{Feuer}~34}} \hfill \mbox{\hyperlink{spruch:46}{\mycombination[0.5em]{Metall}{Feuer}{Erde}~46}} \hfill \mbox{\hyperlink{spruch:58}{\mycombination[0.5em]{Wasser}{Erde}{Metall}~58}}\par
\end{orakelblock}
\newpage

\markright{XI \quad Die List}
\begin{orakelblock}
\musterkopf{XI}{Die List}{3}{4}{Sich stellen}{Annehmen}{Bändigen}
\textbf{Du stellst dich dem Stärkeren, und er weicht zurück auf seinen eigenen Grund. Dort hast du die Oberhand.}\par
\smallskip
\lichtschatten{Gegen diese Härte kommst du nicht mit Gewalt an. Du hältst stand, ohne anzugreifen, bis sich die Gegenseite auf ihr Vertrautes zurückzieht. Genau dort, auf ihrem eigenen Boden, kannst du ihr Form geben. Klugheit ist auch eine Form von Stärke.}
  {Wer immer über Bande spielt, weiß am Ende selbst nicht mehr, wo er steht. Wer nur frontal angreift, rennt gegen die stärkste Stelle.}
\linkteil
{\small Alle fünf Tripel:}\\*[0.5ex] \mbox{\hyperlink{spruch:11}{\mycombination[0.5em]{Holz}{Metall}{Erde}~11}} \hfill \mbox{\hyperlink{spruch:23}{\mycombination[0.5em]{Feuer}{Wasser}{Metall}~23}} \hfill \mbox{\hyperlink{spruch:35}{\mycombination[0.5em]{Erde}{Holz}{Wasser}~35}} \hfill \mbox{\hyperlink{spruch:47}{\mycombination[0.5em]{Metall}{Feuer}{Holz}~47}} \hfill \mbox{\hyperlink{spruch:59}{\mycombination[0.5em]{Wasser}{Erde}{Feuer}~59}}\par
\end{orakelblock}
\newpage

\markright{XII \quad Der Aufstand}
\begin{orakelblock}
\musterkopf{XII}{Der Aufstand}{3}{3}{Sich stellen}{Sich stellen}{Nähren}
\textbf{Du leistest Widerstand, und dein Beispiel ermutigt andere, ihrerseits Widerstand zu leisten. Am Ende nährst du eine Bewegung, die größer ist als du.}\par
\smallskip
\lichtschatten{Zweimal Widerstand: Das ist der wildeste aller Wege. Du stehst auf, und andere stehen mit dir auf, jeder gegen das, was ihn bedrängt. Daraus kann etwas Neues entstehen, und du hast es genährt. Frag dich, wofür, nicht nur wogegen.}
  {Ein Aufstand, der nur gegen etwas ist, reißt nieder und baut nichts auf. Wer seinen Zorn immer nur herunterschluckt, wird bitter.}
\linkteil
{\small Alle fünf Tripel:}\\*[0.5ex] \mbox{\hyperlink{spruch:12}{\mycombination[0.5em]{Holz}{Metall}{Feuer}~12}} \hfill \mbox{\hyperlink{spruch:24}{\mycombination[0.5em]{Feuer}{Wasser}{Erde}~24}} \hfill \mbox{\hyperlink{spruch:36}{\mycombination[0.5em]{Erde}{Holz}{Metall}~36}} \hfill \mbox{\hyperlink{spruch:48}{\mycombination[0.5em]{Metall}{Feuer}{Wasser}~48}} \hfill \mbox{\hyperlink{spruch:60}{\mycombination[0.5em]{Wasser}{Erde}{Holz}~60}}\par
\end{orakelblock}
\newpage

\addchap{Die 60 Tripel}\markboth{Die 60 Tripel}{\mbox{}}\label{kap:tripel}

\begin{einleitung}
Schließt an ein Paar ein zweites an, das mit dessen zweitem Element beginnt, entsteht ein Tripel: drei verschiedene Elemente in fester Reihenfolge. In ihm kommt alles zusammen, was die vorigen Kapitel beschreiben: Das erste und das zweite Element bilden ein Paar, das zweite und das dritte ein zweites. Die Richtungen dieser beiden Paare ergeben zusammen eine der zwölf Wendungen, und die Wendung verbindet auch Anfang und Ende durch eine Richtung. So steht zum Beispiel im Tripel Holz, Feuer, Erde zweimal das Nähren (\textit{Holz nährt Feuer}, \textit{Feuer nährt Erde}), und vom Holz zur Erde führt das Bändigen: Wendung I, Der Strom.

Fünf Elemente für den Anfang und zwölf Wendungen für den Weg ergeben die 60 Tripel. Zusammen bilden sie ein vollständiges Weltsystem: Jedes Element, jede Richtung, jedes Paar und jede Wendung kommt darin vor, und alle greifen ineinander wie im Wandlungskreis selbst. Jeder Spruch dieses Buches ist eine Facette dieses einen Systems: Jeder zeigt es von einer anderen Seite, im Licht und im Schatten.

Auch die Bilder der Sprüche folgen diesem System. Jedes Paar hat einen eigenen Vorgang aus der Natur: Holz brennt, der Deich hält das Wasser, die Axt fällt den Baum. Jede Wendung beschreibt eine Lage, wie Menschen sie erleben: die Prüfung, die Unterweisung, den Aufstand. Ein Tripel fügt die Vorgänge seiner beiden Paare zu einer Szene zusammen, und die Wendung gibt ihr den Sinn. Wer die Paare kennt, erkennt sie in den Tripeln wieder.

Die Tripel sind nach ihrem ersten Element geordnet, und innerhalb jeder Gruppe folgen sie den Wendungen I bis XII.
\end{einleitung}

\newpage
\markright{\element{Holz}}
\spruch{1}{Holz}{Feuer}{Erde}{I Strom}{Das Haus}{Dein Holz heizt den Ziegelofen, und im Feuer wird aus Lehm ein Backstein. Das Gebälk fasst die Ziegel: So entsteht ein Haus.}
  {Du hast etwas angestoßen, und es hat andere entzündet. Daraus ist etwas Festes geworden, ein Ort, an dem Menschen zusammenkommen und bleiben können. Du musst nicht mehr alles selbst tragen. Aber was hier entsteht, folgt deinem Plan, und du gibst ihm den Rahmen.}
  {Ein Fachwerk ohne Gefache ist nur ein Gerippe. Gefache ohne Balken fallen in sich zusammen.}
\newpage
\spruch{2}{Holz}{Feuer}{Metall}{II Prüfung}{Das zweischneidige Schwert}{Als Holzkohle bringst du die Esse zum Glühen, und das Eisen wird weich. Was auf dem Amboss daraus geschmiedet wird, ist ein zweischneidiges Schwert, und am Ende spiegelst du dich in seiner Klinge.}
  {Was du hervorgebracht hast, ist stark geworden, vielleicht stärker, als du wolltest. Es hat eine eigene Schärfe und schneidet nach beiden Seiten, auch zu dir hin. Das ist die Prüfung jedes Schöpfers: zu seinem Werk zu stehen, gerade wenn es ihm fremd wird. Wer es mit Achtung in die Hand nimmt, kann es führen.}
  {Wer vor seinem Werk zurückschreckt und es verstößt, macht es erst zum Feind, wie Viktor Frankenstein sein Geschöpf. Wer sich in sein Werk verliebt wie Pygmalion in seine Statue, sieht nicht, wo es schneidet.}
\newpage
\spruch{3}{Holz}{Feuer}{Wasser}{III Unterweisung}{Das Feuer im Regen}{Du legst Holz nach, damit das Feuer im Regen weiterbrennt. Derselbe Regen tränkt am Ende deine eigenen Wurzeln.}
  {Du hilfst jemandem, an etwas festzuhalten, das ihm Freude macht, auch gegen Zweifel und kalte Worte. Dabei gibst du von dir. Doch was der andere durchsteht, kommt auch dir zugute: Du lernst neu, was dich selbst trägt. Wer andere unterweist, wird selbst genährt.}
  {Wer nur nachlegt, verbrennt sich für ein fremdes Feuer. Wer den Regen von allen fernhält, lässt auch die eigenen Wurzeln vertrocknen.}
\newpage
\spruch{4}{Holz}{Wasser}{Metall}{IV Wurzelsuche}{Die Wurzel am Fels}{Deine Wurzeln folgen dem Wasser in die Tiefe, bis zum Fels, aus dem es kommt. Dort stößt du auf das Harte, das dich begrenzt.}
  {Du suchst nach dem, was dich nährt, und gehst dabei tiefer als gedacht. Unter deiner Herkunft liegt etwas Älteres und Härteres: eine Regel, ein Urteil, ein Verbot. Weich nicht zurück. Eine Wurzel, die den Fels erreicht, sprengt ihn nicht sofort, aber sie findet seine Spalten.}
  {Wer zu tief gräbt, findet nur noch Stein. Wer an der Oberfläche bleibt, weiß nie, worauf er steht.}
\newpage
\spruch{5}{Holz}{Wasser}{Feuer}{V Gesammelte Kraft}{Der Wald nach dem Regen}{Ein Wald, der gut im Saft steht, fängt nicht beim ersten Funken Feuer. Erst wenn sein Holz geschlagen und getrocknet ist, gibt es Wärme, die lange hält.}
  {Etwas in dir will lodern, sofort und mit aller Kraft. Doch du gehst zuerst zurück zu dem, was dich nährt, und dort kommt deine Ungeduld zur Ruhe. Das ist kein Verzicht. Was du danach gibst, brennt ruhiger und länger.}
  {Holz, das nie geschlagen wird, wärmt niemanden. Wer sein Holz grün verfeuert, erntet Rauch statt Glut.}
\newpage
\spruch{6}{Holz}{Wasser}{Erde}{VI Rückhalt}{Die Uferweide}{Die Weide trinkt aus dem Fluss, und der Fluss nagt am Ufer. Doch ihre Wurzeln halten das Ufer fest, und am Ende bestimmt sie, wo das Land bleibt.}
  {Deine Herkunft nährt dich, und sie hat es selbst nicht leicht: An ihren Rändern bröckelt es, alte Gewissheiten werden fortgespült. Du hältst dagegen. Weil du so tief in ihr wurzelst, kannst du festhalten, was bleiben soll, und gibst dem Ufer seine Form.}
  {Wo jede Weide gefällt wird, reißt das Hochwasser das Ufer mit. Wo zu viele stehen, verlandet der Fluss.}
\newpage
\spruch{7}{Holz}{Erde}{Metall}{VII Damm}{Der Holzkeil}{Du treibst einen Keil aus Holz in den Fels und machst ihn nass. Das Holz quillt, der Stein springt, und im Spalt glänzt eine Erzader, härter als dein Keil.}
  {Du hast geduldig gearbeitet und Festes aufgebrochen. Nun liegt vor dir, was darin verborgen war: klar, wertvoll, eigenwillig. Es fragt dich, wozu du das alles getan hast. Weich der Frage nicht aus. Du kannst sie beantworten, und dein Werk wird dadurch nur besser.}
  {Wer jeden Fund wieder zuschüttet, weil er ihn stört, findet nur das Gewohnte. Wer sich von seinem Fund beherrschen lässt, vergisst, wofür er den Fels gespalten hat.}
\newpage
\spruch{8}{Holz}{Erde}{Feuer}{VIII Abkürzung}{Die Tongrube}{Du gräbst den Boden auf und findest Ton. Dein Holz brennt ihn zum Gefäß, und so nährst du ohne Umweg, was dir am Herzen liegt.}
  {Etwas Festgefahrenes stand zwischen dir und dem, was du wolltest. Du hast es beherzt aufgebrochen, und darunter lag genau der Stoff, den du brauchst. Jetzt kannst du deine Kraft direkt hineingeben: Deine Mühe schürt die Glut, und in ihr wird das Gefäß gebrannt.}
  {Wer zu gierig gräbt, reißt auch fruchtbaren Boden mit heraus. Wer zu zaghaft gräbt, kommt nie bis zum Ton.}
\newpage
\spruch{9}{Holz}{Erde}{Wasser}{IX Strenger Weg}{Die Furchen}{Du ziehst Furche um Furche durch den Acker, gerade und mit gleichmäßigem Gefälle. In ihnen läuft das Wasser genau dorthin, wo es gebraucht wird, und am Ende trinken deine Wurzeln.}
  {Es ist mühsame Arbeit, die keine Nachlässigkeit erlaubt: Ist eine Furche zu steil, reißt das Wasser den Boden mit; ist sie zu flach, bleibt es stehen. Doch die strenge Ordnung hat einen Sinn. Wenn alles bestellt ist, merkst du: Sie hat dich nicht eingeengt, sie hat dir das Wasser gebracht.}
  {Wer jeden Tropfen in Kanäle zwingt, hat bald mehr Gräben als Pflanzen. Wer das Wasser laufen lässt, wie es will, sieht es versickern oder den Boden fortspülen.}
\newpage
\spruch{10}{Holz}{Metall}{Wasser}{X Läuterung}{Die Schleuse}{Dein Boot liegt vor dem eisernen Tor der Schleuse, und es geht nicht weiter. Dann strömt durch die Schütze im Tor Wasser herein und hebt dich auf eine neue Höhe.}
  {Vor dir steht eine Grenze, hart und unverrückbar, und du kannst nichts tun als warten und standhalten. Doch gerade diese Grenze gibt dir, was du brauchst: Das Wasser, das sie zurückgehalten hat, trägt dich nach oben. Wenn sich das Tor öffnet, fährst du auf einer neuen Höhe weiter. Was dich aufgehalten hat, hat dich gehoben.}
  {Wer vor dem Tor ungeduldig dagegen anrudert, kommt keinen Meter weiter. Wer die Schleuse fürchtet, bleibt für immer im unteren Becken.}
\newpage
\spruch{11}{Holz}{Metall}{Erde}{XI List}{Der eingewachsene Zaun}{Ein Baum wächst an einem eisernen Zaun empor. Er kämpft nicht gegen den Draht, er umwächst ihn. Mit den Jahren zerfällt das Eisen im Boden zu Rost, und die Wurzeln halten den Grund, auf dem einmal die Grenze stand.}
  {Gegen diese Härte kommst du mit Gewalt nicht an, und du musst es auch nicht. Du wächst weiter, langsam und beharrlich, um das Hindernis herum, bis es ein Teil von dir ist. Was dich einengen sollte, verliert mit der Zeit seine Macht. Am Ende gehört der Boden dir. Hier ist Geduld die klügste List.}
  {Wer jedes Hindernis umwächst, trägt am Ende viel fremdes Eisen in sich. Wer gegen den Draht anwächst, statt ihn zu umfassen, schnürt sich selbst ab.}
\newpage
\spruch{12}{Holz}{Metall}{Feuer}{XII Aufstand}{Die heiße Säge}{Das Holz hält gegen die Säge, und die Säge hält dagegen. Zwischen Blatt und Stamm wird es so heiß, dass das Holz zu rauchen beginnt. Aus Widerstand gegen Widerstand entsteht Hitze, und das Holz nährt ein Feuer.}
  {Man will dich passend machen, und du leistest Widerstand. Die Gegenseite gibt ebenfalls nicht nach. Aus dieser Reibung entsteht Hitze, und sie greift über: Andere entzünden sich an deinem Widerstand. Frag dich, wofür dieses Feuer brennen soll, nicht nur wogegen.}
  {Wer sich zu lange aneinander reibt, hinterlässt Brandspuren und verdirbt die Klinge. Wer jede Reibung meidet, bleibt kalt.}

\newpage
\markright{\element{Feuer}}
\spruch{13}{Feuer}{Erde}{Metall}{I Strom}{Der Glockenguss}{Deine Glut brennt die Form aus Lehm und macht die Bronze flüssig. In der Form nimmt das Metall Gestalt an, und wenn man sie zerschlägt, steht eine Glocke da, die Jahrhunderte überdauert.}
  {Was du mit Leidenschaft begonnen hast, hat eine feste Form gefunden, und in ihr wird etwas Bleibendes daraus. Das gelingt nicht nebenbei: Hitze, Form und Metall müssen genau zusammenpassen, und ein Fehler zeigt sich erst, wenn die Form zerbrochen ist. Doch was so entsteht, überdauert dich. Du stehst am Anfang und am Ende.}
  {Wer zu hastig gießt, bekommt eine Glocke mit Sprung. Wer die Form nicht zerschlagen will, bekommt sein Werk nie zu sehen.}
\newpage
\spruch{14}{Feuer}{Erde}{Wasser}{II Prüfung}{Der Backofen}{Deine Glut härtet den Ofen aus Lehm. Im Teig hält das Mehl das Wasser in einem zähen Netz. Wird der Teig eingeschoben, muss deine Hitze bis in die feuchte Mitte dringen; erst wenn das Brot beim Klopfen hohl klingt, ist es durch.}
  {Du hast viel Wärme gegeben und damit etwas Festes geschaffen, das anderen dient. Jetzt kommt es auf dich selbst an: Das Feuchte und Schwere, das in jeder Sache steckt, will durchgebacken werden, und das kostet Hitze und Geduld. Halte die Hitze im rechten Maß, nicht zu stark und nicht zu schwach. Ob es gelungen ist, zeigt erst die Probe.}
  {Ist der Ofen zu heiß, verbrennt die Kruste, und innen bleibt das Brot klitschig. Ist er zu kühl, läuft der Teig breit und bleibt blass und zäh.}
\newpage
\spruch{15}{Feuer}{Erde}{Holz}{III Unterweisung}{Der Kohlenmeiler}{Unter einer Decke aus Erde glimmst du im gestapelten Holz. Die Erde hält es zusammen und lässt nur wenig Luft herein. So zehrst du langsam vom Holz, Tag um Tag, bis es zu Kohle geworden ist.}
  {Diese Arbeit lernt man nicht aus Büchern, sondern am Meiler, Tag und Nacht neben einem, der es kann. Du gibst weiter, was du weißt: wie man die Decke schließt, wann man Luft gibt, woran man am Rauch erkennt, dass die Kohle gar ist. Und während du es erklärst, siehst du dein eigenes Handwerk neu. Wer unterweist, lernt mit.}
  {Bekommt der Meiler zu viel Luft, brennt das Holz zu Asche. Ist die Decke zu dicht, erstickt die Glut, und das Holz bleibt roh.}
\newpage
\spruch{16}{Feuer}{Holz}{Wasser}{IV Wurzelsuche}{Die Jahresringe}{Du ergreifst ein Scheit und siehst seine Jahresringe: breite aus nassen Jahren, schmale aus trockenen. Ohne Regen wäre es nie gewachsen, und noch jetzt zischt das Wasser in ihm, wenn du es verzehrst.}
  {Du fragst, woher deine Kraft kommt, und gehst ihr nach, weiter zurück als je zuvor. Hinter dem, was dich entzündet hat, findest du ausgerechnet dein Gegenteil: das Stille und Kühle, das dir fremd ist. Ohne es gäbe es dich nicht. Begegne ihm mit Respekt. Gegensätze sind oft enger verbunden, als sie wissen.}
  {Wer sein Gegenteil verachtet, vergisst, wovon er lebt. Wer ihm zu nahe kommt, erlischt.}
\newpage
\spruch{17}{Feuer}{Holz}{Erde}{V Gesammelte Kraft}{Die Herdasche}{Den Winter über zehrst du vom Holz. Im Frühjahr bricht der Keim durch die verkrustete Erde, und deine Asche gibt ihm Kraft.}
  {Bevor du Wärme verschenkst, sammelst du dich. Du nimmst auf, was dich nährt, und dieselbe Kraft bricht auf, was zu hart und fest geworden ist. Erst dann gibst du. Deine Wärme fällt jetzt auf vorbereiteten Boden und macht ihn fruchtbar.}
  {Wer seine Asche auf harten Grund streut, sieht sie vom Wind verweht. Wer sie im Herd lässt, düngt nichts.}
\newpage
\spruch{18}{Feuer}{Holz}{Metall}{VI Rückhalt}{Der Ambossstock}{Dein Feuer nährt sich vom Holz, und auf einem Holzklotz steht der Amboss, der jedem Schlag standhält. So getragen, machst du das Eisen weich und bestimmst seine Form.}
  {Du kannst dich auf das stützen, was dich entzündet hat: Menschen, Ideen, eine Herkunft. Sie fangen die harten Schläge ab, die dich sonst treffen würden. Mit diesem Rückhalt kannst du dich dem Starren zuwenden, und in deiner Glut gibt es nach.}
  {Wer alle Schläge auf andere abwälzt, schmiedet auf deren Kosten. Wer keinen Rückhalt annimmt, spürt jeden Hammerschlag am eigenen Leib.}
\newpage
\spruch{19}{Feuer}{Metall}{Wasser}{VII Damm}{Der Kessel}{Aus dem Eisen, das in deinem Feuer weich wurde, wird ein Kessel geschmiedet, und er füllt sich mit Wasser. Nun hängt er über deinem eigenen Feuer, und wenn er überkocht, zischt es in der Glut.}
  {Du hast etwas Starres mit Leidenschaft in Form gebracht, und daraus ist ein Gefäß geworden, das etwas Neues aufnimmt. Dieses Neue ist kühler, ruhiger, eigensinnig. Es lebt von deiner Wärme und kann sie doch dämpfen. Gib ihm deine Wärme, aber lass dich nicht löschen.}
  {Hängt der Kessel zu tief, kocht er über und löscht die Glut. Hängt er zu hoch, wird das Wasser nie warm.}
\newpage
\spruch{20}{Feuer}{Metall}{Erde}{VIII Abkürzung}{Die Pflugschar}{Du wirfst das Schwert ins Feuer und machst daraus eine Pflugschar. Der edle Stahl geht in die Erde: Was im Feuer neu entstanden ist, dient nun dem Acker.}
  {Etwas Hartes stand zwischen dir und dem, was dir am Herzen liegt: ein Streit, eine Kränkung, eine Waffe in fremder oder eigener Hand. Du hast es eingeschmolzen und neu geformt. Dabei gibst du etwas auf, denn ein Schwert ist aus feinerem Stahl als jede Pflugschar. Doch nun dient es nicht mehr dem Kampf, sondern dem Wachsen.}
  {Wer jede Klinge umschmiedet, steht schutzlos da, wenn er eine braucht. Wer keine umschmiedet, pflügt nie.}
\newpage
\spruch{21}{Feuer}{Metall}{Holz}{IX Strenger Weg}{Der Holzstoß}{Das Eisen glüht in deinem Feuer, auf dem Amboss wird es zur Axt geschmiedet und gehärtet. Im Winter fällt die Axt die Bäume, die Stämme werden entastet, gespalten und gestapelt. Erst wenn das Holz lange genug getrocknet ist, nährt es dein Feuer.}
  {Jetzt ist Arbeit gefragt, und sie verlangt einen langen Atem: erst das Werkzeug, dann Fällen, Entasten, Spalten, Stapeln, und dann das Warten. Nichts davon lässt sich abkürzen. Früher hing das Leben daran, ob der Vorrat für den Winter reichte. Wer so vorsorgt, hat am Ende genug, um lange zu brennen.}
  {Wer nur stapelt und sich nie ans Feuer setzt, friert neben dem vollen Holzstoß. Wer seinen Vorrat zu früh verfeuert, friert, wenn der Winter am härtesten ist.}
\newpage
\spruch{22}{Feuer}{Wasser}{Holz}{X Läuterung}{Das Sommergewitter}{Der Regen fällt auf deine Flamme, und du brennst weiter. Derselbe Regen lässt die Bäume wachsen, die dich in späteren Jahren nähren.}
  {Etwas Kaltes trifft dich: Ernüchterung, Kritik, ein Rückschlag. Du stellst dich ihm und erlischst nicht. Mit der Zeit zeigt sich, dass genau diese Kälte etwas Neues hat wachsen lassen, an dem du dich wieder entzünden kannst. Was dich dämpfen wollte, schenkt dir einen Neubeginn.}
  {Wer jede Kälte als Verrat empfindet, sieht nicht, was in ihr wächst. Wer beim ersten Regen aufgibt, ist nicht mehr da, wenn die Bäume groß sind.}
\newpage
\spruch{23}{Feuer}{Wasser}{Metall}{XI List}{Das Feuersetzen}{Vor dem Fels brennt ein Feuer, bis der Stein durch und durch heiß ist. Dann wird Wasser darauf gegossen: Der Fels springt, und das Erz liegt offen.}
  {Gegen diese Härte kommst du mit Hitze allein nicht an. Doch du weichst der Kälte nicht aus, die dir entgegenschlägt: Du nutzt sie. Wo Hitze und Kälte aufeinandertreffen, bekommt selbst Fels Risse. So erreichst du, was dir verschlossen schien.}
  {Wer immer nur heizt, bringt den Fels nicht zum Springen. Wer das Wasser zu früh gießt, löscht nur sein Feuer.}
\newpage
\spruch{24}{Feuer}{Wasser}{Erde}{XII Aufstand}{Der Vulkan}{Glut bricht aus dem Berg und stellt sich dem Meer, und das Wasser frisst sich in die junge Erde. Am Ende ist kein Boden fruchtbarer als die Asche des Vulkans.}
  {Du lehnst dich auf, und was dir entgegensteht, lehnt sich ebenfalls auf. Es wird laut und heiß, und nichts bleibt, wie es war. Doch aus diesem Aufruhr entsteht neues Land. Wer nach dem Ausbruch Geduld hat, erntet auf diesem Boden reicher als je zuvor.}
  {Ein Ausbruch ohne Ende begräbt alles unter sich. Wer seine Glut immer zurückhält, wird hart wie erkaltetes Gestein.}

\newpage
\markright{\element{Erde}}
\spruch{25}{Erde}{Metall}{Wasser}{I Strom}{Der Brunnen}{In deinem Boden liegt Gestein, und das Wasser, das hindurchsickert, wird darin klar und reich an Mineralien. In einem gemauerten Schacht fasst du es: So wird ein Brunnen daraus.}
  {Was du mit Geduld gepflegt hast, ist zu etwas Klarem geworden, und daraus fließt nun etwas Eigenes: Ideen, Gedanken, vielleicht Menschen, die eigene Wege gehen. Du musst es nicht festhalten. Aber du kannst ihm eine Fassung geben, damit andere daraus schöpfen können.}
  {Ein Brunnen, aus dem zu viel geschöpft wird, fällt trocken. Einer, den niemand benutzt, verschlammt.}
\newpage
\spruch{26}{Erde}{Metall}{Holz}{II Prüfung}{Die Quecke}{Aus deinem Erz wird eine Hacke, und sie räumt das Beet vom Unkraut. Doch die Quecke treibt ihre Wurzeln weiter durch deinen Boden, und mit ihr musst du selbst fertig werden.}
  {Du hast unterstützt, was Ordnung schafft, und es hat aufgeräumt, wo zu viel wucherte. Nun bleibt der Rest bei dir: das Zähe, das sich nicht einfach weghacken lässt. Das ist deine Prüfung. Gründlichkeit und Geduld helfen hier mehr als Gewalt.}
  {Wer die Quecke mit der Hacke zerkleinert, vermehrt sie, denn aus jedem Stück treibt eine neue. Wer sie gewähren lässt, hat bald nur noch Quecke im Beet.}
\newpage
\spruch{27}{Erde}{Metall}{Feuer}{III Unterweisung}{Der Kachelofen}{Aus deinem Erz wird die eiserne Tür des Kachelofens, und sie hält der Glut stand. Die Kacheln aus gebranntem Ton nehmen die Hitze auf und geben sie noch Stunden später als Wärme ab.}
  {Du hast einem anderen mitgegeben, was du hast, und er stellt sich nun einer harten Probe. Du bist dabei und hältst die Hitze mit aus. Was er dabei durchsteht, geht auch in dich über und bleibt in dir, lange nachdem die Probe vorbei ist. Wer unterweist, wird mitgeprüft.}
  {Kacheln, die zu schnell aufgeheizt werden, springen. Ein Ofen, der nie angeheizt wird, bleibt ein kalter Klotz.}
\newpage
\spruch{28}{Erde}{Feuer}{Holz}{IV Wurzelsuche}{Der Brennofen}{Du fragst, wer dich gebrannt hat, und findest das Feuer. Du fragst, was das Feuer genährt hat, und findest das Holz. Es wuchs einst aus deinem eigenen Boden, und seine Wurzeln stecken noch heute in dir.}
  {Was dich geformt hat, hatte selbst eine Herkunft. Hinter der Glut, die dich gefestigt hat, liegt etwas, das in dich hineingewachsen ist und dich bis heute durchdringt: ein Anspruch, ein Ehrgeiz, ein fremder Plan. Stell dich ihm. Du kannst ihn halten, ohne dich von ihm sprengen zu lassen.}
  {Ein Topf, in dem die Wurzeln ungehindert wachsen, springt. Wer keine Wurzel in sich duldet, bleibt leer.}
\newpage
\spruch{29}{Erde}{Feuer}{Metall}{V Gesammelte Kraft}{Der Schmelztiegel}{Im Feuer wirst du zum Tiegel, fest und hitzebeständig. Das Feuer macht das harte Metall flüssig, und du hältst es, bis es gegossen wird.}
  {Bevor du dich um etwas Schwieriges kümmerst, holst du dir Wärme: Zuspruch, Nähe, Freude. Diese Wärme macht weich, was zu hart war, an dir und an der Lage. Dann kannst du tragen und halten, was jetzt Form annehmen will. Du gibst aus der Fülle.}
  {Ein Tiegel, der nie im Feuer war, zerbricht beim ersten Guss. Einer, der zu oft im Feuer war, wird rissig.}
\newpage
\spruch{30}{Erde}{Feuer}{Wasser}{VI Rückhalt}{Der Klinker}{Im Feuer gebrannt, wird dein Ton zu Klinker, und der Regen kann ihm nichts mehr anhaben. So gestärkt, hältst du das Wasser in seinen Grenzen.}
  {Hinter dir steht, was dich gestärkt hat: Wärme, Begeisterung, Menschen, die an dich glauben. Sie halten der Kälte stand, die dich sonst aufweichen würde. Mit diesem Rückhalt kannst du Grenzen ziehen und Fließendes in seine Bahn lenken.}
  {Ungebrannter Ton weicht im ersten Regen auf. Zu heiß gebrannter verzieht sich und taugt nicht mehr für die Mauer.}
\newpage
\spruch{31}{Erde}{Wasser}{Holz}{VII Damm}{Der Baum auf dem Deich}{Dein Deich hält das Wasser, und das Wasser lässt Bäume wachsen, auch auf dem Deich selbst. Ihre Wurzeln dringen in dich ein, und du musst ihnen standhalten.}
  {Du hast eine Grenze gezogen, und in ihrem Schutz ist Leben entstanden. Nun wächst es weiter, in dich hinein, und stellt deine Ordnung infrage. Das ist kein Undank. Halte stand, aber prüfe auch, welche Wurzeln dich festigen und welche dich aushöhlen.}
  {Ein Deich, auf dem Bäume ungehindert wachsen, wird von ihren Wurzeln ausgehöhlt. Einer, auf dem nicht einmal Gras wurzeln darf, wird vom Regen abgetragen.}
\newpage
\spruch{32}{Erde}{Wasser}{Metall}{VIII Abkürzung}{Die Goldwäsche}{Du lenkst das Wasser durch eine Rinne aus Erde, und es spült den Sand fort bis auf den Fels. Dort liegt das Gold, das dein Boden hervorgebracht hat.}
  {Du greifst ein und gibst dem Fließenden eine Richtung. Das Wasser tut den Rest: Es trägt fort, was leicht ist, und lässt liegen, was Gewicht hat. Was übrig bleibt, ist das Wertvolle, das immer schon in dir lag. Der kürzeste Weg dorthin führte durch das Wasser.}
  {Wer zu viel Wasser durch die Rinne jagt, spült das Gold mit fort. Wer gar keins hindurchlässt, findet im Schlamm nichts.}
\newpage
\spruch{33}{Erde}{Wasser}{Feuer}{IX Strenger Weg}{Die Töpferei}{Du hältst das Wasser im Ton, knetest ihn und formst ihn auf der Scheibe. Solange er nass ist, darf kein Feuer an ihn heran. Erst am Ende dieser langen Arbeit wird er im Brand zum Gefäß.}
  {Jetzt gilt Sorgfalt vor Begeisterung: kneten, auf der Scheibe formen, trocknen lassen, Schritt für Schritt. Jede Abkürzung würde sich rächen. Doch am Ende dieses strengen Weges kommt das Feuer zu dir, und du wirst, was du werden sollst: ein Gefäß, das hält.}
  {Wer den Ton zu früh brennt, sieht ihn im Ofen springen. Wer ihn nie brennt, hat nur getrockneten Schlamm.}
\newpage
\spruch{34}{Erde}{Holz}{Feuer}{X Läuterung}{Der Wurzelstock im Ofen}{Du hältst den Wurzeln stand, die in dich eindringen. Ihr Holz brennt schließlich im Ofen, und dich brennt es zu einem festen Gefäß.}
  {Etwas Neues drängt in dein Leben, zerrt an deiner Ordnung, nimmt dir die Ruhe. Du stellst dich ihm. Dann geschieht das Unerwartete: Gerade das, was dich bedrängt hat, wird zum Feuer, das dich härtet. Läuterung heißt, durchs Feuer zu gehen und dabei fester zu werden.}
  {Wer jedes Eindringen fürchtet, bleibt weicher Lehm. Wer sich allem öffnet, wird durchwurzelt und zerfällt.}
\newpage
\spruch{35}{Erde}{Holz}{Wasser}{XI List}{Die Weiden im nassen Grund}{Du hältst die Wurzeln der Weiden fest, die in dich hineinwachsen, und sie ziehen das Wasser heraus, das dich aufweichen würde. So wird der nasse Grund mit der Zeit fest.}
  {Was dich bedrängt, kannst du nicht vertreiben. Also hältst du es fest und lässt es für dich arbeiten. Es nimmt dir ab, was dich schwer und schwammig machte. Klugheit heißt hier: den Gegner nicht bekämpfen, sondern ihm eine Aufgabe geben.}
  {Zu viele Weiden saugen den Boden trocken. Ohne sie bleibt er ein Sumpf.}
\newpage
\spruch{36}{Erde}{Holz}{Metall}{XII Aufstand}{Die Kriegssense}{Der Boden hält die alten Wurzelstöcke fest, und an ihnen wird die Sense stumpf. Da lassen die Bauern sie gerade schmieden und ziehen mit ihr los: Der Acker hat seine eigene Waffe hervorgebracht.}
  {Lange hast du getragen und ausgehalten, und was in dich hineinwuchs, hast du festgehalten. Nun stößt die Härte, die über dich hinwegfahren will, auf Widerstand, und er ist zäher als sie. Aus dem, was du immer schon hattest, machst du das Werkzeug deines Aufstands, wie die Bauern im Bauernkrieg aus ihren Sensen.}
  {Wer jede Sense zur Waffe schmiedet, hat am Ende nichts mehr, womit er ernten kann. Wer sich nie widersetzt, wird gemäht wie Gras.}

\newpage
\markright{\element{Metall}}
\spruch{37}{Metall}{Wasser}{Holz}{I Strom}{Die Baumschule}{Aus deiner blechernen Kanne fließt Wasser auf die jungen Bäume. Wenn sie gewachsen sind, stutzt du sie mit der Schere in Form.}
  {Aus deiner Klarheit ist Ruhe geworden, und aus der Ruhe wächst Neues, das nicht mehr dir gehört. Es entwickelt sich nach eigenen Plänen. Doch weil du am Anfang standest, hast du das Recht und die Aufgabe, ihm das Maß vorzugeben.}
  {Wer alles stutzt, was er hat wachsen lassen, zieht nur Hecken. Wer nie eingreift, erkennt seine Bäume bald nicht wieder.}
\newpage
\spruch{38}{Metall}{Wasser}{Feuer}{II Prüfung}{Der Löschtrog}{Im eisernen Trog neben der Esse steht Wasser; damit hält der Schmied das Feuer im Zaum. Doch am Ende muss das Eisen selbst in die Glut, und es zeigt sich, ob es standhält.}
  {Du hast für Ruhe gesorgt, wo es hitzig wurde, und anderen geholfen, die Fassung zu bewahren. Jetzt kommt die Hitze zu dir: Leidenschaft, Druck, ein Konflikt, der dich selbst betrifft. Das ist deine Prüfung. Bleib bei dir. Eisen, das in die Glut muss, wird nicht schwächer, sondern brauchbarer.}
  {Wer jede Hitze von sich fernhält, wird nie geschmiedet. Wer zu lange in der Glut bleibt, verbrennt.}
\newpage
\spruch{39}{Metall}{Wasser}{Erde}{III Unterweisung}{Das Erz im Flussbett}{Aus deinem Fels rinnt ein Bach, und der Bach gräbt sich durch die Erde. In der Schlucht, die er schneidet, liegt das Gestein offen, aus dem du stammst.}
  {Du gibst weiter, was du klar erkannt hast, und jemand anderes macht daraus seinen eigenen Weg. Er stößt auf Widerstände und arbeitet sich durch. Dabei legt er frei, was du selbst vergessen hattest: woher du kommst und was dich getragen hat.}
  {Ein Lehrer, der nur Antworten gibt, sieht nie, was seine Schüler freilegen. Einer, der nichts weitergibt, vergisst, woher er kommt.}
\newpage
\spruch{40}{Metall}{Erde}{Feuer}{IV Wurzelsuche}{Die Glut im Berg}{Du folgst dem Erz zurück in den Berg und dem Berg zurück in die Tiefe. Dort unten glüht es noch, und du stehst der Hitze gegenüber, aus der alles kam.}
  {Du suchst nach deinem Ursprung, und hinter dem festen Grund findest du etwas Heißes, Ungebändigtes: eine alte Leidenschaft, einen Zorn, eine Liebe, die du lange kühl gehalten hast. Stell dich ihr. Du hältst mehr Hitze aus, als du glaubst, und erkennst dich darin wieder.}
  {Wer immer tiefer in die Glut hinabsteigt, schmilzt. Wer sie verleugnet, bleibt kalt und spröde.}
\newpage
\spruch{41}{Metall}{Erde}{Wasser}{V Gesammelte Kraft}{Die Tropfsteinhöhle}{Du kehrst zurück in den Berg. Er lässt das Regenwasser nur langsam hindurch, und auf seinem Weg durch dein Gestein nimmt es Kalk auf. Tropfen für Tropfen baut es daraus in der Stille Tropfsteine.}
  {Bevor du sprichst oder entscheidest, ziehst du dich an einen festen Ort zurück. Dort beruhigt sich, was in dir und um dich zu heftig strömt. Was du danach gibst, kommt Tropfen für Tropfen, aber es ist klar und hat Bestand.}
  {Wer zu lange im Berg bleibt, gibt nur noch Tropfen, wo ein Strom gebraucht würde. Wer nie hineingeht, hat nichts Klares zu geben.}
\newpage
\spruch{42}{Metall}{Erde}{Holz}{VI Rückhalt}{Die Fallkerbe}{Das Eisen deiner Axt kommt aus dem Berg, und du stehst fest auf deinem Grund. Der Boden hält den Baum, bis die Axt ihn trifft; mit der Fallkerbe bestimmst du, wohin er fällt.}
  {Du stützt dich auf das, woher du kommst: auf Menschen, die dich getragen haben, und auf eine Erfahrung, die dir niemand nehmen kann. Dieser Grund hält stand gegen alles, was wild in ihn hineinwächst. So gestützt, kannst du die Axt sicher führen und entscheiden, wohin es fällt.}
  {Wer ohne Fallkerbe fällt, wird vom eigenen Baum getroffen. Wer nur seinen Stand prüft und nie zuschlägt, fällt keinen Baum.}
\newpage
\spruch{43}{Metall}{Holz}{Feuer}{VII Damm}{Die schartige Axt}{Deine Axt fällt Bäume, und aus dem Holz brennt der Köhler Kohle. Wird die Axt stumpf und schartig, muss sie selbst in diese Glut, bevor der Schmied sie neu ausschmieden kann.}
  {Du hast klar entschieden und beendet, was zu viel war. Daraus ist etwas Lebendiges geworden, das sich nicht mehr in deine Ordnung fügt. Irgendwann musst du selbst in dieses Feuer, um wieder scharf zu werden. Halte stand. Es ist dein eigenes Werk.}
  {Wer jedes Feuer fürchtet, das er selbst entfacht hat, wird stumpf. Wer zu lange in der Glut bleibt, verliert seine Härte.}
\newpage
\spruch{44}{Metall}{Holz}{Wasser}{VIII Abkürzung}{Der Einbaum}{Deine Axt höhlt einen Stamm aus, und der Einbaum gleitet ins Wasser. Er saugt sich voll, bis er dicht ist, und trägt dich ohne Umweg ans andere Ufer.}
  {Etwas Sperriges lag dir im Weg. Mit Axt und Dechsel hast du daraus ein Fahrzeug gemacht, und nun trägt es dich dorthin, wohin du wolltest: in die Stille, in die Tiefe, ans andere Ufer. Der kürzeste Weg führte nicht um das Hindernis herum, sondern mitten hindurch.}
  {Wer den Stamm zu dünn aushöhlt, schlägt beim ersten Stoß leck. Ein Einbaum, der zu lange an Land liegt, trocknet aus und reißt.}
\newpage
\spruch{45}{Metall}{Holz}{Erde}{IX Strenger Weg}{Der Stampflehm}{Deine Axt macht aus Stämmen Bretter, und zwischen den Brettern wird Lehm festgestampft, Schicht um Schicht. Fällt die Schalung, steht eine Wand aus Erde, die dich schützt.}
  {Eine Ordnung folgt auf die andere: erst die Bretter, dann der Lehm, Tag für Tag, Schicht um Schicht. Das ist mühsam und lässt wenig Raum für Eigenes. Doch am Ende dieses Weges steht etwas, das dich trägt und schützt, und du bist heimgekehrt zu dem, woraus du kommst.}
  {Ist der Lehm zu nass, reißt die Wand beim Trocknen. Ist er zu trocken, bindet er nicht.}
\newpage
\spruch{46}{Metall}{Feuer}{Erde}{X Läuterung}{Die Schlacke}{Du hältst der Glut stand, und das Feuer trennt dich von der Schlacke. Was abfällt, kehrt in die Erde zurück, und du findest neuen Grund.}
  {Du bist durch eine harte Zeit gegangen, in der vieles an dir verbrannt ist. Du hast standgehalten. Jetzt merkst du: Was abgefallen ist, war nicht das Wertvolle. Es geht zurück in den Boden und wird dort zu etwas Nützlichem. Und du stehst klarer da als vorher, auf festerem Grund.}
  {Wer nach jedem Feuer der Schlacke nachtrauert, sieht nicht, wie rein er geworden ist. Wer das Feuer meidet, trägt seine Schlacke weiter mit sich.}
\newpage
\spruch{47}{Metall}{Feuer}{Holz}{XI List}{Die Brandschneise}{Gegen das Feuer kommst du nicht an, also gehst du dorthin, wovon es lebt. Mit Axt und Säge schlägst du eine Schneise in den Wald, und wenn die Flammen sie erreichen, finden sie keine Nahrung mehr.}
  {Die Hitze, die dir entgegenschlägt, kannst du nicht direkt löschen. Also hältst du stand und siehst genau hin, woher sie ihre Kraft nimmt. Dort setzt du an, klar und entschieden. Ein Konflikt, dem man die Nahrung entzieht, brennt von selbst aus.}
  {Wer jede Schneise zu breit schlägt, verwüstet mehr als das Feuer. Wer keine schlägt, sieht zu, wie es sich ausbreitet.}
\newpage
\spruch{48}{Metall}{Feuer}{Wasser}{XII Aufstand}{Der Dampfkessel}{Der eiserne Kessel hält der Glut stand, und das Feuer kämpft gegen das Wasser darin an. Aus diesem doppelten Widerstand steigt Dampf, und er treibt mehr an, als du allein vermocht hättest.}
  {Du hältst stand, und was dich bedrängt, wird selbst bedrängt. Zwischen den Fronten entsteht Druck. Gefasst und gelenkt, wird dieser Druck zu einer Kraft, die vieles in Bewegung bringt. Frag dich, wohin sie führen soll.}
  {Ein Kessel ohne Ventil explodiert. Einer, der nie unter Druck steht, treibt nichts an.}

\newpage
\markright{\element{Wasser}}
\spruch{49}{Wasser}{Holz}{Feuer}{I Strom}{Die Holztrift}{Dein Fluss trägt die geschlagenen Stämme zu Tal, und in der Stadt am Ufer brennen sie in vielen Herden. Bricht dort ein Feuer aus, holt man das Löschwasser wieder aus dir.}
  {Aus deiner Ruhe ist etwas gewachsen, und es hat andere begeistert. Das Feuer, das jetzt brennt, ist nicht mehr deins. Aber du warst am Anfang da, und deine Ruhe wird gebraucht, wenn es zu heiß wird. Bleib in der Nähe.}
  {Wer jedes Feuer gleich löscht, verbietet anderen die Freude. Wer sich ganz zurückzieht, sieht zu, wie die Stadt brennt.}
\newpage
\spruch{50}{Wasser}{Holz}{Erde}{II Prüfung}{Der Auwald}{Du tränkst die Bäume, und ihre Wurzeln halten den Boden. Doch beim Hochwasser stehst du selbst dem Ufer gegenüber und musst deinen Weg durch die Erde finden.}
  {Du hast andere gestärkt, und sie haben Ordnung geschaffen, wo es eng und hart war. Jetzt aber stehst du selbst vor einer Grenze, und keiner bahnt dir den Weg. Das ist deine Prüfung. Du hast mehr Kraft, als dir bewusst ist. Wasser findet seinen Weg.}
  {Wer nie über die Ufer tritt, erfährt nie seine Kraft. Wer bei jedem Hochwasser alles mitreißt, zerstört auch, was ihn hält.}
\newpage
\spruch{51}{Wasser}{Holz}{Metall}{III Unterweisung}{Der Baum im Fels}{Du tränkst einen jungen Baum, der in einer Felsspalte wurzelt. Seine Wurzeln drängen gegen den harten Stein und weiten den Spalt, und in diesem Spalt sammelst du dich wieder.}
  {Du hast jemandem gegeben, was du weißt, und er geht damit seinen eigenen Weg, eigensinnig und widerständig. Du musst ihn nicht schützen. Gerade dort, wo er sich behauptet, zeigt er dir, woraus du selbst gemacht bist: aus Grundsätzen, an die du dich kaum noch erinnert hast.}
  {Wer seinen Schüler vor jedem harten Stein bewahrt, lässt ihn nie stark werden. Wer ihn allein lässt, sieht ihn im Fels vertrocknen.}
\newpage
\spruch{52}{Wasser}{Metall}{Erde}{IV Wurzelsuche}{Der unterirdische Fluss}{Du folgst deiner Quelle zurück in den Fels und dem Fels zurück in den Berg. Dort fließt du als unterirdischer Fluss und musst dir deinen Weg durch die Erde erst bahnen.}
  {Du fragst, woher deine Gewissheiten kommen, und gehst tiefer als je zuvor. Hinter allem Klaren liegt etwas Schweres, Dunkles: eine alte Last, eine Pflicht, die nie jemand infrage gestellt hat. Stell dich ihr. Wasser bahnt sich seinen Weg, auch durch den Berg, langsam, aber sicher.}
  {Wer immer tiefer in den Berg fließt, kommt nie mehr ans Licht. Wer an der Oberfläche bleibt, trocknet im ersten heißen Sommer aus.}
\newpage
\spruch{53}{Wasser}{Metall}{Holz}{V Gesammelte Kraft}{Der Obstgarten}{Im steinernen Becken sammelt sich das Regenwasser. Der Gärtner entfernt mit dem Messer die wilden Triebe, und erst danach gießt er den Baum, damit er Frucht trägt.}
  {Bevor du anderen hilfst zu wachsen, kehrst du zu deinen Grundsätzen zurück. Sie machen dir klar, was weichen muss, damit Wachstum gelingt. Erst danach gibst du, ruhig und gezielt. Was du so nährst, wächst nicht wild, sondern trägt Frucht.}
  {Wer gießt, ohne auszulichten, zieht Gestrüpp. Wer nur stutzt und nie gießt, hat bald einen dürren Stumpf.}
\newpage
\spruch{54}{Wasser}{Metall}{Feuer}{VI Rückhalt}{Die Felswand}{Du entspringst dem Fels, und der Fels hält der Glut stand, die sich ihm nähert. In seinem Schutz kannst du das Feuer löschen.}
  {Hinter dir steht etwas Festes: Grundsätze, die sich bewährt haben, Menschen mit klarem Urteil. Sie halten stand, wo die Hitze am größten ist. In ihrem Schutz behältst du die Ruhe und bringst zum Erlöschen, was außer Kontrolle geraten ist.}
  {Wer sich nur hinter dem Fels verschanzt, löscht nie selbst. Wer ohne Rückhalt ins Feuer geht, verdampft.}
\newpage
\spruch{55}{Wasser}{Feuer}{Erde}{VII Damm}{Der verschlammte Bach}{Du löschst den Brand, und die Asche wird zu neuem Boden. Doch der Regen spült ihn in dein Bett, und du musst dir deinen Weg neu graben.}
  {Du hast etwas Hitziges beruhigt, und daraus ist Neues entstanden, fruchtbar, aber schwer. Jetzt liegt es dir selbst im Weg. Es ist kein Undank, dass dein eigenes Werk dich herausfordert. Grab dich hindurch; der neue Lauf wird klarer sein als der alte.}
  {Wer jeden neuen Boden wegspült, lässt nichts wachsen. Wer sich vom Schlamm aufhalten lässt, wird zum Tümpel.}
\newpage
\spruch{56}{Wasser}{Feuer}{Holz}{VIII Abkürzung}{Der Wolkenbruch}{Du fällst als Regen auf das Feuer, das nach dem Wald greift. Wo es erlischt, trinken die Bäume, und du nährst sie ohne Umweg.}
  {Etwas Hitziges bedrohte, was dir am Herzen liegt. Du hast es beherzt gelöscht. Nun kannst du deine Ruhe direkt dorthin geben, wo sie gebraucht wird: zu dem, was wachsen will. Der kürzeste Weg zum Wachstum war, dem Brand die Hitze zu nehmen.}
  {Wer jedes Feuer im Wald löscht, lässt das dürre Holz sich häufen, bis es umso schlimmer brennt. Wer nie löscht, hat bald nichts mehr zu tränken.}
\newpage
\spruch{57}{Wasser}{Feuer}{Metall}{IX Strenger Weg}{Das Härten}{Im Feuer wird der Stahl zum Glühen gebracht, bis er die richtige Farbe hat; dann nimmt ihm dein Wasser zischend die Hitze. So wird er hart, und am Ende findest du Halt an ihm.}
  {Es ist eine Zeit strenger Übung: Hitze und Kälte wechseln, jeder Schritt verlangt Genauigkeit. Das ist anstrengend und lässt wenig Raum für Träume. Doch was du so bearbeitest, wird fest und klar. Am Ende dieses Weges hast du Grundsätze, auf die du dich stützen kannst.}
  {Wer zu schroff abschreckt, bekommt Risse im Stahl. Wer nie abschreckt, behält weichen Stahl.}
\newpage
\spruch{58}{Wasser}{Erde}{Metall}{X Läuterung}{Die Klamm}{Du gräbst dich durch die Erde, die dich aufhält, bis du auf Fels triffst. Dort wirst du klar, und in den Wänden glänzt das Erz.}
  {Lange hast du gegen etwas Schweres angekämpft, das dich trübte und bremste. Du hast dich ihm gestellt und dich hindurchgearbeitet. Jetzt stößt du auf festen Grund, auf etwas Klares, das dir niemand mehr nehmen kann. Aus der Mühsal ist ein Neubeginn geworden.}
  {Wer sich immer tiefer eingräbt, sieht nur noch die Wände seiner Schlucht. Wer der Erde ausweicht, bleibt trüb.}
\newpage
\spruch{59}{Wasser}{Erde}{Feuer}{XI List}{Der Wasserkrug}{Du nagst am Ufer und legst den Ton frei; ein Töpfer formt daraus einen Krug und brennt ihn. In diesem Krug trägt man dich zum Feuer, und du löschst es.}
  {Gegen das Feste kommst du nicht direkt an. Also arbeitest du geduldig an ihm, bis es sich wandelt und eine neue Gestalt annimmt. In dieser Gestalt aber trägt es dich, und du gelangst genau dorthin, wo du gebraucht wirst. Klugheit heißt, sich von seinem Gegner tragen zu lassen.}
  {Wer sich immer tragen lässt, vergisst, dass er selbst fließen kann. Wer jeden Krug zerschlägt, erreicht das Feuer nie.}
\newpage
\spruch{60}{Wasser}{Erde}{Holz}{XII Aufstand}{Das Hochwasser}{Du trittst über die Ufer, und der Boden klammert sich an die Wurzeln, um nicht fortgerissen zu werden. Wenn das Wasser sinkt, grünt das getränkte Land, fruchtbarer als zuvor.}
  {Du lehnst dich auf gegen das, was dich einengt, und es leistet seinerseits Widerstand. Eine Zeit lang ist alles in Aufruhr. Doch was danach bleibt, ist fruchtbar: Neues kann wachsen, wo vorher nichts war, und du hast es genährt.}
  {Ein Fluss, der ständig über die Ufer tritt, macht das Land zum Sumpf. Einer, der es nie tut, lässt die Aue vertrocknen.}

\endgroup

\addchap{Schlüssel}\markboth{Schlüssel}{\mbox{}}\label{kap:register}

\begin{einleitung}
Mit diesen Tafeln findest du jede Kombination, ohne zu blättern. Suche die erste Farbe in der linken Spalte und die nächste in der Kopfzeile; die Nummer im Schnittpunkt verweist auf das Paar oder das Tripel. Die Übersicht der Wendungen nennt zu jeder ihre beiden Richtungen, die Richtung von Anfang zu Ende und ihr Dreieck.
\end{einleitung}

\addsec*{Paare}\markright{Paare}
\begin{center}\renewcommand{\arraystretch}{1.6}\setlength{\tabcolsep}{5pt}
\begin{tabular}{r|ccccc}
{\footnotesize 1.\ \textbackslash\ 2.} & \kreis{Holz} & \kreis{Feuer} & \kreis{Erde} & \kreis{Metall} & \kreis{Wasser} \\ \hline
\kreis{Holz} & -- & \hyperlink{paar:Holz-Feuer}{\textcolor{ueberschriftfarbe}{\textbf{1}}} & \hyperlink{paar:Holz-Erde}{\textcolor{ueberschriftfarbe}{\textbf{3}}} & \hyperlink{paar:Holz-Metall}{\textcolor{ueberschriftfarbe}{\textbf{4}}} & \hyperlink{paar:Holz-Wasser}{\textcolor{ueberschriftfarbe}{\textbf{2}}} \\
\kreis{Feuer} & \hyperlink{paar:Feuer-Holz}{\textcolor{ueberschriftfarbe}{\textbf{6}}} & -- & \hyperlink{paar:Feuer-Erde}{\textcolor{ueberschriftfarbe}{\textbf{5}}} & \hyperlink{paar:Feuer-Metall}{\textcolor{ueberschriftfarbe}{\textbf{7}}} & \hyperlink{paar:Feuer-Wasser}{\textcolor{ueberschriftfarbe}{\textbf{8}}} \\
\kreis{Erde} & \hyperlink{paar:Erde-Holz}{\textcolor{ueberschriftfarbe}{\textbf{12}}} & \hyperlink{paar:Erde-Feuer}{\textcolor{ueberschriftfarbe}{\textbf{10}}} & -- & \hyperlink{paar:Erde-Metall}{\textcolor{ueberschriftfarbe}{\textbf{9}}} & \hyperlink{paar:Erde-Wasser}{\textcolor{ueberschriftfarbe}{\textbf{11}}} \\
\kreis{Metall} & \hyperlink{paar:Metall-Holz}{\textcolor{ueberschriftfarbe}{\textbf{15}}} & \hyperlink{paar:Metall-Feuer}{\textcolor{ueberschriftfarbe}{\textbf{16}}} & \hyperlink{paar:Metall-Erde}{\textcolor{ueberschriftfarbe}{\textbf{14}}} & -- & \hyperlink{paar:Metall-Wasser}{\textcolor{ueberschriftfarbe}{\textbf{13}}} \\
\kreis{Wasser} & \hyperlink{paar:Wasser-Holz}{\textcolor{ueberschriftfarbe}{\textbf{17}}} & \hyperlink{paar:Wasser-Feuer}{\textcolor{ueberschriftfarbe}{\textbf{19}}} & \hyperlink{paar:Wasser-Erde}{\textcolor{ueberschriftfarbe}{\textbf{20}}} & \hyperlink{paar:Wasser-Metall}{\textcolor{ueberschriftfarbe}{\textbf{18}}} & -- \\
\end{tabular}
\end{center}

\addsec*{Wendungen}\markright{Wendungen}
\begin{center}
\footnotesize\renewcommand{\arraystretch}{1.15}\setlength{\tabcolsep}{3pt}
\begin{tabular}{@{}rlccl@{}}
\toprule
\multicolumn{2}{c}{\textbf{Wendung}} & \textbf{Richtungen} & \textbf{Anfang\ra Ende} & \textbf{Dreieck} \\
\midrule
I    & \hyperlink{muster:I}{Der Strom} & +1, +1 & +2 & stumpf, Yin \\
II   & \hyperlink{muster:II}{Die Prüfung} & +1, +2 & −2 & spitz, Yang \\
III  & \hyperlink{muster:III}{Die Unterweisung} & +1, −2 & −1 & stumpf, Yin \\
IV   & \hyperlink{muster:IV}{Die Wurzelsuche} & −1, −1 & −2 & stumpf, Yin \\
V    & \hyperlink{muster:V}{Die gesammelte Kraft} & −1, +2 & +1 & stumpf, Yin \\
VI   & \hyperlink{muster:VI}{Der Rückhalt} & −1, −2 & +2 & spitz, Yang \\
VII  & \hyperlink{muster:VII}{Der Damm} & +2, +1 & −2 & spitz, Yang \\
VIII & \hyperlink{muster:VIII}{Die Abkürzung} & +2, −1 & +1 & stumpf, Yin \\
IX   & \hyperlink{muster:IX}{Der strenge Weg} & +2, +2 & −1 & spitz, Yang \\
X    & \hyperlink{muster:X}{Die Läuterung} & −2, +1 & −1 & stumpf, Yin \\
XI   & \hyperlink{muster:XI}{Die List} & −2, −1 & +2 & spitz, Yang \\
XII  & \hyperlink{muster:XII}{Der Aufstand} & −2, −2 & +1 & spitz, Yang \\
\bottomrule
\end{tabular}
\end{center}

\newpage
\addsec*{Tripel}\markright{Tripel}
\begin{center}\renewcommand{\arraystretch}{1.25}\setlength{\tabcolsep}{6pt}
\begin{tabular}{r|ccccc}
{\footnotesize 1.\,2.\ \textbackslash\ 3.} & \kreis{Holz} & \kreis{Feuer} & \kreis{Erde} & \kreis{Metall} & \kreis{Wasser} \\ \hline
\mypair[0.5em]{Holz}{Feuer} &  &  & \hyperlink{spruch:1}{\textbf{1}} & \hyperlink{spruch:2}{\textbf{2}} & \hyperlink{spruch:3}{\textbf{3}} \\
\mypair[0.5em]{Holz}{Erde} &  & \hyperlink{spruch:8}{\textbf{8}} &  & \hyperlink{spruch:7}{\textbf{7}} & \hyperlink{spruch:9}{\textbf{9}} \\
\mypair[0.5em]{Holz}{Metall} &  & \hyperlink{spruch:12}{\textbf{12}} & \hyperlink{spruch:11}{\textbf{11}} &  & \hyperlink{spruch:10}{\textbf{10}} \\
\mypair[0.5em]{Holz}{Wasser} &  & \hyperlink{spruch:5}{\textbf{5}} & \hyperlink{spruch:6}{\textbf{6}} & \hyperlink{spruch:4}{\textbf{4}} &  \\
\hline
\mypair[0.5em]{Feuer}{Holz} &  &  & \hyperlink{spruch:17}{\textbf{17}} & \hyperlink{spruch:18}{\textbf{18}} & \hyperlink{spruch:16}{\textbf{16}} \\
\mypair[0.5em]{Feuer}{Erde} & \hyperlink{spruch:15}{\textbf{15}} &  &  & \hyperlink{spruch:13}{\textbf{13}} & \hyperlink{spruch:14}{\textbf{14}} \\
\mypair[0.5em]{Feuer}{Metall} & \hyperlink{spruch:21}{\textbf{21}} &  & \hyperlink{spruch:20}{\textbf{20}} &  & \hyperlink{spruch:19}{\textbf{19}} \\
\mypair[0.5em]{Feuer}{Wasser} & \hyperlink{spruch:22}{\textbf{22}} &  & \hyperlink{spruch:24}{\textbf{24}} & \hyperlink{spruch:23}{\textbf{23}} &  \\
\hline
\mypair[0.5em]{Erde}{Holz} &  & \hyperlink{spruch:34}{\textbf{34}} &  & \hyperlink{spruch:36}{\textbf{36}} & \hyperlink{spruch:35}{\textbf{35}} \\
\mypair[0.5em]{Erde}{Feuer} & \hyperlink{spruch:28}{\textbf{28}} &  &  & \hyperlink{spruch:29}{\textbf{29}} & \hyperlink{spruch:30}{\textbf{30}} \\
\mypair[0.5em]{Erde}{Metall} & \hyperlink{spruch:26}{\textbf{26}} & \hyperlink{spruch:27}{\textbf{27}} &  &  & \hyperlink{spruch:25}{\textbf{25}} \\
\mypair[0.5em]{Erde}{Wasser} & \hyperlink{spruch:31}{\textbf{31}} & \hyperlink{spruch:33}{\textbf{33}} &  & \hyperlink{spruch:32}{\textbf{32}} &  \\
\hline
\mypair[0.5em]{Metall}{Holz} &  & \hyperlink{spruch:43}{\textbf{43}} & \hyperlink{spruch:45}{\textbf{45}} &  & \hyperlink{spruch:44}{\textbf{44}} \\
\mypair[0.5em]{Metall}{Feuer} & \hyperlink{spruch:47}{\textbf{47}} &  & \hyperlink{spruch:46}{\textbf{46}} &  & \hyperlink{spruch:48}{\textbf{48}} \\
\mypair[0.5em]{Metall}{Erde} & \hyperlink{spruch:42}{\textbf{42}} & \hyperlink{spruch:40}{\textbf{40}} &  &  & \hyperlink{spruch:41}{\textbf{41}} \\
\mypair[0.5em]{Metall}{Wasser} & \hyperlink{spruch:37}{\textbf{37}} & \hyperlink{spruch:38}{\textbf{38}} & \hyperlink{spruch:39}{\textbf{39}} &  &  \\
\hline
\mypair[0.5em]{Wasser}{Holz} &  & \hyperlink{spruch:49}{\textbf{49}} & \hyperlink{spruch:50}{\textbf{50}} & \hyperlink{spruch:51}{\textbf{51}} &  \\
\mypair[0.5em]{Wasser}{Feuer} & \hyperlink{spruch:56}{\textbf{56}} &  & \hyperlink{spruch:55}{\textbf{55}} & \hyperlink{spruch:57}{\textbf{57}} &  \\
\mypair[0.5em]{Wasser}{Erde} & \hyperlink{spruch:60}{\textbf{60}} & \hyperlink{spruch:59}{\textbf{59}} &  & \hyperlink{spruch:58}{\textbf{58}} &  \\
\mypair[0.5em]{Wasser}{Metall} & \hyperlink{spruch:53}{\textbf{53}} & \hyperlink{spruch:54}{\textbf{54}} & \hyperlink{spruch:52}{\textbf{52}} &  &  \\
\end{tabular}
\end{center}

% Anhang (eigene Datei); zum Abschalten ein % vor die nächste Zeile setzen
% anhang.tex: Anhang zum Pentagame-Orakel
% Eingebunden im Hauptdokument mit % anhang.tex: Anhang zum Pentagame-Orakel
% Eingebunden im Hauptdokument mit % anhang.tex: Anhang zum Pentagame-Orakel
% Eingebunden im Hauptdokument mit \include{anhang} (mit % davor abschalten).

\addchap{Anhang: Starke und schwache Siege}\markboth{Anhang}{\mbox{}}\label{kap:anhang}

\begin{einleitung}
Dieser Anhang verbindet das Orakel mit einer Beobachtung aus dem \textit{Pentagame Compendium} (Teil~IV, Kapitel~12). Er ist für Leser gedacht, die die Wendungen auch auf dem Brett weiterverfolgen möchten. Für die Befragung braucht man ihn nicht.
\end{einleitung}

\addsec*{Starke und schwache Siege}

\begin{einleitung}
Jede Wendung bildet ein stumpfes oder ein spitzes Dreieck (siehe \hyperref[kap:muster]{„Die zwölf Wendungen“}). Auf diese Form kommt es an, sobald man auch die Reihenfolge betrachtet, in der die Figuren herauskommen. Liegen die Ziele der ersten beiden Figuren nebeneinander (die erste Richtung ist ±1), bleibt für die dritte nur eine Möglichkeit, ein spitzes Dreieck zu bilden, aber zwei für ein stumpfes. Liegen sie übereck (±2), ist es umgekehrt. Der seltenere Abschluss ist der schwerere; das Compendium nennt ihn den stärkeren Sieg. Stark ist ein Sieg also, wenn auf zwei benachbarte Figuren ein spitzes Dreieck folgt oder auf zwei übereck liegende ein stumpfes.

Die erste Richtung und die Form des Dreiecks stecken vollständig in der Wendung. Ob ein Sieg stark oder schwach ist, hängt daher allein von der Wendung seines Tripels ab. Mit der Tabelle im Kapitel \hyperref[kap:muster]{„Die zwölf Wendungen“} ergibt sich: Stark sind \hyperlink{muster:II}{II, Die Prüfung}, \hyperlink{muster:VI}{VI, Der Rückhalt}, \hyperlink{muster:VIII}{VIII, Die Abkürzung} und \hyperlink{muster:X}{X, Die Läuterung}; die übrigen acht ergeben einen schwachen Sieg.

Alle vier starken wechseln zwischen Kreis und Stern: Auf eine Richtung von ±1 folgt eine von ±2, oder umgekehrt. Keine Wendung, die zweimal dieselbe Richtung geht, führt zu einem starken Sieg.

Das Compendium schlug vor, diese Unterscheidung bei einem Unentschieden (3:3) als Feinwertung zu nutzen: Ein starker Sieg schlägt dann einen schwachen. Im Spiel hat sich das nicht durchgesetzt. Für das Orakel zeigt es aber, dass die Wendungen nicht nur Deutungen sind, sondern auch auf dem Brett eine eigene Gestalt haben.
\end{einleitung} (mit % davor abschalten).

\addchap{Anhang: Starke und schwache Siege}\markboth{Anhang}{\mbox{}}\label{kap:anhang}

\begin{einleitung}
Dieser Anhang verbindet das Orakel mit einer Beobachtung aus dem \textit{Pentagame Compendium} (Teil~IV, Kapitel~12). Er ist für Leser gedacht, die die Wendungen auch auf dem Brett weiterverfolgen möchten. Für die Befragung braucht man ihn nicht.
\end{einleitung}

\addsec*{Starke und schwache Siege}

\begin{einleitung}
Jede Wendung bildet ein stumpfes oder ein spitzes Dreieck (siehe \hyperref[kap:muster]{„Die zwölf Wendungen“}). Auf diese Form kommt es an, sobald man auch die Reihenfolge betrachtet, in der die Figuren herauskommen. Liegen die Ziele der ersten beiden Figuren nebeneinander (die erste Richtung ist ±1), bleibt für die dritte nur eine Möglichkeit, ein spitzes Dreieck zu bilden, aber zwei für ein stumpfes. Liegen sie übereck (±2), ist es umgekehrt. Der seltenere Abschluss ist der schwerere; das Compendium nennt ihn den stärkeren Sieg. Stark ist ein Sieg also, wenn auf zwei benachbarte Figuren ein spitzes Dreieck folgt oder auf zwei übereck liegende ein stumpfes.

Die erste Richtung und die Form des Dreiecks stecken vollständig in der Wendung. Ob ein Sieg stark oder schwach ist, hängt daher allein von der Wendung seines Tripels ab. Mit der Tabelle im Kapitel \hyperref[kap:muster]{„Die zwölf Wendungen“} ergibt sich: Stark sind \hyperlink{muster:II}{II, Die Prüfung}, \hyperlink{muster:VI}{VI, Der Rückhalt}, \hyperlink{muster:VIII}{VIII, Die Abkürzung} und \hyperlink{muster:X}{X, Die Läuterung}; die übrigen acht ergeben einen schwachen Sieg.

Alle vier starken wechseln zwischen Kreis und Stern: Auf eine Richtung von ±1 folgt eine von ±2, oder umgekehrt. Keine Wendung, die zweimal dieselbe Richtung geht, führt zu einem starken Sieg.

Das Compendium schlug vor, diese Unterscheidung bei einem Unentschieden (3:3) als Feinwertung zu nutzen: Ein starker Sieg schlägt dann einen schwachen. Im Spiel hat sich das nicht durchgesetzt. Für das Orakel zeigt es aber, dass die Wendungen nicht nur Deutungen sind, sondern auch auf dem Brett eine eigene Gestalt haben.
\end{einleitung} (mit % davor abschalten).

\addchap{Anhang: Starke und schwache Siege}\markboth{Anhang}{\mbox{}}\label{kap:anhang}

\begin{einleitung}
Dieser Anhang verbindet das Orakel mit einer Beobachtung aus dem \textit{Pentagame Compendium} (Teil~IV, Kapitel~12). Er ist für Leser gedacht, die die Wendungen auch auf dem Brett weiterverfolgen möchten. Für die Befragung braucht man ihn nicht.
\end{einleitung}

\addsec*{Starke und schwache Siege}

\begin{einleitung}
Jede Wendung bildet ein stumpfes oder ein spitzes Dreieck (siehe \hyperref[kap:muster]{„Die zwölf Wendungen“}). Auf diese Form kommt es an, sobald man auch die Reihenfolge betrachtet, in der die Figuren herauskommen. Liegen die Ziele der ersten beiden Figuren nebeneinander (die erste Richtung ist ±1), bleibt für die dritte nur eine Möglichkeit, ein spitzes Dreieck zu bilden, aber zwei für ein stumpfes. Liegen sie übereck (±2), ist es umgekehrt. Der seltenere Abschluss ist der schwerere; das Compendium nennt ihn den stärkeren Sieg. Stark ist ein Sieg also, wenn auf zwei benachbarte Figuren ein spitzes Dreieck folgt oder auf zwei übereck liegende ein stumpfes.

Die erste Richtung und die Form des Dreiecks stecken vollständig in der Wendung. Ob ein Sieg stark oder schwach ist, hängt daher allein von der Wendung seines Tripels ab. Mit der Tabelle im Kapitel \hyperref[kap:muster]{„Die zwölf Wendungen“} ergibt sich: Stark sind \hyperlink{muster:II}{II, Die Prüfung}, \hyperlink{muster:VI}{VI, Der Rückhalt}, \hyperlink{muster:VIII}{VIII, Die Abkürzung} und \hyperlink{muster:X}{X, Die Läuterung}; die übrigen acht ergeben einen schwachen Sieg.

Alle vier starken wechseln zwischen Kreis und Stern: Auf eine Richtung von ±1 folgt eine von ±2, oder umgekehrt. Keine Wendung, die zweimal dieselbe Richtung geht, führt zu einem starken Sieg.

Das Compendium schlug vor, diese Unterscheidung bei einem Unentschieden (3:3) als Feinwertung zu nutzen: Ein starker Sieg schlägt dann einen schwachen. Im Spiel hat sich das nicht durchgesetzt. Für das Orakel zeigt es aber, dass die Wendungen nicht nur Deutungen sind, sondern auch auf dem Brett eine eigene Gestalt haben.
\end{einleitung}

\newpage

\addchap{Quellen und Literatur}\markboth{Quellen und Literatur}{\mbox{}}

\begin{einleitung}
Klassische chinesische Texte werden nach Titel und Kapitel zitiert, etwa \textit{Suwen} 68 oder \textit{Lunyu} 12.22; die Kapitelzählung ist in allen Ausgaben gleich. Die Originaltexte sind frei zugänglich im \href{https://ctext.org}{Chinese Text Project} (ctext.org). Übersetzungen der Zitate, wo nicht anders angegeben: eigene, mit Hilfe von KI angefertigt und am chinesischen Text geprüft.
\end{einleitung}

\addsec*{Klassische Texte}

\begin{itemize}[leftmargin=*]
  \item \textit{Shangshu} ({\cjk 尚書}), das \textit{Buch der Urkunden}, Kapitel „Hongfan“ ({\cjk 洪範}, „Der große Plan“): die fünf Elemente und ihre Natur. Englisch: James Legge, \textit{The Shoo King}, The Chinese Classics III, 1865.
  \item \textit{Lunyu} ({\cjk 論語}), die \textit{Gespräche} des Konfuzius, 12.22: rén als Liebe zu den Menschen. Deutsch: Richard Wilhelm, \textit{Kungfutse. Gespräche (Lun Yü)}, 1910.
  \item \textit{Huainanzi} ({\cjk 淮南子}), 139 v.\,Chr. Englisch: John S. Major u.\,a., \textit{The Huainanzi}, Columbia University Press 2010.
  \item \textit{Chunqiu fanlu} ({\cjk 春秋繁露}), Dong Zhongshu zugeschrieben, Kapitel 58 „Wuxing xiangsheng“: Hervorbringen und Besiegen. Englisch: Sarah A. Queen, John S. Major, \textit{Luxuriant Gems of the Spring and Autumn}, Columbia University Press 2016.
  \item \textit{Huangdi neijing suwen} ({\cjk 黃帝內經素問}), Kapitel 5 (die fünf Gemütsbewegungen), 67 (Besiegtes und Nicht-Besiegtes, Übermaß und Mangel) und 68 (Übermaß und Grenze). Englisch: Paul U. Unschuld, Hermann Tessenow, \textit{Huang Di Nei Jing Su Wen}, University of California Press 2011.
  \item \textit{Nanjing} ({\cjk 難經}), Schwierigkeiten 69 und 75: Mutter und Kind. Englisch: Paul U. Unschuld, \textit{Nan-ching. The Classic of Difficult Issues}, University of California Press 1986.
  \item \textit{Lingqijing} ({\cjk 靈棋經}), das Vorbild dieses Orakels. Englisch: Ralph D. Sawyer, Mei-chün Lee Sawyer, \textit{Ling Ch’i Ching. A Classic Chinese Oracle}, Shambhala 1995.
\end{itemize}

\addsec*{Weitere Quellen}

\begin{itemize}[leftmargin=*]
  \item Jan Suchanek, \textit{Pentagame Compendium}, 1.~Auflage 2022, \href{https://pentagame.org}{pentagame.org}: Regeln, Geschichte, Notation und Theorie des Spiels, auf dem Stand von 2022 (Brett mit 100 Feldern und früherer Farbordnung); zum Wu Xing Teil~III, Kapitel~9.
  \item \href{https://iep.utm.edu/wuxing/}{Wuxing (Internet Encyclopedia of Philosophy)}
  \item \href{https://www.britannica.com/topic/wuxing}{Wuxing (Encyclopaedia Britannica)}
  \item \href{https://plato.stanford.edu/entries/chinese-phil-medicine/}{Chinese Philosophy and Chinese Medicine (Stanford Encyclopedia of Philosophy)}
  \item \href{https://en.wikipedia.org/wiki/Lingqijing}{Lingqijing (Wikipedia)}
  \item \href{https://russellcottrell.com/VirtualYarrowStalks/LingQiJing.htm}{The Mathematics of the Ling Qi Jing (Russell Cottrell)}
  \item \href{https://muse.jhu.edu/article/397150/summary}{Rezension von Sawyer/Sawyer, \textit{Ling Ch’i Ching}, China Review International 1997}
\end{itemize}

\end{document}